% Media and communication — Anglais Tle Tech — Cours signé OnBuch+
% Programme officiel MINESEC. Compiler avec : tectonic N05-media-and-communication.tex
\newcommand{\DOCMATIERE}{Anglais}
\newcommand{\DOCNIVEAU}{Tle Tech}
\newcommand{\DOCMODULE}{Anglais – classe de Terminale technique}
\newcommand{\DOCLECON}{N05}
\newcommand{\DOCTITRE}{Media and communication}
% OnBuch+ — préambule « cours signés » ENRICHI (compilé avec tectonic / XeLaTeX)
% Système visuel « Pop » d'OnBuch+ : fond crème (jamais blanc), bordures encre
% épaisses, ombres « dures » décalées, titres Archivo Black en capitales.
% Version MATHÉMATIQUES : graphes (pgfplots/tikz), boîtes pédagogiques
% (définition, propriété, méthode, attention, le savais-tu, exercice résolu,
% exercice, TP, bilan), page de garde et signature OnBuch+ ancrée en bas.
%
% Macros attendues dans main.tex AVANT \input{preamble} :
%   \DOCMATIERE  \DOCNIVEAU  \DOCMODULE  \DOCLECON (numéro)  \DOCTITRE
\documentclass[11pt]{article}

\usepackage{fontspec}
\usepackage[french]{babel}
\usepackage[a4paper,margin=2cm,top=3.4cm,bottom=2.8cm,headheight=1.4cm,footskip=1.3cm]{geometry}
\usepackage{amsmath,amssymb,amsfonts}
\usepackage{mathtools} % \xmapsto, \coloneqq... souvent utilisés par les modèles
\usepackage{cancel} % simplifications barrées (bilans de forces)
% \abs{z} (module, valeur absolue) et \norm{u} (norme), utilisés par les modèles
\providecommand{\abs}{}\let\abs\relax\DeclarePairedDelimiter{\abs}{\lvert}{\rvert}
\providecommand{\norm}{}\let\norm\relax\DeclarePairedDelimiter{\norm}{\lVert}{\rVert}
\usepackage[table]{xcolor} % [table] : fournit aussi \rowcolors (alternance de lignes)
\usepackage{pagecolor}
\usepackage{tikz}
\usetikzlibrary{arrows.meta,positioning,calc,shapes.geometric,shapes.misc,shapes.arrows,decorations.pathmorphing,decorations.pathreplacing,decorations.markings,patterns,shadows,mindmap,trees,fit,backgrounds,matrix,intersections,angles,quotes}
\usepackage{pgfplots}
\usepackage[version=4]{mhchem} % équations nucléaires \ce{^{235}_{92}U}
\usepackage[european,siunitx]{circuitikz} % schémas électriques
\pgfplotsset{compat=1.18}
\usepgfplotslibrary{fillbetween,groupplots} % aires entre courbes (calcul intégral)
\usepackage{enumitem}
\usepackage{fancyhdr}
% [most] charge toutes les bibliothèques tcolorbox (listings, minted, raster,
% documentation...) et gonfle inutilement la mémoire TeX sur un document long
% (~300 boîtes) : on ne charge que ce qui est réellement utilisé.
\usepackage[skins,breakable]{tcolorbox}
\usepackage{graphicx}
\usepackage{colortbl}
\usepackage{booktabs}
\usepackage{tabularx}
\usepackage{multirow}
\usepackage{diagbox} % case d'en-tête diagonale des tableaux à double entrée
% Colonnes extensibles centrée / gauche / droite (C, L, R), souvent utilisées par les modèles
\newcolumntype{C}{>{\centering\arraybackslash}X}
\newcolumntype{L}{>{\raggedright\arraybackslash}X}
\newcolumntype{R}{>{\raggedleft\arraybackslash}X}
\usepackage{array}
\usepackage{multicol}
\usepackage{siunitx}
% parse-numbers=false : accepte aussi \SI{2,5 \times 10^{-3}}{...}
\sisetup{locale=FR,output-decimal-marker={,},group-minimum-digits=4,parse-numbers=false}
\DeclareSIUnit\ohm{\ensuremath{\symup{\Omega}}} % U+2126 absent de la police
\DeclareSIUnit\division{div} % réglages d'oscilloscope (ms/div, V/div)
\DeclareSIUnit\tr{tr} % tours (vitesse de rotation, ex. \SI{1500}{\tr\per\minute})
\DeclareSIUnit\molar{M} % concentration molaire (ex. \SI{3}{\molar}, \SI{1}{\micro\molar})
\DeclareSIUnit\an{an} % année (le modèle confond parfois avec \year, un compteur TeX) — ex. \SI{5}{\kilo\meter\per\an}
\DeclareSIUnit\FCFA{FCFA} % franc CFA (ex. \SI{500}{\FCFA\per\kilo\gram})
\DeclareSIUnit\degreeBrix{\textdegree Brix} % teneur en sucre d'un sirop/jus (réfractométrie)
\DeclareSIUnit\month{mois} % le modèle utilise parfois \month, un compteur TeX, comme unité de durée
\DeclareSIUnit\microliter{\micro\liter} % le modèle écrit parfois \microliter en un seul token
\DeclareSIUnit\million{millions} % dénombrement (ex. \SI{15}{\million\per\mL} spermatozoïdes)
\usepackage{qrcode}
% \degree : commande fréquemment utilisée par les modèles pour un angle ou une
% température en dehors de \SI{}{} (non fournie par les paquets chargés).
\providecommand{\degree}{^{\circ}}
% \qed (fin de démonstration), utilisé par les modèles sans amsthm
\providecommand{\qed}{\hfill$\square$}
% \barre : double barre des valeurs interdites dans un tableau de variations (tkz-tab)
\providecommand{\barre}{\ensuremath{\|}}
% environnement proof (amsthm non chargé) utilisé par les modèles
\ifdefined\proof\else\newenvironment{proof}[1][Démonstration]{\par\noindent\textit{#1.}\ }{\qed\par}\fi
\usepackage{lastpage}
\usepackage{hyperref}
\hypersetup{hidelinks}

% ============================================================
% Palette « Pop » OnBuch+ — mêmes hex que la classe P (plus_ui.dart)
% ============================================================
\definecolor{poporange}{HTML}{F59321}
\definecolor{popdark}{HTML}{E07A0C}
\definecolor{popcream}{HTML}{FFF6E8}
\definecolor{popink}{HTML}{1C1714}
\definecolor{poppurple}{HTML}{7A5AE0}
\definecolor{popgreen}{HTML}{2E9E6A}
\definecolor{poppink}{HTML}{E0457B}
\definecolor{popblue}{HTML}{2F7DE1}
\definecolor{popgold}{HTML}{C98A12}
\definecolor{popmuted}{HTML}{8A7F76}
\definecolor{popline}{HTML}{EADFD2}
\definecolor{popgray}{HTML}{8A7F76}
\colorlet{popgrey}{popgray}
% Alias tolérants (le modèle nomme parfois les couleurs différemment)
\colorlet{popwhite}{white}
\colorlet{popblack}{popink}
\colorlet{popred}{poppink}
\colorlet{popyellow}{popgold}
% Teintes claires (fonds de tableaux/encadrés) — le modèle nomme parfois
% les couleurs avec un suffixe L (ex. popblueL) sans les avoir définies.
\colorlet{poporangeL}{poporange!20}
\colorlet{popdarkL}{popdark!20}
\colorlet{poppurpleL}{poppurple!20}
\colorlet{popgreenL}{popgreen!20}
\colorlet{poppinkL}{poppink!20}
\colorlet{popblueL}{popblue!20}
\colorlet{popgoldL}{popgold!20}
\colorlet{popredL}{popred!20}
\colorlet{popgrayL}{popgray!20}
% Teintes claires pour fonds de tableaux / zones
\colorlet{popblueL}{popblue!10!white}
\colorlet{poporangeL}{poporange!14!white}
\colorlet{popgreenL}{popgreen!12!white}
\colorlet{poppurpleL}{poppurple!12!white}
\colorlet{poppinkL}{poppink!12!white}
\colorlet{popgoldL}{popgold!14!white}

\pagecolor{popcream}

% ============================================================
% Typographie
% ============================================================
\defaultfontfeatures{Ligatures=TeX}
\setmainfont{PlusJakartaSans-Medium.ttf}[Path=../../../fonts/,BoldFont=PlusJakartaSans-Bold.ttf]
% Maths en sans-serif (Fira Math) avec chiffres et lettres droites de Plus
% Jakarta, pour que les formules se fondent dans le texte.
\usepackage[math-style=ISO,bold-style=ISO]{unicode-math}
\setmathfont{FiraMath-Regular.otf}[Path=../../../fonts/]
\setmathfont{PlusJakartaSans-Medium.ttf}[Path=../../../fonts/,range={up/num,up/{Latin,latin},\mathup}]
\usepackage{icomma} % virgule décimale sans espace en mode math
% Caractères mathématiques Unicode absents des polices de texte (Plus Jakarta,
% Archivo Black) : rendus par leur équivalent mathématique, en texte comme en
% formule — sinon ils s'affichent en carré vide (ex. « Arithmétique dans ℤ »).
\usepackage{newunicodechar}
\newunicodechar{ℝ}{\ensuremath{\mathbb{R}}}
\newunicodechar{ℤ}{\ensuremath{\mathbb{Z}}}
\newunicodechar{ℂ}{\ensuremath{\mathbb{C}}}
\newunicodechar{ℕ}{\ensuremath{\mathbb{N}}}
\newunicodechar{ℚ}{\ensuremath{\mathbb{Q}}}
\newunicodechar{𝔻}{\ensuremath{\mathbb{D}}}
\newunicodechar{α}{\ensuremath{\alpha}}
\newunicodechar{β}{\ensuremath{\beta}}
\newunicodechar{γ}{\ensuremath{\gamma}}
\newunicodechar{δ}{\ensuremath{\delta}}
\newunicodechar{ε}{\ensuremath{\varepsilon}}
\newunicodechar{θ}{\ensuremath{\theta}}
\newunicodechar{λ}{\ensuremath{\lambda}}
\newunicodechar{μ}{\ensuremath{\mu}}
\newunicodechar{σ}{\ensuremath{\sigma}}
\newunicodechar{τ}{\ensuremath{\tau}}
\newunicodechar{φ}{\ensuremath{\varphi}}
\newunicodechar{ψ}{\ensuremath{\psi}}
\newunicodechar{ω}{\ensuremath{\omega}}
\newunicodechar{Σ}{\ensuremath{\Sigma}}
% majuscules produites par \MakeUppercase dans les titres (α-aminés -> Α)
\newunicodechar{Α}{\ensuremath{\alpha}}
\newunicodechar{Β}{\ensuremath{\beta}}
\newunicodechar{Γ}{\ensuremath{\Gamma}}
\newunicodechar{Δ}{\ensuremath{\Delta}}
\newunicodechar{Ε}{\ensuremath{\varepsilon}}
\newunicodechar{Θ}{\ensuremath{\Theta}}
\newunicodechar{Λ}{\ensuremath{\Lambda}}
\newunicodechar{Μ}{\ensuremath{\mu}}
\newunicodechar{Τ}{\ensuremath{\tau}}
\newunicodechar{Φ}{\ensuremath{\Phi}}
\newunicodechar{Ψ}{\ensuremath{\Psi}}
\newunicodechar{Ω}{\ensuremath{\Omega}}
\newunicodechar{↦}{\ensuremath{\mapsto}}
\newunicodechar{∈}{\ensuremath{\in}}
\newunicodechar{∉}{\ensuremath{\notin}}
\newunicodechar{∧}{\ensuremath{\wedge}}
\newunicodechar{⇔}{\ensuremath{\Leftrightarrow}}
\newunicodechar{⟺}{\ensuremath{\Longleftrightarrow}}
\newunicodechar{⇒}{\ensuremath{\Rightarrow}}
\newunicodechar{⟹}{\ensuremath{\Longrightarrow}}
\newunicodechar{∘}{\ensuremath{\circ}}
\newunicodechar{∪}{\ensuremath{\cup}}
\newunicodechar{∩}{\ensuremath{\cap}}
\newunicodechar{≡}{\ensuremath{\equiv}}
\newunicodechar{⊂}{\ensuremath{\subset}}
\newunicodechar{⊥}{\ensuremath{\perp}}
\newunicodechar{∃}{\ensuremath{\exists}}
\newunicodechar{∀}{\ensuremath{\forall}}
\newunicodechar{∛}{\ensuremath{\sqrt[3]{\,}}}
\newunicodechar{∜}{\ensuremath{\sqrt[4]{\,}}}
\newunicodechar{⃗}{\ensuremath{{}^{\to}}}
\newunicodechar{ⁿ}{\ifmmode^{n}\else\textsuperscript{\ensuremath{n}}\fi}
\newunicodechar{ˣ}{\ifmmode^{x}\else\textsuperscript{\ensuremath{x}}\fi}
\newunicodechar{ᵇ}{\ifmmode^{b}\else\textsuperscript{\ensuremath{b}}\fi}
\newunicodechar{ʳ}{\ifmmode^{r}\else\textsuperscript{\ensuremath{r}}\fi}
\newunicodechar{ᵘ}{\ifmmode^{u}\else\textsuperscript{\ensuremath{u}}\fi}
\newunicodechar{ʸ}{\ifmmode^{y}\else\textsuperscript{\ensuremath{y}}\fi}
\newunicodechar{ᵃ}{\ifmmode^{a}\else\textsuperscript{\ensuremath{a}}\fi}
\newunicodechar{ᵉ}{\ifmmode^{e}\else\textsuperscript{\ensuremath{e}}\fi}
\newunicodechar{ᵏ}{\ifmmode^{k}\else\textsuperscript{\ensuremath{k}}\fi}
\newunicodechar{ᵖ}{\ifmmode^{p}\else\textsuperscript{\ensuremath{p}}\fi}
\newunicodechar{ᴺ}{\ifmmode^{N}\else\textsuperscript{\ensuremath{N}}\fi}
\newunicodechar{ᶜ}{\ifmmode^{c}\else\textsuperscript{\ensuremath{c}}\fi}
\newunicodechar{ʰ}{\ifmmode^{h}\else\textsuperscript{\ensuremath{h}}\fi}
\newunicodechar{ᵠ}{\ifmmode^{q}\else\textsuperscript{\ensuremath{q}}\fi}
\newunicodechar{ₙ}{\ifmmode_{n}\else\textsubscript{\ensuremath{n}}\fi}
\newunicodechar{ₐ}{\ifmmode_{a}\else\textsubscript{\ensuremath{a}}\fi}
\newunicodechar{ᵢ}{\ifmmode_{i}\else\textsubscript{\ensuremath{i}}\fi}
\newunicodechar{ᵣ}{\ifmmode_{r}\else\textsubscript{\ensuremath{r}}\fi}
\newunicodechar{ₖ}{\ifmmode_{k}\else\textsubscript{\ensuremath{k}}\fi}
\newunicodechar{ᵦ}{\ifmmode_{\beta}\else\textsubscript{\ensuremath{\beta}}\fi}
\newunicodechar{ₓ}{\ifmmode_{x}\else\textsubscript{\ensuremath{x}}\fi}
\newfontfamily\popdisplay{ArchivoBlack-Regular.ttf}[Path=../../../fonts/]
\newfontfamily\poplogo{SpaceGrotesk-Bold.ttf}[Path=../../../fonts/]
\newfontfamily\popsemi{PlusJakartaSans-SemiBold.ttf}[Path=../../../fonts/]
\newfontfamily\popxbold{PlusJakartaSans-ExtraBold.ttf}[Path=../../../fonts/]
\newfontfamily\popxboldls{PlusJakartaSans-ExtraBold.ttf}[Path=../../../fonts/,LetterSpace=3.0]

\setlist[enumerate]{leftmargin=1.7em,itemsep=5pt,topsep=4pt}
\setlist[itemize]{leftmargin=1.7em,itemsep=4pt,topsep=4pt,label={\color{poporange}\textbullet}}
\setlength{\parindent}{0pt}
\setlength{\parskip}{5pt}
\renewcommand{\arraystretch}{1.25}

% pgfplots : style Pop par défaut
\pgfplotsset{
  every axis/.append style={axis line style={popink,line width=1pt},tick style={popink},
    grid style={popline,line width=0.5pt},label style={font=\small},tick label style={font=\footnotesize}},
  popcurve/.style={poporange,line width=1.8pt,no markers,smooth},
}
% Même style disponible dans un \draw TikZ simple (hors pgfplots)
\tikzset{popcurve/.style={poporange,line width=1.8pt}}
\tikzset{popgrid/.style={popline,very thin}}
\colorlet{popcurve}{poporange} % le nom du style est parfois utilisé comme couleur

% ============================================================
% Pill / squiggle
% ============================================================
\newcommand{\poppill}[3]{%
  \tikz[baseline=-0.55ex]\node[draw=popink,line width=1.2pt,rounded corners=2.6pt,%
    fill=#2,inner xsep=6.5pt,inner ysep=3pt]{%
    \popxboldls\fontsize{7}{7}\selectfont\color{#3}\MakeUppercase{#1}};%
}
\newcommand{\popsquiggle}[1]{%
  \tikz[baseline=-0.4ex]{
    \draw[#1,line width=2.6pt,line cap=round]
      plot [smooth] coordinates {(0,0.09) (0.34,0.01) (0.68,0.21) (1.02,0.01) (1.36,0.10)};
  }%
}
% Mot-clé surligné (marqueur orange sous le texte)
\newcommand{\cle}[1]{{\popxbold\color{popdark}#1}}

% ============================================================
% En-tête / pied
% ============================================================
\pagestyle{fancy}
\fancyhf{}
\renewcommand{\headrulewidth}{0pt}
\renewcommand{\footrulewidth}{0pt}
\lhead{\raisebox{-0.15ex}{\poplogo\fontsize{13}{13}\selectfont\color{popink}OnBuch\color{poporange}+}}
\rhead{\poppill{\DOCMATIERE\ \textbullet\ \DOCNIVEAU\ \textbullet\ Leçon \DOCLECON}{white}{popink}}
\lfoot{\popxbold\fontsize{7.6}{7.6}\selectfont\color{popmuted}\MakeUppercase{Cours signé OnBuch+}\ \textbullet\ {\popsemi\DOCTITRE}}
\rfoot{\tikz[baseline=-0.6ex]\node[rounded corners=8.5pt,fill=popink,minimum height=17pt,minimum width=17pt,inner xsep=5pt,inner sep=0]{\popxbold\fontsize{8}{8}\selectfont\color{popcream}\thepage\,/\,\pageref*{LastPage}};}

% ============================================================
% Titres
% ============================================================
\newcommand{\chaptitre}[2]{%
  \begin{tcolorbox}[enhanced,colback=poporange,colframe=popink,boxrule=2.6pt,arc=11pt,
      drop shadow=popink,left=15pt,right=15pt,top=12pt,bottom=11pt,before skip=3pt,after skip=18pt]
    \poppill{#2}{popink}{popcream}\par\vspace{9pt}
    {\raggedright\hyphenpenalty=10000\exhyphenpenalty=10000\popdisplay\fontsize{22}{24}\selectfont\color{popcream}\MakeUppercase{#1}\par}
  \end{tcolorbox}
}
% Section numérotée : pastille orange avec le numéro + titre Archivo
\newcounter{popsec}
\newcommand{\coursec}[1]{%
  \stepcounter{popsec}\setcounter{popsub}{0}%
  \par\vspace{16pt}\noindent
  \begin{minipage}{\linewidth}
  \tikz[baseline=-0.7ex]\node[draw=popink,line width=1.6pt,fill=poporange,rounded corners=4pt,minimum size=22pt,inner sep=0]{\popdisplay\fontsize{12}{12}\selectfont\color{popcream}\thepopsec};%
  \hspace{8pt}{\popdisplay\fontsize{15}{17}\selectfont\color{popink}\MakeUppercase{#1}}\par
  \vspace{4pt}\noindent{\color{poporange}\rule{36pt}{3.6pt}}
  \end{minipage}\par\nopagebreak\vspace{8pt}
}
\newcounter{popsub}
\newcommand{\courssub}[1]{%
  \stepcounter{popsub}%
  \par\vspace{9pt}\noindent{\popxbold\fontsize{11.5}{13}\selectfont\color{popink}{\color{poporange}\thepopsec.\thepopsub}\hspace{6pt}#1}\par\nopagebreak\vspace{4pt}
}

% ============================================================
% Boîtes pédagogiques — même recette : fond blanc, bordure épaisse colorée,
% ombre dure encre, badge plein. Titre optionnel [#1].
% ============================================================
\tcbset{popbase/.style={breakable,enhanced,colback=white,boxrule=2.3pt,arc=9pt,
  shadow={3pt}{-3pt}{0pt}{popink},left=14pt,right=14pt,top=11pt,bottom=11pt,before skip=11pt,after skip=15pt,
  fontupper=\popsemi\fontsize{10.6}{14.5}\selectfont\color{popink},
  fontlower=\popsemi\fontsize{10.6}{14.5}\selectfont\color{popink}}}
% Coche dessinée (le glyphe ✓ n'existe pas dans les polices du document)
\newcommand{\popcheck}[1]{\tikz[baseline=-0.5ex]\draw[#1,line width=1.6pt,line cap=round,line join=round](-0.13,0.0)--(-0.04,-0.09)--(0.13,0.1);}
\providecommand{\coche}{\popcheck{popgreen}}
\providecommand{\resultat}[1]{\par\smallskip\noindent\fcolorbox{popgreen}{popgreenL}{\parbox{\dimexpr\linewidth-2\fboxsep-2\fboxrule\relax}{\textbf{Résultat.} #1}}\par\smallskip}
\newcommand{\popboxhead}[3]{\poppill{#1}{#2}{white}\ifx\relax#3\relax\else\hspace{6pt}{\popxbold\fontsize{10.6}{12}\selectfont\color{popink}#3}\fi\par\vspace{7pt}}

\newtcolorbox{aretenir}[1][]{popbase,colframe=popgreen,before upper={\popboxhead{À retenir}{popgreen}{#1}}}
\newtcolorbox{exemplebox}[1][]{popbase,colframe=poporange,before upper={\popboxhead{Exemple}{poporange}{#1}}}
\newtcolorbox{definition}[1][]{popbase,colframe=popblue,before upper={\popboxhead{Définition}{popblue}{#1}}}
\newtcolorbox{propriete}[1][]{popbase,colframe=poppurple,before upper={\popboxhead{Propriété}{poppurple}{#1}}}
\newtcolorbox{methode}[1][]{popbase,colframe=popgold,before upper={\popboxhead{Méthode}{popgold}{#1}}}
\newtcolorbox{attention}[1][]{popbase,colframe=poppink,before upper={\popboxhead{Attention}{poppink}{#1}}}
\newtcolorbox{experience}[1][]{popbase,colframe=popblue,colback=popblueL,before upper={\popboxhead{Expérience}{popblue}{#1}}}
% « Le savais-tu ? » : carte violette pleine (contexte, culture, Cameroun)
\newtcolorbox{savaistu}[1][]{popbase,colback=poppurple,colframe=popink,
  fontupper=\popsemi\fontsize{10.4}{14.2}\selectfont\color{white},
  before upper={\poppill{Le savais-tu\,?}{white}{poppurple}\ifx\relax#1\relax\else\hspace{6pt}{\popxbold\color{white}#1}\fi\par\vspace{7pt}}}
% Exercice résolu : énoncé en haut, \tcblower puis la solution sur fond crème
\newcounter{popexo}\newcounter{popexr}
\newtcolorbox{exoresolu}[1][]{popbase,colframe=poporange,colbacklower=popcream,
  segmentation style={popink,line width=1.2pt,dashed},
  before upper={\stepcounter{popexr}\popboxhead{Exercice résolu \thepopexr}{poporange}{#1}},
  before lower={\poppill{Solution}{popink}{popcream}\par\vspace{6pt}}}
% Exercice (non résolu en place) — la correction est dans la partie Corrigés
\newtcolorbox{exercice}[1][]{popbase,colframe=popink,
  before upper={\stepcounter{popexo}\popboxhead{Exercice \thepopexo}{popink}{#1}}}
% Corrigé
\newtcolorbox{corrige}[1][]{popbase,colframe=popgreen,colback=popgreenL,
  before upper={\popboxhead{Corrigé}{popgreen}{#1}}}
% Objectifs / prérequis / bilan
\newtcolorbox{objectifs}{popbase,colframe=popink,colback=poporangeL,
  before upper={\popboxhead{Objectifs de la leçon}{popink}{}}}
\newtcolorbox{prerequis}{popbase,colframe=popink,colback=popblueL,
  before upper={\popboxhead{Prérequis}{popblue}{}}}
\newtcolorbox{bilan}{popbase,colframe=popink,colback=white,boxrule=2.8pt,
  before upper={\popboxhead{Fiche bilan}{poporange}{L'essentiel à connaître pour l'examen}}}
% Figure encadrée avec légende
\newtcolorbox{popfigure}[1][]{enhanced,breakable=false,colback=white,colframe=popline,boxrule=1.4pt,arc=8pt,
  left=10pt,right=10pt,top=10pt,bottom=8pt,before skip=10pt,after skip=12pt,halign=center,
  lower separated=false,halign lower=center,
  fontlower=\popsemi\fontsize{9}{11.5}\selectfont\color{popmuted},
  #1}
\newcommand{\legende}[1]{\par\vspace{6pt}{\popsemi\fontsize{9}{11.5}\selectfont\color{popmuted}#1}}

% ============================================================
% Signature OnBuch+
% ============================================================
\newcommand{\poplogoinline}[1]{{\poplogo\fontsize{#1}{#1}\selectfont\color{popink}OnBuch\color{poporange}+}}
\newcommand{\signatureonbuch}{%
  \begin{tcolorbox}[enhanced,colback=popink,colframe=popink,boxrule=2.2pt,arc=10pt,
      drop shadow=poporange,left=14pt,right=14pt,top=10pt,bottom=10pt,before skip=0pt,after skip=0pt]
    \tikz[baseline=-0.7ex]\node[circle,fill=popgreen,draw=popcream,line width=1.3pt,minimum size=22pt,inner sep=0]{\popcheck{white}};%
    \hspace{9pt}\begin{minipage}[c]{0.62\linewidth}
      {\popxbold\fontsize{12}{13}\selectfont\color{popcream}Signé\ \textbullet\ {\poplogo OnBuch\color{poporange}+}}\par\vspace{2pt}
      {\raggedright\popsemi\fontsize{8}{10.5}\selectfont\color{popline}\MakeUppercase{Équipe pédagogique OnBuch+}\\\MakeUppercase{Conforme au programme officiel MINESEC}\par}
    \end{minipage}\hfill
    {\poplogo\fontsize{22}{22}\selectfont\color{popcream}OnBuch\color{poporange}+}
  \end{tcolorbox}%
}

% ============================================================
% Page de garde
% #1 titre · #2 matière · #3 niveau · #4 module · #5 n° leçon · #6 durée ·
% #7 description / objectifs (texte ou itemize)
% ============================================================
\newcommand{\pagegarde}[7]{%
  \thispagestyle{empty}
  \newgeometry{a4paper,margin=1.7cm,top=1.6cm,bottom=1.4cm}%
  \begin{tikzpicture}[remember picture,overlay]
    \fill[poporange!18] ([xshift=-3.2cm,yshift=-3.6cm]current page.north east) circle (4.6cm);
    \fill[poppurple!14] ([xshift=2.4cm,yshift=7.6cm]current page.south west) circle (3.8cm);
  \end{tikzpicture}%
  {\poplogo\fontsize{30}{30}\selectfont\color{popink}OnBuch\color{poporange}+}\hfill
  \raisebox{0.6ex}{\poppill{Cours signé \textbullet\ Premium}{popink}{popcream}}\par
  \vspace{1pt}\popsquiggle{poppurple}\par
  \vspace{8pt}
  \begin{tcolorbox}[enhanced,colback=poporange,colframe=popink,boxrule=2.8pt,arc=13pt,
      drop shadow=popink,left=18pt,right=18pt,top=12pt,bottom=13pt,after skip=10pt]
    \poppill{#2 \textbullet\ #3}{popink}{popcream}\hspace{5pt}\poppill{#4}{white}{popink}\par\vspace{8pt}
    {\popdisplay\fontsize{14}{14}\selectfont\color{popink}LEÇON #5}\par\vspace{4pt}
    {\raggedright\hyphenpenalty=10000\exhyphenpenalty=10000\popdisplay\fontsize{25}{27}\selectfont\color{popcream}\MakeUppercase{#1}\par}
    \vspace{8pt}
    {\popxbold\fontsize{9.5}{9.5}\selectfont\color{popink}\MakeUppercase{Durée conseillée : #6}}
  \end{tcolorbox}
  \begin{tcolorbox}[enhanced,colback=white,colframe=popink,boxrule=2.2pt,arc=10pt,
      drop shadow=popink,left=16pt,right=16pt,top=10pt,bottom=10pt,after skip=8pt]
    \poppill{Dans cette leçon}{poppurple}{white}\par\vspace{5pt}
    \popsemi\fontsize{9.3}{12.6}\selectfont\color{popink}#7
  \end{tcolorbox}
  \noindent\begin{minipage}[t]{0.487\linewidth}
    \begin{tcolorbox}[enhanced,colback=popgreen,colframe=popink,boxrule=2.2pt,arc=10pt,
        drop shadow=popink,left=11pt,right=11pt,top=10pt,bottom=10pt,height=3.1cm,valign=center]
      \tcbox[colback=white,colframe=popink,boxrule=1.3pt,arc=5pt,boxsep=0pt,nobeforeafter,
        left=3pt,right=3pt,top=3pt,bottom=3pt,valign=center]{\qrcode[height=1.9cm]{https://play.google.com/store/apps/details?id=com.luvvix.onbuch&hl=fr}}%
      \hspace{8pt}\begin{minipage}[c]{0.5\linewidth}
        {\popdisplay\fontsize{11}{12.5}\selectfont\color{white}\MakeUppercase{Télécharge OnBuch}}\par\vspace{4pt}
        {\raggedright\popsemi\fontsize{8.4}{11}\selectfont\color{white}Gratuit \textbullet\ Google Play\\Tuteur IA, annales, concours, résultats\par}
      \end{minipage}
    \end{tcolorbox}
  \end{minipage}\hfill
  \begin{minipage}[t]{0.487\linewidth}
    \begin{tcolorbox}[enhanced,colback=poppurple,colframe=popink,boxrule=2.2pt,arc=10pt,
        drop shadow=popink,left=11pt,right=11pt,top=10pt,bottom=10pt,height=3.1cm,valign=center]
      \tikz[baseline=-0.7ex]\node[circle,fill=white,draw=popink,line width=1.3pt,minimum size=1.6cm,inner sep=0]{\popdisplay\fontsize{24}{24}\selectfont\color{poppurple}+};%
      \hspace{8pt}\begin{minipage}[c]{0.6\linewidth}
        {\popdisplay\fontsize{11}{12.5}\selectfont\color{white}\MakeUppercase{Passe à OnBuch+}}\par\vspace{4pt}
        {\raggedright\popsemi\fontsize{8.4}{11}\selectfont\color{white}Tous les cours signés, Tuteur IA illimité, fiches PDF — onglet OnBuch+ de l'app\par}
      \end{minipage}
    \end{tcolorbox}
  \end{minipage}
  \vspace{8pt}
  \popcontenu
  \vfill
  \signatureonbuch
  \restoregeometry
  \newpage
}

% Rangée « Ce cours contient » (page de garde)
\newcommand{\poptile}[3]{% #1 couleur #2 gros texte #3 légende
  \begin{tcolorbox}[enhanced,colback=#1,colframe=popink,boxrule=2pt,arc=9pt,shadow={2.5pt}{-2.5pt}{0pt}{popink},
      left=6pt,right=6pt,top=9pt,bottom=9pt,halign=center,valign=center,height=1.75cm,width=\linewidth,nobeforeafter]
    {\popdisplay\fontsize{12.5}{13}\selectfont\color{white}\MakeUppercase{#2}}\par\vspace{3pt}
    {\centering\popsemi\fontsize{7.8}{9.5}\selectfont\color{white}#3\par}
  \end{tcolorbox}}
\newcommand{\popcontenu}{%
  \par\noindent{\popxboldls\fontsize{7.5}{7.5}\selectfont\color{popmuted}CE COURS SIGNÉ CONTIENT}\par\vspace{7pt}
  \noindent\begin{minipage}[t]{0.235\linewidth}\poptile{popblue}{Cours}{complet, illustré, pas à pas}\end{minipage}\hfill
  \begin{minipage}[t]{0.235\linewidth}\poptile{poporange}{Méthodes}{et exercices résolus}\end{minipage}\hfill
  \begin{minipage}[t]{0.235\linewidth}\poptile{poppink}{Exercices}{type examen + corrigés}\end{minipage}\hfill
  \begin{minipage}[t]{0.235\linewidth}\poptile{popgreen}{Fiche bilan}{l'essentiel à retenir}\end{minipage}\par}

% Sommaire
\newtcolorbox{sommaire}{enhanced,colback=white,colframe=popink,boxrule=2.2pt,arc=10pt,
  drop shadow=popink,left=18pt,right=18pt,top=13pt,bottom=13pt,before skip=6pt,after skip=20pt,
  before upper={\poppill{Sommaire}{poppurple}{white}\par\vspace{9pt}\popsemi\fontsize{10.6}{14.5}\selectfont\color{popink}}}

% Signature de fin de cours — poussée tout en bas de la dernière page
\newcommand{\findecours}{%
  \par\vfill\nopagebreak
  \begin{center}\popsquiggle{poporange}\end{center}\vspace{4pt}
  \signatureonbuch
}
\AtBeginDocument{\def\llbracket{\mathopen{[\mkern-3.5mu[}}\def\rrbracket{\mathclose{]\mkern-3.5mu]}}} % intervalles d'entiers (glyphe ⟦ absent de la police)
% Glyphes absents de Fira Math : redéfinis à partir de glyphes disponibles
\AtBeginDocument{%
  \def\setminus{\mathbin{\backslash}}\let\smallsetminus\setminus
  \def\vdots{\mathord{\vbox{\baselineskip3.2pt\lineskiplimit0pt\kern4pt\hbox{.}\hbox{.}\hbox{.}}}}%
  \def\ddots{\mathinner{\mkern1mu\raise6.5pt\hbox{.}\mkern2mu\raise3.6pt\hbox{.}\mkern2mu\raise0.7pt\hbox{.}\mkern1mu}}%
  \def\triangle{\mathord{\scalebox{0.9}{$\symup{\Delta}$}}}%
  \def\nabla{\mathord{\raisebox{\depth}{\scalebox{1}[-1]{$\symup{\Delta}$}}}}%
  \def\checkmark{\popcheck{popdark}}%
  \def\star{\mathbin{\ast}}\def\bigstar{\mathord{\ast}}%
  \def\diamond{\mathbin{\scalebox{0.6}{$\Diamond$}}}%
  \def\bigcup{\mathop{\mathchoice{\raisebox{-0.3em}{\scalebox{1.7}{$\cup$}}}{\scalebox{1.25}{$\cup$}}{\cup}{\cup}}\displaylimits}%
  \def\bigcap{\mathop{\mathchoice{\raisebox{-0.3em}{\scalebox{1.7}{$\cap$}}}{\scalebox{1.25}{$\cap$}}{\cap}{\cap}}\displaylimits}%
  \def\lesseqqgtr{\mathrel{\lessgtr}}\def\bigoplus{\mathop{\oplus}\displaylimits}%
}
\providecommand{\ssi}{si et seulement si} % abréviation fréquente
\AtBeginDocument{\providecommand{\wideparen}{\overparen}} % arc de cercle

%% FIGFIX-BEGIN v3 — correctifs globaux des figures (docs/onbuch-generation/tools/figfix)
\usepackage{adjustbox}\usepackage{etoolbox}
% toute figure TikZ/pgfplots plus large que la ligne est réduite à la largeur disponible
\BeforeBeginEnvironment{tikzpicture}{\begin{adjustbox}{max width=\linewidth}}
\AfterEndEnvironment{tikzpicture}{\end{adjustbox}}
% légendes pgfplots lisibles par-dessus les courbes
\pgfplotsset{every axis/.append style={legend style={fill=white,fill opacity=0.92,text opacity=1,draw=black!15,inner sep=3pt}}}
% étiquettes d'arêtes (syntaxe edge["..."]) avec fond blanc
\tikzset{every edge quotes/.append style={fill=white,inner sep=1.2pt}}
% bulles de cartes mentales : texte à la ligne dans la bulle, bulle un peu plus grande
\tikzset{every concept/.append style={text width=2.0cm,align=center,minimum size=2.3cm,inner sep=2pt}}
%% FIGFIX-END
\begin{document}

\pagegarde{Media and communication}{Anglais}{Tle Tech}{Anglais – classe de Terminale technique}{N05}{9 séances (fiche de progression harmonisée)}{This lesson develops students' ability to communicate about media and ICTs in English: expressing preferences when subscribing to telephone and internet packages, understanding texts on ICTs for work, learning and recreation, and writing a well-structured composition on the recreational use of ICTs at home and in the community.
\vspace{4pt}
\begin{multicols}{2}\small
\begin{itemize}
\item À la fin de cette leçon, je sais utiliser le vocabulaire anglais lié aux abonnements téléphoniques et internet (bundle, data, airtime, subscription, provider...)
\item À la fin de cette leçon, je sais exprimer mes préférences et justifier un choix de forfait téléphonique ou internet (I'd rather..., I prefer... because...)
\item À la fin de cette leçon, je sais comprendre un texte audio sur l'utilisation des TIC pour le travail ou les études et en relever les informations essentielles
\item À la fin de cette leçon, je sais lire et exploiter un texte sur les usages récréatifs des TIC à la maison ou dans la communauté
\item À la fin de cette leçon, je sais réviser et mobiliser les structures grammaticales clés (comparatifs, modaux, present perfect, linkers) dans des échanges oraux et écrits
\item À la fin de cette leçon, je sais rédiger une composition structurée (introduction, développement, conclusion) sur les usages récréatifs des TIC
\end{itemize}\end{multicols}}

\begin{sommaire}
\begin{multicols}{2}
\begin{enumerate}
\item Subscribing to service packages: vocabulary and speaking
\item Grammar general revision
\item Listening: ICTs to enhance work and learning
\item Reading: ICTs as recreational resources
\item Writing: composition on recreational use of ICTs
\item Integration activities, evaluation and remediation
\item Activité d'intégration
\item Méthodes et exercices résolus
\item Exercices
\item Corrigés des exercices
\item Fiche bilan
\end{enumerate}
\end{multicols}
\end{sommaire}

\begin{objectifs}
\begin{itemize}
\item À la fin de cette leçon, je sais utiliser le vocabulaire anglais lié aux abonnements téléphoniques et internet (bundle, data, airtime, subscription, provider...)
\item À la fin de cette leçon, je sais exprimer mes préférences et justifier un choix de forfait téléphonique ou internet (I'd rather..., I prefer... because...)
\item À la fin de cette leçon, je sais comprendre un texte audio sur l'utilisation des TIC pour le travail ou les études et en relever les informations essentielles
\item À la fin de cette leçon, je sais lire et exploiter un texte sur les usages récréatifs des TIC à la maison ou dans la communauté
\item À la fin de cette leçon, je sais réviser et mobiliser les structures grammaticales clés (comparatifs, modaux, present perfect, linkers) dans des échanges oraux et écrits
\item À la fin de cette leçon, je sais rédiger une composition structurée (introduction, développement, conclusion) sur les usages récréatifs des TIC
\item À la fin de cette leçon, je sais interpréter des données chiffrées sur la connectivité au Cameroun et les utiliser pour justifier une réponse
\item À la fin de cette leçon, je sais adopter un comportement responsable et critique face aux médias et réseaux sociaux
\end{itemize}
\end{objectifs}

\begin{prerequis}
\begin{itemize}
\item Vocabulaire de base des médias et des TIC vu en classe de Première (phone, internet, computer, social media)
\item Expression des goûts et préférences (like, love, hate, enjoy + V-ing)
\item Comparatifs et superlatifs (cheaper than, the best offer)
\item Modaux de base (can, must, should) pour exprimer capacité, obligation et conseil
\item Structure d'un paragraphe et connecteurs simples (and, but, because, so)
\end{itemize}
\end{prerequis}

\newpage

\coursec{Subscribing to service packages: vocabulary and speaking}

\courssub{Lead-in: Aminatou's dilemma}

Aminatou is a final-year student at Lycée Technique de Douala. Next month, she must hand in a research project for her Probatoire, and she needs a good internet connection to download documents, watch tutorial videos and exchange files with her classmates. She stands in front of an MTN agency on Avenue de Gaulle, but she hesitates: MTN, Orange and Camtel all advertise attractive offers. Which one should she choose? At home, the situation is the same for everyone: her father uses WhatsApp video calls to follow cocoa prices with his partners in Kumba, while her younger brother spends his evenings on TikTok and online games. The same phone, the same network, but very different uses.

This raises two questions for this lesson: \emph{Are ICTs in Cameroon mainly tools for work and learning, or for entertainment?} And more practically: \emph{How can one choose a service package wisely and use digital tools responsibly?} In this first part of the lesson, you will learn the vocabulary of telephone and internet subscriptions, and you will practise expressing and justifying your preferences like a real customer in a Cameroonian agency.

\begin{aretenir}[Lesson objectives]
By the end of this section, you will be able to:
\begin{itemize}
\item name the main elements of a telephone or internet subscription (airtime, data bundle, validity, top up...);
\item compare service packages using price, volume, validity and network quality;
\item express and justify a preference orally (I prefer... to..., I'd rather..., I'd go for... because...);
\item perform a short dialogue between a customer and a sales agent.
\end{itemize}
\end{aretenir}

\courssub{Vocabulary: the language of subscriptions}

Before choosing a package, you must understand the words on the posters and in the USSD menus. Read the following table carefully; the French gloss will help you memorise each term.

\begin{center}
\begin{tabularx}{\linewidth}{|l|X|l|}\hline
\rowcolor{popblueL} \textbf{English term} & \textbf{Meaning} & \textbf{French} \\ \hline
airtime / credit & money loaded on a phone line to make calls, send SMS or buy bundles & crédit téléphonique \\ \hline
data bundle & a fixed amount of internet data bought for a given price and validity period & forfait internet \\ \hline
SMS pack & a number of text messages bought at a reduced price & forfait SMS \\ \hline
call plan & a tariff or package for voice calls & forfait d'appels \\ \hline
subscription & an agreement to pay regularly for a service & abonnement \\ \hline
provider / operator & the company that sells the service (MTN, Orange, Camtel, Nexttel) & opérateur / fournisseur \\ \hline
network coverage & the geographical area where the signal is available & couverture réseau \\ \hline
validity period & the time during which a bundle can be used (1 day, 7 days, 30 days) & période de validité \\ \hline
unlimited & without a fixed limit of data or calls & illimité \\ \hline
to top up / to recharge & to add credit to a phone line & recharger \\ \hline
roaming & using your phone on a foreign network when travelling abroad & itinérance \\ \hline
Wi-Fi hotspot & a place or device that gives wireless internet access & point d'accès Wi-Fi \\ \hline
USSD code & a short code you dial to access services, e.g. *123\# & code USSD \\ \hline
\end{tabularx}
\end{center}

\begin{definition}[Data bundle]
A \cle{data bundle} is a fixed amount of internet data (for example 200 MB, 1 GB or 5 GB) that a customer buys from a provider for a given price and for a given validity period. When the data is finished or the validity period expires, the customer must buy a new bundle.
\end{definition}

\begin{exemplebox}[Using the vocabulary in context]
Complete sentences with the words above:
\begin{itemize}
\item ``My bundle has expired. I need to \textbf{top up} my line and buy a new one.''
\item ``There is no MTN \textbf{network coverage} in my village, so I use Orange.''
\item ``This weekly bundle has a \textbf{validity period} of seven days.''
\item ``When my uncle travels to Nigeria, he activates \textbf{roaming} to keep his Cameroonian number active.''
\end{itemize}
\end{exemplebox}

\courssub{Cameroonian providers and their offers}

In Cameroon, four main \cle{providers} share the market: \textbf{MTN} (whose brand colour is yellow and whose greeting is ``Y'ello''), \textbf{Orange}, \textbf{Camtel} (the historical state operator, also present in fibre optic internet) and \textbf{Nexttel}. Each provider sells different categories of bundles:

\begin{itemize}
\item \textbf{daily bundles}: small volumes (e.g. 100--500 MB), valid for 24 hours, cheap;
\item \textbf{weekly bundles}: medium volumes (e.g. 1--2 GB), valid for 7 days;
\item \textbf{monthly bundles}: large volumes (e.g. 5--30 GB), valid for 30 days, often cheaper per GB;
\item \textbf{social bundles}: data restricted to applications like WhatsApp, Facebook or TikTok, cheaper but limited;
\item \textbf{full internet bundles}: data usable on all websites and applications.
\end{itemize}

To subscribe, customers usually dial a \cle{USSD code} (for example *123\# or \#150\#), navigate a menu, choose an offer and confirm. An SMS then confirms the subscription.

\begin{popfigure}
\begin{center}
\begin{tikzpicture}
\node[draw=popdark, thick, rounded corners, fill=popcream, minimum width=\linewidth, minimum height=0.6cm] {\textbf{Compare these three fictitious offers (oral exploitation)}};
\end{tikzpicture}
\vspace{0.2cm}
\begin{tikzpicture}[x=1cm,y=1cm]
% Header row
\fill[popblue] (0,0) rectangle (3.4,-0.8);
\fill[popblue] (3.4,0) rectangle (7.0,-0.8);
\fill[popblue] (7.0,0) rectangle (10.6,-0.8);
\fill[popblue] (10.6,0) rectangle (14.2,-0.8);
\node[white,font=\bfseries] at (1.7,-0.4) {Criteria};
\node[white,font=\bfseries] at (5.2,-0.4) {Offer A: Daily};
\node[white,font=\bfseries] at (8.8,-0.4) {Offer B: Weekly};
\node[white,font=\bfseries] at (12.4,-0.4) {Offer C: Monthly};
% Rows - values with commas braced for foreach parsing
\foreach \y/\lab/\a/\b/\c in {
 -0.8/Price/200 FCFA/{1,000 FCFA}/{4,000 FCFA},
 -1.6/Data volume/200 MB/2 GB/8 GB,
 -2.4/Validity/24 hours/7 days/30 days,
 -3.2/Price per day/200 FCFA/{about 143 FCFA}/{about 133 FCFA},
 -4.0/Type/Social bundle/Full internet/Full internet}{
  \draw[popline] (0,\y) rectangle (14.2,\y-0.8);
  \draw[popline] (3.4,\y) -- (3.4,\y-0.8);
  \draw[popline] (7.0,\y) -- (7.0,\y-0.8);
  \draw[popline] (10.6,\y) -- (10.6,\y-0.8);
  \node[font=\bfseries\small, anchor=west] at (0.15,\y-0.4) {\lab};
  \node[font=\small] at (5.2,\y-0.4) {\a};
  \node[font=\small] at (8.8,\y-0.4) {\b};
  \node[font=\small] at (12.4,\y-0.4) {\c};
}
\draw[popdark, thick] (0,0) rectangle (14.2,-4.8);
\end{tikzpicture}
\end{center}
\legende{Comparative table of three fictitious bundles. Use it to practise: ``Which offer do you prefer? Why?'' Prices in FCFA, volumes in MB/GB.}
\end{popfigure}

\begin{exoresolu}[Choosing the best offer]
\textbf{Task:} Aminatou needs internet for 7 days of research. She hesitates between Offer A (daily, 200 MB at 200 FCFA) and Offer B (weekly, 2 GB at 1,000 FCFA). Which one is more economical? Justify in one sentence.
\tcblower
\textbf{Solution:} For 7 days, Offer A would cost $7 \times 200 = 1\,400$ FCFA and give only $7 \times 200 = 1\,400$ MB $= \num{1,4}$ GB. Offer B costs 1\,000 FCFA and gives 2 GB. Offer B is cheaper and gives more data. Sentence: \emph{``I'd rather take the 2 GB weekly bundle at 1,000 FCFA because it is cheaper per day than the daily 200 MB bundle at 200 FCFA.''}
\end{exoresolu}

\courssub{Expressing and justifying preferences}

Choosing a package is not enough: at the Probatoire and Baccalauréat, you must \emph{say why} you choose it. Here is the functional language you need.

\begin{center}
\begin{tabularx}{\linewidth}{|l|X|}\hline
\rowcolor{popblueL} \textbf{Function} & \textbf{Structures and examples} \\ \hline
Expressing a preference & I \textbf{prefer} MTN \textbf{to} Orange. / I \textbf{'d rather} take the weekly bundle. / I\textbf{'d go for} the monthly plan. \\ \hline
Asking for a preference & \textbf{What about} the social bundle? / \textbf{Would you rather} have more data or more call time? \\ \hline
Giving an opinion & \textbf{In my opinion}, coverage matters more than price. / \textbf{I think} Camtel fibre is the fastest. \\ \hline
Justifying & ... \textbf{because} it is cheaper. / ... \textbf{since} my village has no 4G. / ... \textbf{as} the validity is longer. \\ \hline
Comparing & cheaper \textbf{than}, more data \textbf{than}, the \textbf{most} reliable, the \textbf{best} value for money \\ \hline
\end{tabularx}
\end{center}

\begin{methode}[How to express a justified preference]
Follow these three steps every time you speak:
\begin{enumerate}
\item \textbf{State your preference} with a correct structure: \emph{I prefer X to Y} / \emph{I'd rather take X} / \emph{I'd go for X}.
\item \textbf{Add a linker}: \emph{because}, \emph{since} or \emph{as}.
\item \textbf{Give a concrete argument}: price (FCFA), volume (MB/GB), validity, network coverage or bonus offers.
\end{enumerate}
Example: \emph{``I'd go for the monthly bundle \textbf{because} it is cheaper per gigabyte and \textbf{since} I use the internet every day.''}
\end{methode}

\begin{attention}[Common mistake: ``prefer... than'']
Many candidates write or say: \emph{``I prefer MTN \textbf{than} Orange.''} This is \textbf{wrong}. The correct structure is:
\[ \text{I prefer } X \textbf{ to } Y. \]
Example: \emph{I prefer the weekly bundle \textbf{to} the daily one.} Note also: \emph{I'd rather take A \textbf{than} B} (with ``rather'', we use ``than'', never ``to'').
\end{attention}

\courssub{Grammar checkpoint: comparatives and superlatives}

Since comparing offers is central to this lesson, here is a quick revision of the forms you need.

\begin{center}
\begin{tabular}{|l|l|l|}\hline
\rowcolor{popblueL} \textbf{Type} & \textbf{Rule} & \textbf{Examples} \\ \hline
Short adjectives (1 syllable) & \texttt{-er / -est} & cheap $\rightarrow$ cheaper $\rightarrow$ the cheapest \\ \hline
Long adjectives (2+ syllables) & \texttt{more / most} & expensive $\rightarrow$ more expensive $\rightarrow$ the most expensive \\ \hline
Two syllables ending in \texttt{-y} & \texttt{-ier / -iest} & heavy $\rightarrow$ heavier $\rightarrow$ the heaviest \\ \hline
Irregular & -- & good $\rightarrow$ better $\rightarrow$ best \quad bad $\rightarrow$ worse $\rightarrow$ worst \\ \hline
\end{tabular}
\end{center}

\begin{exemplebox}[Practice]
\begin{itemize}
\item Offer A is \textbf{cheaper than} Offer B \textbf{but} Offer B has \textbf{more data}.
\item Offer C is \textbf{the most expensive} \textbf{but} it offers \textbf{the best} value for money per GB.
\item My connection is \textbf{better} today \textbf{than} yesterday.
\end{itemize}
\end{exemplebox}

\courssub{Guided dialogue: subscribing by USSD code}

Read this dialogue aloud in pairs. Aminatou finally decides to subscribe from home, using a USSD code.

\begin{exemplebox}[Dialogue: Aminatou and her friend Divine]
\textbf{Divine:} Have you bought your bundle for the research project yet?\par
\textbf{Aminatou:} Not yet. I hesitate between the daily and the weekly offer. What do you think?\par
\textbf{Divine:} In my opinion, the weekly bundle is better value for money. How much data do you need?\par
\textbf{Aminatou:} About 2 GB, to download PDFs and watch tutorials.\par
\textbf{Divine:} Then I'd go for the 2 GB weekly bundle at 1,000 FCFA. Just dial *123\#, choose option 2, then confirm.\par
\textbf{Aminatou:} OK... I dial *123\#... I select ``Internet bundles''... then ``Weekly''... then ``2 GB -- 1,000 FCFA''... I confirm.\par
\textbf{Divine:} Did you receive the confirmation SMS?\par
\textbf{Aminatou:} Yes! ``You have successfully subscribed to the 2 GB weekly bundle, valid for 7 days.'' Thank you, Divine!\par
\textbf{Divine:} You're welcome. Now you can work on your Probatoire project!
\end{exemplebox}

\begin{popfigure}
\begin{center}
\begin{tikzpicture}[node distance=0.55cm,
stepbox/.style={draw=popblue, thick, rounded corners, fill=popblueL, text width=4.6cm, align=center, minimum height=1cm, font=\small},
arr/.style={-{Stealth[length=3mm]}, thick, popdark}]
\node[stepbox] (s1) {\textbf{Step 1:} Dial the USSD code (e.g. *123\#)};
\node[stepbox, below=of s1] (s2) {\textbf{Step 2:} Choose the offer in the menu (daily / weekly / monthly)};
\node[stepbox, below=of s2] (s3) {\textbf{Step 3:} Confirm the purchase (airtime is deducted)};
\node[stepbox, below=of s3, fill=popgreenL, draw=popgreen] (s4) {\textbf{Step 4:} Receive the confirmation SMS};
\draw[arr] (s1) -- (s2);
\draw[arr] (s2) -- (s3);
\draw[arr] (s3) -- (s4);
\end{tikzpicture}
\end{center}
\legende{Steps to subscribe to a bundle by USSD code: dial the code, choose the offer, confirm, then receive the confirmation SMS.}
\end{popfigure}

\courssub{Pair work: role-play at the agency}

Now it is your turn to speak. Work in pairs: one student is the \textbf{customer}, the other is the \textbf{sales agent} at an agency in Yaoundé or Bafoussam. Use the comparative table of the three fictitious offers, the vocabulary list and the preference structures.

\begin{exercice}[Role-play: at the MTN agency in Bafoussam]
\textbf{Customer:} You want a bundle for one month of online courses. You have 4,000 FCFA. Ask about prices, volumes, validity and coverage in your quarter. Express your preference and justify it.\par
\textbf{Sales agent:} Present the three offers (daily, weekly, monthly). Answer the customer's questions, suggest the best offer and explain the subscription steps by USSD code.\par
\textbf{Useful phrases:} \emph{How much is the monthly bundle? -- What about the validity? -- I'd rather... because... -- In my opinion... -- I prefer... to...}
\end{exercice}

\begin{exercice}[Written consolidation]
\begin{enumerate}
\item Complete: ``I prefer the weekly bundle \_\_\_\_ the daily one because it is cheaper per day.''
\item Correct the mistake: ``I prefer Orange than MTN.''
\item Write two sentences comparing Offer A and Offer C using \emph{I'd rather... because...}.
\item Order the steps of a USSD subscription: receive SMS / dial the code / confirm / choose the offer.
\end{enumerate}
\end{exercice}

\begin{corrige}[Exercice: Written consolidation]
\begin{enumerate}
\item I prefer the weekly bundle \textbf{to} the daily one because it is cheaper per day.
\item I prefer Orange \textbf{to} MTN.
\item Possible answers: \emph{I'd rather take Offer C because 8 GB is enough for a whole month.} / \emph{I'd rather avoid Offer A because 200 MB is not enough for my online courses.}
\item Correct order: dial the code $\rightarrow$ choose the offer $\rightarrow$ confirm $\rightarrow$ receive the SMS.
\end{enumerate}
\end{corrige}

\begin{savaistu}[Mobile telephony in Cameroon]
Did you know? Cameroon counts more than \textbf{20 million mobile subscribers} -- almost one line per inhabitant! However, internet \cle{coverage} remains unequal: big cities like Douala, Yaoundé and Bafoussam enjoy 4G, while many rural areas still depend on 2G or have no signal at all. This is why comparing \emph{network coverage} is just as important as comparing prices when you choose a provider.
\end{savaistu}

\begin{aretenir}[Summary]
\begin{itemize}
\item A \cle{data bundle} = a fixed volume of data + a price + a validity period.
\item Compare packages with four criteria: \textbf{price} (FCFA), \textbf{volume} (MB/GB), \textbf{validity} and \textbf{coverage}.
\item Express a justified preference: \emph{I prefer X to Y / I'd rather... / I'd go for...} + \emph{because/since/as} + argument.
\item Never say ``prefer... than''; say ``prefer... \textbf{to}''.
\item To subscribe by USSD: dial the code, choose the offer, confirm, receive the SMS.
\end{itemize}
\end{aretenir}

\coursec{Grammar general revision}

In the previous section, you learned how to talk about telephone and internet service packages: bundles, subscriptions, data volume, validity. You practised sentences like \emph{``I would like to subscribe to the weekly bundle''}. Now we will sharpen the grammar tools you need to compare packages, ask polite questions at an agency, give advice about safe ICT use, and write coherent compositions. This is a \cle{general revision} (révision générale) : all these structures were introduced in earlier classes, but at the Baccalauréat technique they must be \textbf{accurate and automatic} (précis et automatiques).

\courssub{Comparatives and superlatives: comparing packages and devices}

\textbf{Intuition.} When you choose a bundle, you constantly compare: \emph{Which network is faster? Which package is the cheapest?} English has two tools for this: the \cle{comparative} (comparatif, comparing two items) and the \cle{superlative} (superlatif, singling out one item from a group).

\begin{propriete}[Comparatives and superlatives: the rules]
\begin{itemize}
\item \textbf{Short adjectives} (adjectifs courts : one syllable, sometimes two): add \cle{-er} for the comparative and \cle{the + -est} for the superlative: \emph{fast $\rightarrow$ faster $\rightarrow$ the fastest}; \emph{cheap $\rightarrow$ cheaper $\rightarrow$ the cheapest}. Adjectives ending in a single consonant double it: \emph{big $\rightarrow$ bigger $\rightarrow$ the biggest}.
\item \textbf{Long adjectives} (adjectifs longs : two syllables or more): use \cle{more + adjective} for the comparative and \cle{the most + adjective} for the superlative: \emph{reliable $\rightarrow$ more reliable $\rightarrow$ the most reliable}; \emph{expensive $\rightarrow$ more expensive $\rightarrow$ the most expensive}.
\item \textbf{Adjectives ending in -y}: \emph{easy $\rightarrow$ easier $\rightarrow$ the easiest}.
\item \textbf{Irregular forms} (formes irrégulières, to memorise): \emph{good $\rightarrow$ better $\rightarrow$ the best}; \emph{bad $\rightarrow$ worse $\rightarrow$ the worst}; \emph{far $\rightarrow$ further/farther $\rightarrow$ the furthest/farthest}.
\item The comparative is followed by \cle{than} (que): \emph{This router is faster \textbf{than} the old one.}
\end{itemize}
\end{propriete}

\begin{exemplebox}[Comparing two data packages]
Package A gives 2 GB for 1\,000 FCFA; Package B gives 5 GB for 2\,000 FCFA.
\begin{itemize}
\item \emph{Package B is \textbf{more economical than} Package A for heavy users.}
\item \emph{Package A is \textbf{cheaper than} Package B, but it offers \textbf{less} data.}
\item \emph{Of all the packages, the monthly bundle is \textbf{the most convenient}.}
\end{itemize}
\end{exemplebox}

\begin{attention}[A very frequent error: ``more cheaper'']
Never combine \emph{more} with an \emph{-er} form. \emph{More cheaper}, \emph{more faster}, \emph{the most cheapest} are all wrong. Choose \textbf{one} system only: \emph{cheaper} or \emph{more expensive}. In the same way, never write \emph{``the most fastest''} — write \emph{the fastest}. This error is heavily penalised in the writing paper.
\end{attention}

\courssub{Modals: can/could, should/ought to, must/mustn't}

\textbf{Intuition.} Modals are small auxiliary verbs that colour the main verb: they express \cle{ability} (capacité), \cle{request} (demande), \cle{advice} (conseil) or \cle{obligation} (obligation). They never take \emph{-s} in the third person and are always followed by the \textbf{bare infinitive} (infinitif sans \emph{to}).

\begin{center}
\begin{tabularx}{\linewidth}{|l|l|X|}
\hline
\rowcolor{popblueL} \textbf{Modal} & \textbf{Function} & \textbf{Example in an ICT context} \\
\hline
can & ability & \emph{You can buy a bundle with mobile money.} \\
\hline
could & polite request & \emph{Could you show me the price list, please?} \\
\hline
should / ought to & advice & \emph{You should use a strong password.} \\
\hline
must & strong obligation & \emph{Students must switch off their phones during exams.} \\
\hline
mustn't & prohibition & \emph{You mustn't share your PIN code with anyone.} \\
\hline
\end{tabularx}
\end{center}

\begin{popfigure}
\begin{center}
\begin{tikzpicture}[font=\small]
\node[draw=popblue, fill=popblueL, rounded corners, minimum width=4.2cm, minimum height=1cm] (ability) at (0,3) {\textbf{Ability / Request}};
\node[draw=popgreen, fill=popgreenL, rounded corners, minimum width=4.2cm, minimum height=1cm] (advice) at (0,1.5) {\textbf{Advice}};
\node[draw=poporange, fill=poporangeL, rounded corners, minimum width=4.2cm, minimum height=1cm] (rules) at (0,0) {\textbf{Rules / Prohibition}};
\node[draw=poppurple, fill=poppurpleL, rounded corners, minimum width=5.5cm, align=left] (ex1) at (6.2,3) {\emph{Can I top up here?}\\ \emph{Could you help me, please?}};
\node[draw=poppurple, fill=poppurpleL, rounded corners, minimum width=5.5cm, align=left] (ex2) at (6.2,1.5) {\emph{You should back up your files.}\\ \emph{You ought to log out after use.}};
\node[draw=poppurple, fill=poppurpleL, rounded corners, minimum width=5.5cm, align=left] (ex3) at (6.2,0) {\emph{You must pay before browsing.}\\ \emph{You mustn't download illegal files.}};
\draw[-{Stealth}, thick, popblue] (ability) -- (ex1);
\draw[-{Stealth}, thick, popgreen] (advice) -- (ex2);
\draw[-{Stealth}, thick, poporange] (rules) -- (ex3);
\end{tikzpicture}
\end{center}
\legende{Summary map of modal verbs: each function (ability/request, advice, rules) is matched with typical ICT-context sentences. Note that \emph{mustn't} expresses prohibition, not the absence of obligation.}
\end{popfigure}

\begin{attention}[``Mustn't'' does not mean ``don't have to'']
\emph{You mustn't use your phone} = it is \textbf{forbidden} (interdit). \emph{You don't have to print the form} = it is \textbf{not necessary} (pas nécessaire, vous pouvez le faire si vous voulez). Confusing these two changes the meaning completely — a classic trap in the grammar section of the exam.
\end{attention}

\courssub{Present perfect vs simple past: talking about ICT experiences}

\textbf{Intuition.} Both tenses talk about the past, but from different points of view. The \cle{simple past} (prétérit simple) presents a finished event at a definite time; the \cle{present perfect} (présent parfait) connects the past to the present (experience, recent result, or an action continuing until now).

\begin{definition}[Form and use]
\begin{itemize}
\item \textbf{Simple past}: subject + verb in the past (\emph{used, subscribed, bought}). Use it with finished time markers: \emph{yesterday, last week, in 2023, two days ago}. \emph{I subscribed to the monthly bundle \textbf{last week}.}
\item \textbf{Present perfect}: subject + \cle{have/has} + \cle{past participle} (participe passé : \emph{have used, has subscribed, have bought}). Use it with \emph{ever, never, already, just, yet, for, since} when no finished time is given: \emph{I have \textbf{never} used a cybercafé.} / \emph{She has had this phone \textbf{for} two years.} / \emph{We have known this provider \textbf{since} 2020.}
\end{itemize}
\end{definition}

\begin{methode}[Transforming a sentence into the present perfect]
\begin{enumerate}
\item Identify the subject and choose \emph{have} (I, you, we, they) or \emph{has} (he, she, it).
\item Find the past participle of the main verb: regular verbs add \emph{-ed} (\emph{use $\rightarrow$ used}); learn the irregular ones (\emph{buy $\rightarrow$ bought}, \emph{write $\rightarrow$ written}, \emph{go $\rightarrow$ gone}).
\item Add a suitable marker if required: \emph{for} + duration, \emph{since} + starting point, \emph{ever/never} for experience, \emph{just/already/yet} for recent news.
\item Remove any finished time expression (\emph{last week, yesterday}): they are incompatible with the present perfect.
\end{enumerate}
\end{methode}

\begin{exoresolu}[Simple past or present perfect?]
Complete: \emph{``I (never/use) video conferencing, but my brother (try) it last month during his training.''}
\tcblower
\textbf{Solution.} The first part expresses a life experience with \emph{never} and no definite time: present perfect $\rightarrow$ \emph{I \textbf{have never used} video conferencing}. The second part has a finished time marker (\emph{last month}): simple past $\rightarrow$ \emph{my brother \textbf{tried} it last month}. Full sentence: \emph{``I have never used video conferencing, but my brother tried it last month during his training.''}
\end{exoresolu}

\courssub{Linkers of cause, contrast and addition}

\textbf{Intuition.} Linkers (connecteurs logiques) are the ``glue'' of a composition. Without them, your text is a list of disconnected sentences; with them, the examiner sees a logical argument.

\begin{center}
\begin{tabularx}{\linewidth}{|l|l|X|}
\hline
\rowcolor{popblueL} \textbf{Function} & \textbf{Linkers} & \textbf{Example} \\
\hline
Cause / reason & because, since, as & \emph{I chose this network \textbf{because} it covers my village.} \\
\hline
Consequence & so, therefore & \emph{The bundle was cheap, \textbf{so} I bought two.} \\
\hline
Contrast & although, however, whereas & \emph{\textbf{Although} the monthly bundle is expensive, it is more economical for heavy users.} \\
\hline
Addition & moreover, furthermore, also & \emph{The package is affordable; \textbf{moreover}, it includes free night data.} \\
\hline
\end{tabularx}
\end{center}

\begin{exemplebox}[Contrast in context]
\emph{Although the monthly bundle is expensive, it is more economical for heavy users.} Notice the structure: \emph{although} + clause, comma, main clause. With \emph{however}, the contrast comes in a second sentence: \emph{The monthly bundle is expensive. \textbf{However}, it is more economical for heavy users.}
\end{exemplebox}

\begin{attention}[Punctuation traps]
\begin{itemize}
\item Never use \emph{although} and \emph{but} in the same sentence: \emph{Although it is expensive, but it is useful} is wrong.
\item \emph{Because} introduces a reason; \emph{so} introduces a result. Do not combine them: \emph{Because it rained, so I stayed home} is wrong.
\item \emph{However} and \emph{moreover} usually start a new sentence and are followed by a comma.
\end{itemize}
\end{attention}

\begin{savaistu}[Linkers and your Bac mark]
In the Baccalauréat technique writing paper, one explicit marking criterion is \cle{coherence} (organisation and linking of ideas). Examiners report that candidates who correctly use \emph{although, however, moreover, furthermore} systematically score higher on composition, even with the same ideas. Learning five linkers well is one of the cheapest ways to gain marks!
\end{savaistu}

\courssub{Question forms and polite requests in service situations}

\textbf{Intuition.} At a phone agency or cybercafé, a direct question (\emph{``What is the price?''}) can sound rude. English uses longer, indirect formulas to be polite.

\begin{definition}[Structures of polite questions]
\begin{itemize}
\item \textbf{Direct question}: Wh-word + auxiliary + subject + verb? \emph{How much \textbf{is} the weekly bundle?} / \emph{When \textbf{does} the offer expire?}
\item \textbf{Indirect (polite) question}: \emph{Could you tell me} + Wh-word + subject + verb (no inversion!): \emph{\textbf{Could you tell me} how much the weekly bundle \textbf{is}?} / \emph{\textbf{Could you tell me} when the offer \textbf{expires}?}
\item \textbf{Yes/no polite request}: \emph{Could you + verb...?} / \emph{Would you mind + -ing...?} \emph{Could you check my balance, please?} / \emph{Would you mind helping me configure my modem?}
\end{itemize}
\end{definition}

\begin{attention}[No inversion after ``Could you tell me'']
The most common mistake is keeping the question order: \emph{Could you tell me how much \textbf{is it}?} is wrong. After the polite introducer, use statement order: \emph{Could you tell me how much \textbf{it is}?}
\end{attention}

\courssub{Consolidation exercises}

\begin{exercice}[Gap-fill: modals and linkers]
Complete with \emph{can, could, should, mustn't, because, although, however, moreover}:
\begin{enumerate}
\item \emph{\ldots\ldots} you tell me where I can buy a SIM card, please?
\item You \ldots\ldots not use your phone during the examination; it is forbidden.
\item \ldots\ldots the 4G package is expensive, it is very fast.
\item The connection was slow, \ldots\ldots I went to the cybercafé.
\item You \ldots\ldots change your password regularly to stay safe online.
\end{enumerate}
\end{exercice}

\begin{corrige}[Exercice 1]
1. \emph{Could} (polite request). 2. \emph{must} (\emph{must not} = prohibition). 3. \emph{Although} (contrast at the start of the sentence). 4. \emph{so} (consequence). 5. \emph{should} (advice).
\end{corrige}

\begin{exercice}[Sentence transformation]
\begin{enumerate}
\item Rewrite in the present perfect: \emph{I used a computer for the first time.} (never $\rightarrow$ negative experience) — Begin: \emph{I have\ldots}
\item Make this question polite: \emph{How much does the monthly bundle cost?} — Begin: \emph{Could you tell me\ldots}
\item Transform with a comparative: \emph{The fibre connection is fast. The ADSL connection is not so fast.}
\end{enumerate}
\end{exercice}

\begin{corrige}[Exercice 2]
1. \emph{I have never used a computer.} (have + past participle \emph{used}, with \emph{never}). 2. \emph{Could you tell me how much the monthly bundle costs?} (statement order after the introducer, no inversion). 3. \emph{The fibre connection is faster than the ADSL connection.}
\end{corrige}

\begin{exercice}[Error correction]
Each sentence contains one mistake. Correct it.
\begin{enumerate}
\item \emph{This package is more cheaper than the other one.}
\item \emph{Although the data is expensive, but the speed is good.}
\item \emph{Could you tell me where is the nearest agency?}
\item \emph{I have subscribed to this bundle last Saturday.}
\end{enumerate}
\end{exercice}

\begin{corrige}[Exercice 3]
1. \emph{This package is \textbf{cheaper} than the other one.} (never \emph{more + -er}). 2. \emph{Although the data is expensive, the speed is good.} (remove \emph{but}). 3. \emph{Could you tell me where the nearest agency \textbf{is}?} (no inversion). 4. \emph{I subscribed to this bundle last Saturday} (simple past with a finished time) or \emph{I have subscribed to this bundle} (present perfect without the time marker).
\end{corrige}

\begin{aretenir}[Key points of this revision]
\begin{itemize}
\item Short adjectives: \emph{-er / the -est}; long adjectives: \emph{more / the most}. Never mix the two systems.
\item Modals + bare infinitive: \emph{can/could} (ability, request), \emph{should/ought to} (advice), \emph{must/mustn't} (obligation, prohibition).
\item Present perfect (\emph{have/has + past participle}) for experience and links to the present; simple past for finished events with a definite time.
\item Linkers (\emph{because, so, although, however, moreover, furthermore}) build coherence — a marked criterion at the Bac.
\item Polite questions: after \emph{Could you tell me\ldots}, use statement word order.
\end{itemize}
\end{aretenir}

\coursec{Listening: ICTs to enhance work and learning}

After revising key grammar structures in the previous section, we now turn our attention to \textbf{listening comprehension}. In today's digital world, understanding spoken English about technology is essential. Whether you are following an online tutorial, participating in a virtual meeting, or listening to a podcast about new apps, you need strategies to catch the main ideas and specific details. This section will train your ear through a realistic dialogue about how students, teachers, farmers, and traders in Cameroon use ICTs (Information and Communication Technologies) to improve their work and studies.

\courssub{Pre-listening: Brainstorming on ICT uses in Cameroon}

Before listening, it is crucial to activate your background knowledge. This helps you predict vocabulary and ideas, making the listening task much easier.

\begin{methode}[How to prepare for a listening task]
\begin{enumerate}
\item \textbf{Read the questions first.} Know what information you are hunting for (names, tools, purposes, benefits).
\item \textbf{Brainstorm vocabulary.} Think of words linked to the topic (e.g., \cle{smartphone}, \cle{WhatsApp}, \cle{mobile money}, \cle{e-learning}).
\item \textbf{Predict the context.} Who is speaking? Where? Why?
\item \textbf{During the first listening:} Focus on \textbf{global understanding} (who, what, general topic). Do not try to write everything.
\item \textbf{During the second listening:} Focus on \textbf{details}. Fill in your note-taking grid using \textbf{abbreviations} (e.g., \texttt{stdt} for student, \texttt{w/} for with, \texttt{b/c} for because).
\item \textbf{After listening:} Check your notes, complete missing info, and write a short summary.
\end{enumerate}
\end{methode}

\begin{exemplebox}[Brainstorming session: How do Cameroonians use ICTs?]
Work in small groups. Complete the mind map below with as many ideas as possible in 3 minutes.

\begin{popfigure}
\begin{tikzpicture}[
    mindmap,
    concept color=popblue,
    level 1 concept/.append style={sibling angle=90, level distance=4cm},
    level 2 concept/.append style={sibling angle=45, level distance=3cm},
    every node/.style={font=\small}
]
\node[concept] {ICT Uses in Cameroon}
    child[concept color=popgreen] {node[concept] {Students}
        child {node[concept] {e-learning platforms}}
        child {node[concept] {WhatsApp study groups}}
        child {node[concept] {Online tutorials (YouTube)}}
        child {node[concept] {Digital libraries}}}
    child[concept color=poporange] {node[concept] {Teachers}
        child {node[concept] {Virtual classrooms}}
        child {node[concept] {Sharing resources via email/Drive}}
        child {node[concept] {Online marking tools}}
        child {node[concept] {CRTV school programmes}}}
    child[concept color=poppurple] {node[concept] {Farmers}
        child {node[concept] {Checking market prices (SMS/Apps)}}
        child {node[concept] {Weather forecasts}}
        child {node[concept] {Mobile money for sales}}
        child {node[concept] {Agricultural advice via WhatsApp}}}
    child[concept color=poppink] {node[concept] {Traders}
        child {node[concept] {Online markets (Jumia, Kaymu)}}
        child {node[concept] {Mobile money transactions}}
        child {node[concept] {WhatsApp Business for orders}}
        child {node[concept] {Digital advertising}}};
\end{tikzpicture}
\legende{Mind map: ICT uses by different socio-professional groups in Cameroon}
\end{popfigure}
\end{exemplebox}

\textbf{Key vocabulary to remember:}
\begin{center}
\begin{tabularx}{\linewidth}{|l|X|l|}\hline
\textbf{English term} & \textbf{Definition / Context} & \textbf{French gloss} \\ \hline
e-learning & Learning via digital platforms, often at a distance & apprentissage en ligne \\ \hline
online course & A course delivered entirely on the internet & cours en ligne \\ \hline
tutorial & A step-by-step guide (video or text) to learn a skill & tutoriel \\ \hline
digital skills & Abilities to use digital tools effectively & compétences numériques \\ \hline
productivity & Efficiency in completing tasks & productivité \\ \hline
to download / to upload & To receive / send data from/to a server & télécharger / envoyer (vers un serveur) \\ \hline
search engine & Tool to find information online (Google, Bing) & moteur de recherche \\ \hline
virtual classroom & Online space where teachers and students meet live & classe virtuelle \\ \hline
mobile money & Financial service via phone (MTN MoMo, Orange Money) & argent mobile \\ \hline
\end{tabularx}
\end{center}

\courssub{Listening Task 1: Global Comprehension}

You will hear a dialogue between \textbf{three speakers} discussing how they use ICTs for work and learning. Listen once and answer the questions below.

\begin{exercice}[Global Comprehension Questions]
\begin{enumerate}
\item How many speakers are there? What are their roles (student, teacher, farmer, trader)?
\item What is the general topic of the conversation?
\item What is the purpose of the dialogue? (To complain? To share experiences? To advertise a product?)
\item Which ICT tools are mentioned? Tick the ones you hear: \quad $\square$ WhatsApp \quad $\square$ Zoom \quad $\square$ Facebook \quad $\square$ Mobile Money \quad $\square$ Moodle \quad $\square$ YouTube \quad $\square$ Google \quad $\square$ SMS
\end{enumerate}
\end{exercice}

\begin{corrige}[Exercice 1 -- Global Comprehension]
\begin{enumerate}
\item \textbf{Three speakers:} \textit{Amina} (university student), \textit{Mr. Njoya} (secondary school teacher), \textit{Papa Kofi} (cocoa farmer).
\item \textbf{Topic:} How they use ICTs (smartphones, apps, internet) to improve their studies, teaching, and farming business.
\item \textbf{Purpose:} To share experiences and show the benefits of digital tools in daily professional life.
\item \textbf{Tools mentioned:} WhatsApp, Zoom, Mobile Money, Moodle, YouTube, Google. (Facebook and SMS are \textbf{not} mentioned in this extract.)
\end{enumerate}
\end{corrige}

\courssub{Listening Task 2: Detailed Comprehension -- Note-taking Grid}

Listen a second time. Complete the grid below with precise information. Use abbreviations to write faster.

\begin{popfigure}
\begin{tikzpicture}
% Note-taking grid using TikZ matrix
\matrix (grid) [
    matrix of nodes,
    nodes in empty cells,
    nodes={anchor=center, minimum height=1.2cm, text width=3.2cm, align=center, font=\small},
    column sep=-\pgflinewidth,
    row sep=-\pgflinewidth,
    column 1/.style={nodes={fill=popblue!20, font=\bfseries\small, text width=2.2cm}},
    column 2/.style={nodes={fill=popcream}},
    column 3/.style={nodes={fill=popcream}},
    column 4/.style={nodes={fill=popcream}},
    column 5/.style={nodes={fill=popcream}},
    row 1/.style={nodes={fill=popblue, font=\bfseries\small, text=white, minimum height=1.4cm}},
] {
|[text width=2.2cm]| \textbf{Speaker} & \textbf{Role} & \textbf{ICT Tool(s) Used} & \textbf{Purpose} & \textbf{Benefit / Outcome} \\
Amina & & & & \\
Mr. Njoya & & & & \\
Papa Kofi & & & & \\
};
% Draw grid lines
\draw[popline, line width=1pt] (grid.north west) rectangle (grid.south east);
\foreach \i in {1,...,4} {
    \draw[popline] (grid-\i-1.north west) -- (grid-\i-5.north east);
}
\foreach \j in {1,...,5} {
    \draw[popline] (grid-1-\j.north west) -- (grid-4-\j.south west);
}
\end{tikzpicture}
\legende{Note-taking grid for detailed listening comprehension (to be completed during the 2nd listening)}
\end{popfigure}

\begin{exoresolu}[Completed Note-taking Grid (Teacher's Version)]
\textbf{Task:} Complete the grid based on the second listening.
\tcblower
\begin{center}
\begin{tabularx}{\linewidth}{|l|X|X|X|X|}\hline
\textbf{Speaker} & \textbf{Role} & \textbf{ICT Tool(s) Used} & \textbf{Purpose} & \textbf{Benefit / Outcome} \\ \hline
Amina & University student (Level 3, Law) & Moodle (university platform), WhatsApp group, YouTube tutorials & Follow online courses, discuss cases with classmates, watch legal tutorials & Flexibility: studies while doing internship; access to recorded lectures anytime \\ \hline
Mr. Njoya & Secondary school teacher (Physics) & Zoom, Google Drive, WhatsApp Broadcast list & Teach virtual classes during holidays, share lesson notes \& videos, send homework reminders & Reaches students at home; saves travel time; builds digital skills for students \\ \hline
Papa Kofi & Cocoa farmer (Bertoua region) & Smartphone, Mobile Money (MoMo), WhatsApp, Market price app (Esoko) & Check cocoa prices before selling, receive payments via MoMo, join farmer coop group & Better prices (no middlemen cheating); instant payment; advice on plant diseases \\ \hline
\end{tabularx}
\end{center}
\end{exoresolu}

\courssub{Post-listening: Oral Report}

Now, use your completed grid to give a short oral report (1 minute per speaker). Follow this structure:

\begin{methode}[Structure for a short oral report]
\begin{enumerate}
\item \textbf{Introduction:} "I'm going to talk about [Name], a [role] from [place]."
\item \textbf{Tools:} "He/She uses [tool 1] and [tool 2]..."
\item \textbf{Purpose:} "...in order to [verb phrase] / for [noun phrase]."
\item \textbf{Benefit:} "Thanks to this, he/she can [result]. This improves his/her [productivity / income / learning]."
\item \textbf{Conclusion:} "This shows how ICTs can really help [category of people] in Cameroon."
\end{enumerate}
\end{methode}

\begin{exemplebox}[Model Oral Report: Amina]
\textbf{Introduction:} "I'm going to talk about Amina, a university law student in Yaoundé."

\textbf{Tools:} "She uses Moodle, her university's e-learning platform, a WhatsApp group with her classmates, and YouTube tutorials."

\textbf{Purpose:} "She uses Moodle to follow online courses and download lecture notes. The WhatsApp group is for discussing case studies and sharing summaries. YouTube helps her understand difficult legal concepts through tutorials."

\textbf{Benefit:} "Thanks to these tools, Amina can study flexibly. She does an internship during the day and watches recorded lectures at night. She says her digital skills have improved a lot."

\textbf{Conclusion:} "This shows how e-learning helps Cameroonian students balance work and studies."
\end{exemplebox}

\courssub{Grammar \& Vocabulary Focus from the Listening}

\begin{definition}[E-learning]
\textbf{E-learning} (electronic learning) refers to \textbf{learning through digital tools and the internet, often at a distance}. It includes online courses, virtual classrooms, video tutorials, and digital resources. It allows learners to study at their own pace, anywhere.
\begin{itemize}
\item \textit{Example:} "Amina follows an \textbf{e-learning} programme on Moodle."
\item \textit{Related terms:} \cle{online learning}, \cle{distance learning}, \cle{digital education}.
\end{itemize}
\end{definition}

\begin{attention}[Common Mistake: \textbf{Learn} vs \textbf{Teach}]
Many Francophone learners confuse these two verbs because French uses \textit{apprendre} for both.
\begin{center}
\begin{tabular}{|l|l|l|}\hline
\textbf{Verb} & \textbf{Meaning} & \textbf{Example} \\ \hline
\textbf{LEARN} & Acquire knowledge/skill (apprendre) & "Students \textbf{learn} physics on Zoom." \\ \hline
\textbf{TEACH} & Give knowledge/skill to others (enseigner) & "Mr. Njoya \textbf{teaches} physics on Zoom." \\ \hline
\end{tabular}
\end{center}
\textbf{Rule:} The \textbf{teacher teaches}, the \textbf{student learns}. Never say "The teacher learns the students" $\rightarrow$ \textbf{The teacher \emph{teaches} the students.}
\end{attention}

\begin{exemplebox}[Vocabulary in Context: Papa Kofi the Farmer]
\textbf{Original sentence from the dialogue:}
"The farmer in Bafoussam uses a smartphone to check maize prices before selling."

\textbf{Analysis:}
\begin{itemize}
\item \textbf{check} = verify, look up (vérifier, consulter)
\item \textbf{maize} = corn (maïs)
\item \textbf{before selling} = \textit{before} + \textbf{-ing} form (avant de vendre)
\item \textbf{smartphone} = téléphone intelligent
\end{itemize}
\textbf{Your turn:} Write a similar sentence for a trader.
\textit{Example:} "The trader in Douala uses WhatsApp Business to receive orders from customers."
\end{exemplebox}

\courssub{Extension: Distance Learning in Cameroon -- Discussion}

\begin{savaistu}[Distance Learning in Cameroon: A Brief History]
\begin{itemize}
\item \textbf{2020 (COVID-19):} Schools closed nationwide. The Ministry of Secondary Education (MINESEC) and Ministry of Higher Education (MINESUP) launched emergency remote teaching.
\item \textbf{CRTV "École à la maison":} Daily radio and TV lessons for exam classes (3\textsuperscript{e}, Terminale, Probatoire, Bac). Reached students without internet.
\item \textbf{University platforms:} State universities (Yaoundé I, II, Douala, Buea, Dschang, Ngaoundéré, Maroua) adopted \textbf{Moodle} and \textbf{Google Classroom}. Some created their own portals (e.g., \textit{UNIV-YAOUNDE I e-campus}).
\item \textbf{Private initiatives:} \textit{Kamerbigbang}, \textit{Edtech Cameroon}, \textit{ScholarX} offer online courses and career guidance.
\item \textbf{Current challenge:} Internet access in rural areas, cost of data, electricity cuts, and low digital literacy among some teachers/parents.
\end{itemize}
\end{savaistu}

\begin{exercice}[Discussion Questions -- Work in Pairs]
Discuss the following questions in English. Use the vocabulary from this section.
\begin{enumerate}
\item Did you follow CRTV lessons or university online courses during COVID-19? What was your experience?
\item What are the \textbf{advantages} and \textbf{disadvantages} of e-learning for a student in a rural area (e.g., Mokolo, Batouri)?
\item How can \textbf{mobile money} help a small business owner (trader, farmer, tailor) grow their activity?
\item Imagine you are the Minister of Posts and Telecommunications. What \textbf{three measures} would you take to improve digital access for education?
\end{enumerate}
\end{exercice}

\begin{corrige}[Exercice 3 -- Suggested Answers for Discussion]
\begin{enumerate}
\item \textit{Personal answers. Example:} "I followed CRTV lessons for Probatoire. It was helpful but hard to ask questions. The TV signal was sometimes weak in my village."
\item \textbf{Advantages:} Access to quality resources, no travel cost, flexible schedule, digital skills development. \textbf{Disadvantages:} Poor internet connectivity, high data cost, no electricity, lack of devices, isolation, need for self-discipline.
\item Mobile money allows: instant payments (no cash risk), savings, access to micro-loans, paying suppliers remotely, receiving money from diaspora customers.
\item \textit{Example measures:} (1) Subsidize data bundles for students/teachers. (2) Build more community telecentres with solar power in rural areas. (3) Train teachers in digital pedagogy (how to teach online effectively).
\end{enumerate}
\end{corrige}

\courssub{Integration Mini-Project: ICT Profile Interview}

\begin{methode}[Mini-Project Steps]
\begin{enumerate}
\item \textbf{Prepare} 5 questions about ICT use (tools, frequency, purpose, benefits, challenges).
\item \textbf{Interview} a classmate, family member, or neighbour (in English if possible, otherwise translate key answers).
\item \textbf{Fill} a note-taking grid like the one used in Listening Task 2.
\item \textbf{Write} a 150-word profile: "How [Name] uses ICTs to work/learn better."
\item \textbf{Present} your profile to the class (2 minutes).
\end{enumerate}
\end{methode}

\begin{exemplebox}[Sample Interview Questions]
\begin{enumerate}
\item What is your main occupation? (student, teacher, farmer, trader, other)
\item Which digital tools do you use regularly? (WhatsApp, Facebook, Mobile Money, Google, Zoom, specialized apps...)
\item How do these tools help you in your work or studies? Give a concrete example.
\item What difficulties do you face? (network, cost, skills, electricity...)
\item What new digital skill would you like to learn?
\end{enumerate}
\end{exemplebox}

\begin{aretenir}[Key Takeaways: Listening to ICT Dialogues]
\begin{itemize}
\item \textbf{Pre-listening:} Predict vocabulary and context (brainstorming).
\item \textbf{1st listening:} Global idea (who, what, why). Don't write everything.
\item \textbf{2nd listening:} Details $\rightarrow$ use a \textbf{structured grid} (Who / Tool / Purpose / Benefit).
\item \textbf{Note-taking:} Use abbreviations, symbols (\&, +, $\rightarrow$, $\uparrow$), keywords only.
\item \textbf{Post-listening:} Oral report $\rightarrow$ Structure: Intro $\rightarrow$ Tools $\rightarrow$ Purpose $\rightarrow$ Benefit $\rightarrow$ Conclusion.
\item \textbf{Vocabulary:} Master the list: \cle{e-learning}, \cle{virtual classroom}, \cle{mobile money}, \cle{download/upload}, \cle{search engine}, \cle{digital skills}, \cle{productivity}, \cle{tutorial}.
\item \textbf{Grammar trap:} \textbf{Teach} (enseigner) $\neq$ \textbf{Learn} (apprendre).
\end{itemize}
\end{aretenir}

\coursec{Reading: ICTs as recreational resources}

In the previous section, we listened to texts about using ICTs to \emph{enhance work and learning}: computers and phones as tools for study and business. But life is not only work! In this section, we move to the other side of the screen: ICTs as \cle{recreational resources} (des ressources de loisirs). You will read an authentic-style text about how young Cameroonians use ICTs for fun, learn key vocabulary, answer comprehension questions like at the Probatoire and Baccalauréat technique, and exploit real data about screen time.

\courssub{Pre-reading: class survey}

Before reading, let us activate what you already know. Look at this question and discuss it in pairs:

\begin{center}
\emph{``What do you do for fun with your phone or computer?''}
\end{center}

Possible answers:
\begin{itemize}
\item I listen to \cle{music} on my phone.
\item I watch \cle{films} and series.
\item I play \cle{online games} (des jeux en ligne).
\item I chat with friends on \cle{social media} (WhatsApp, Facebook, TikTok).
\item I watch football matches through \cle{sports streaming}.
\end{itemize}

\textbf{Class survey task:} Walk around the class, ask five classmates the question, and complete this table. Then report: \emph{``Three students out of five use their phones to listen to music.''}

\begin{center}
\begin{tabularx}{\linewidth}{|l|X|X|}\hline
\rowcolor{popblueL} \textbf{Name} & \textbf{Favourite activity} & \textbf{Hours per day} \\ \hline
1. & & \\ \hline
2. & & \\ \hline
3. & & \\ \hline
4. & & \\ \hline
5. & & \\ \hline
\end{tabularx}
\end{center}

\courssub{Vocabulary: entertainment and leisure}

Here are the key words of this lesson. Learn them with their French glosses.

\begin{definition}[Key vocabulary]
\begin{itemize}
\item \cle{entertainment} (n.): amusement, diversion — \emph{divertissement}.
\item \cle{leisure} (n.): free time, time for relaxation — \emph{loisir, temps libre}.
\item \cle{streaming} (n.): watching or listening to content directly on the internet without downloading — \emph{diffusion en continu}.
\item \cle{social network} (n.): a platform to connect with people online (Facebook, WhatsApp) — \emph{réseau social}.
\item \cle{online game} (n.): a game played on the internet — \emph{jeu en ligne}.
\item \cle{to chat} (v.): to exchange messages informally — \emph{discuter, tchatter}.
\item \cle{to post / to share} (v.): to publish / to send content to others — \emph{publier / partager}.
\item \cle{addiction} (n.): the impossibility to stop a habit — \emph{addiction, dépendance}.
\item \cle{cybercafé} (n.): a shop where people pay to use computers and internet — \emph{cybercafé}.
\end{itemize}
\end{definition}

\begin{definition}[Screen time]
\cle{Screen time} is \textbf{the amount of time spent in front of a screen} (phone, TV, computer or tablet) during a day. Example: \emph{``My screen time yesterday was five hours.''} Doctors recommend limiting recreational screen time to protect the eyes and the brain.
\end{definition}

\courssub{Reading text: How young Cameroonians use ICTs at home and in the community}

Read the text silently, then aloud with your teacher. Pay attention to the main idea of each paragraph.

\begin{exemplebox}[Reading text]
\textbf{How young Cameroonians use ICTs at home and in the community}

\textbf{Paragraph 1.} In Cameroon today, the mobile phone is the king of leisure. In towns like Douala, Yaoundé and Bamenda, most young people own a smartphone. At home, they use it to listen to music, watch videos on YouTube and chat with friends on WhatsApp and Facebook. For many students, the phone is the first thing they touch in the morning and the last thing they see at night.

\textbf{Paragraph 2.} In the community, ICT leisure has its own special places. \emph{Video clubs} are small rooms where young people pay a hundred francs to watch action films or Nigerian series on a big television screen. \emph{Gaming centres} offer PlayStation games, especially football games like FIFA, where friends compete against each other. On weekends, \emph{football viewing centres} are full: supporters pay to watch the English Premier League or the Indomitable Lions on satellite television, shouting and celebrating together.

\textbf{Paragraph 3.} \emph{Community radio} also plays an important role. In rural areas where internet is expensive, young people listen to music programmes, football commentaries and phone-in shows on local radio stations. Some even send text messages to the radio to greet their friends or request a song.

\textbf{Paragraph 4.} However, this digital leisure has two faces. On one hand, ICTs bring relaxation, information and social links: a student in Bafoussam can laugh with a cousin in France. On the other hand, excessive screen time can lead to addiction, poor school results and exposure to fake news. The challenge for Cameroonian youth is to enjoy ICTs without becoming their slave.
\end{exemplebox}

\courssub{Comprehension tasks}

\begin{methode}[How to answer a comprehension question]
\begin{enumerate}
\item Read the question carefully and underline the key words.
\item Find the paragraph and the exact line in the text that contains the answer.
\item \textbf{Quote the line} (cite la ligne du texte) or refer to it: \emph{``According to paragraph 2...''}.
\item \textbf{Reformulate with your own words} (reformule avec tes propres mots): do not copy long sentences from the text.
\item For True/False questions, always write True or False AND the justification from the text.
\end{enumerate}
\end{methode}

\begin{attention}[Frequent mistake at the Bac]
Many candidates answer \emph{``Yes''} or \emph{``No''} (or \emph{``True''} / \emph{``False''}) \textbf{without justifying with the text}. At the Probatoire and Baccalauréat technique, a True/False answer without a quotation or a clear reference to the text usually receives \textbf{zero}, even if the answer is correct! Always add: \emph{``...because the text says that...''} or \emph{``(paragraph 3, line 2)''}.
\end{attention}

\begin{exoresolu}[True/False with justification]
\textbf{Statement:} ``Young Cameroonians watch football matches only at home.'' True or False? Justify.
\tcblower
\textbf{Solution:} \emph{False.} According to paragraph 2, on weekends ``football viewing centres are full: supporters pay to watch the English Premier League or the Indomitable Lions on satellite television.'' This means that many young people watch football in community viewing centres, not only at home.
\end{exoresolu}

\begin{exercice}[Comprehension questions]
Answer in complete sentences and justify from the text.
\begin{enumerate}
\item \textbf{True/False:} a) Community radio is important where internet is expensive. b) Video clubs are free of charge. c) The writer thinks ICTs are only dangerous.
\item \textbf{Matching:} Match the words to their definitions: \emph{video club, gaming centre, viewing centre, community radio} — a) a place to watch football on satellite TV; b) a local station with music and phone-in shows; c) a room to watch films for a small fee; d) a place to play PlayStation games.
\item \textbf{Reference:} In paragraph 4, what does the word ``their'' (``without becoming their slave'') refer to?
\item \textbf{Main ideas:} Write one sentence giving the main idea of each paragraph (1 to 4).
\end{enumerate}
\end{exercice}

\begin{corrige}[Exercice: Comprehension questions]
\begin{enumerate}
\item a) \emph{True} — paragraph 3: ``In rural areas where internet is expensive, young people listen to... local radio stations.'' b) \emph{False} — paragraph 2: young people ``pay a hundred francs'' to watch films. c) \emph{False} — paragraph 4 says ICTs also ``bring relaxation, information and social links''.
\item video club $\rightarrow$ c; gaming centre $\rightarrow$ d; viewing centre $\rightarrow$ a; community radio $\rightarrow$ b.
\item ``Their'' refers to \emph{ICTs} (information and communication technologies).
\item Paragraph 1: the smartphone is the main leisure tool at home. Paragraph 2: community places offer digital leisure (video clubs, gaming and viewing centres). Paragraph 3: community radio entertains rural youth. Paragraph 4: digital leisure has advantages and dangers.
\end{enumerate}
\end{corrige}

\courssub{Critical thinking: advantages vs dangers}

The text ends with a balanced idea: ICTs have ``two faces''. Let us organise this debate clearly.

\begin{popfigure}
\begin{center}
\begin{tikzpicture}[
box/.style={rounded corners=3pt, draw=#1, fill=#1!12, text width=5.6cm, minimum height=0.85cm, align=center, font=\small},
title/.style={rounded corners=3pt, draw=#1, fill=#1!35, text width=5.6cm, align=center, font=\small\bfseries}]
\node[title=popgreen] (A) at (0,0) {ADVANTAGES (les avantages)};
\node[title=poporange] (D) at (7.2,0) {DANGERS (les dangers)};
\node[box=popgreen] at (0,-1.1) {Relaxation and entertainment after work or school};
\node[box=popgreen] at (0,-2.2) {Access to information and culture (music, films, news)};
\node[box=popgreen] at (0,-3.3) {Social links: chatting with family and friends far away};
\node[box=popgreen] at (0,-4.4) {Learning languages and new skills through videos};
\node[box=poporange] at (7.2,-1.1) {Addiction and excessive screen time};
\node[box=poporange] at (7.2,-2.2) {Fake news and wrong information};
\node[box=poporange] at (7.2,-3.3) {Cybercrime: scams, theft of personal data};
\node[box=poporange] at (7.2,-4.4) {Exam failure: less time for studies and sleep};
\draw[-{Stealth}, thick, popdark] (3.0,-5.2) -- (4.2,-5.2) node[midway, below, font=\small\itshape]{balance!};
\end{tikzpicture}
\end{center}
\legende{Advantages versus dangers of recreational ICTs: a balanced view. A good citizen enjoys the advantages and protects himself against the dangers.}
\end{popfigure}

\textbf{Pair work discussion:} Choose one advantage and one danger. Give a concrete example from your quarter or your school. Use the structure: \emph{``In my neighbourhood, ... For example, ... That is why I think...''}

\courssub{Data exploitation: screen time among Cameroonian youth}

At the Baccalauréat technique, you may have to read and comment on a chart. Observe the bar chart below, drawn from a class survey of 100 students in a technical high school in Yaoundé.

\begin{popfigure}
\begin{center}
\begin{tikzpicture}
\begin{axis}[
ybar, bar width=0.7cm,
width=0.85\linewidth, height=6.5cm,
ylabel={Average daily screen time (hours)},
symbolic x coords={Social media, Games, Music, Studies},
xtick=data,
ymin=0, ymax=4.5,
ytick={0,1,2,3,4},
nodes near coords,
every node near coord/.append style={font=\small},
ymajorgrids=true, grid style={popline!40},
axis line style={popdark},
]
\addplot[fill=popblue!70, draw=popblue] coordinates {(Social media,3.2) (Games,1.8) (Music,1.4) (Studies,1.1)};
\end{axis}
\end{tikzpicture}
\end{center}
\legende{Average daily screen time of students by activity (survey of 100 students, technical high school, Yaoundé, 2026). Social media dominates with 3,2 hours per day, while studies represent only 1,1 hour.}
\end{popfigure}

\begin{exemplebox}[Exploiting the data: a model sentence]
\emph{``According to the chart, the average daily screen time for social media is 3,2 hours, which is the highest of all activities. Studies receive only 1,1 hour on average. This shows that recreational screen time is much higher than study screen time among young Cameroonians.''}
\end{exemplebox}

\begin{methode}[How to comment on a chart]
\begin{enumerate}
\item \textbf{Introduce} the chart: what does it show? (source, place, population).
\item \textbf{Quote figures}: the highest bar, the lowest bar, a comparison (\emph{twice as much as...}).
\item \textbf{Interpret}: what does it mean for young people's lives?
\item \textbf{Conclude} with your opinion or a recommendation.
\end{enumerate}
\end{methode}

\begin{exercice}[Chart commentary]
Using the bar chart above, answer in full sentences:
\begin{enumerate}
\item Which activity has the longest screen time? Quote the figure.
\item Compare games and studies: which one is higher, and by how much?
\item Write a short paragraph (5 lines) commenting on the chart, following the method above.
\end{enumerate}
\end{exercice}

\begin{corrige}[Exercice: Chart commentary]
\begin{enumerate}
\item Social media has the longest screen time: 3,2 hours per day on average.
\item Games (1,8 hours) are higher than studies (1,1 hours): students spend 0,7 hour more on games than on studies.
\item Model paragraph: \emph{``The chart shows the average daily screen time of 100 students in a technical high school in Yaoundé. Social media comes first with 3,2 hours, followed by games with 1,8 hours and music with 1,4 hours. Studies come last with only 1,1 hours. This means that students use their screens much more for fun than for learning. In my opinion, young people should reduce their social media time and increase their study time to succeed at the Baccalauréat technique.''}
\end{enumerate}
\end{corrige}

\begin{savaistu}[Digital leisure in Cameroonian neighbourhoods]
Did you know? Even with smartphones everywhere, \emph{video clubs} and \emph{gaming centres} remain very popular leisure places in Cameroonian quarters. In many neighbourhoods of Douala and Yaoundé, a PlayStation game costs about 100 to 200 FCFA, and football viewing centres become real stadiums when the Indomitable Lions play! These places show that in Cameroon, digital leisure is often a \textbf{collective} experience: people prefer to play, watch and celebrate \emph{together} rather than alone.
\end{savaistu}

\begin{aretenir}[What to remember]
\begin{itemize}
\item Key vocabulary: entertainment, leisure, streaming, social network, online game, to chat, to post/share, screen time, addiction, cybercafé.
\item \cle{Screen time} = the amount of time spent in front of a screen.
\item Reading skill: always justify True/False answers with a quotation or reference from the text.
\item Chart skill: introduce, quote figures, interpret, conclude.
\item Critical mind: recreational ICTs bring relaxation, information and social links, but also addiction, fake news, cybercrime and exam failure. Enjoy them with balance!
\end{itemize}
\end{aretenir}

\coursec{Writing: composition on recreational use of ICTs}

In the previous sections, you read texts about how young people use ICTs as recreational resources at home and in the community: playing video games, watching films online, chatting on social networks, listening to music. You have collected many ideas and much vocabulary. Now it is time to \emph{produce} your own text. At the Probatoire and the Baccalauréat technique, you will be asked to write a composition of about 250 to 300 words on a topic like this one. This section teaches you, step by step, how to plan, organise, write and correct a good composition.

\courssub{Understanding the task: analysing a Bac-type topic}

Before writing a single sentence, you must \cle{analyse the topic} (analyser le sujet). Many candidates lose marks because they write about a different subject from the one given.

\begin{exemplebox}[A typical Bac technique topic]
\emph{« Write a 250--300 word composition on the advantages and disadvantages of using ICTs for recreation in your community. »}
\end{exemplebox}

Let us analyse this topic together. Three questions must be answered:

\begin{enumerate}
\item \textbf{What is the subject?} Using ICTs (smartphones, computers, internet, games consoles) for \emph{recreation} (leisure, fun / loisir, divertissement) --- not for work or study.
\item \textbf{What is the task?} Discuss \emph{both} advantages \emph{and} disadvantages. If you write only about advantages, you answer only half of the question.
\item \textbf{What are the limits?} Between 250 and 300 words, and the context is \emph{your community} (your town, your neighbourhood, your school / votre ville, votre quartier, votre école).
\end{enumerate}

\begin{attention}[Do not copy the topic as your introduction]
A very frequent error at the Bac is to copy the subject word for word as the first sentence of the composition. The examiner sees immediately that you have not understood the task. Instead, \cle{reformulate} (reformuler) the topic in your own words. For example, do not write « The advantages and disadvantages of using ICTs for recreation in my community are many », but rather « In Douala today, almost every young person relaxes with a phone or a computer. Is this new habit good or bad for our community? »
\end{attention}

\courssub{Planning: brainstorming and building an outline}

Good writers never start writing immediately. They \cle{brainstorm} (remue-méninges) first: they write down quickly all the ideas that come to their mind, in English or even in French, without judging them.

\begin{exemplebox}[Brainstorming on our topic]
\textbf{Advantages:} relaxation after work (détente après le travail), video games develop reflexes (les jeux vidéo développent les réflexes), social networks keep friends connected (les réseaux sociaux maintiennent les amis connectés), music and films are cheap entertainment (la musique et les films sont un divertissement bon marché), learning English through films (apprendre l'anglais grâce aux films), community youth centres with computers (centres de jeunesse communautaires avec ordinateurs).\\
\textbf{Disadvantages:} addiction to games (dépendance aux jeux), less time for studies (moins de temps pour les études), eye problems (problèmes de vue), expensive data bundles (forfaits data coûteux), cyberbullying (cyberharcèlement), less physical sport (moins de sport physique), bad content for children (contenus inappropriés pour les enfants).
\end{exemplebox}

Then you \cle{select} the best 3 or 4 ideas (two advantages, two disadvantages, for example) and you build an \cle{outline} (plan détaillé): the skeleton of your composition.

\begin{methode}[Building an outline in five steps]
\begin{enumerate}
\item Read the topic twice and underline the key words (lisez le sujet deux fois et soulignez les mots-clés).
\item Brainstorm for five minutes: list every idea (faites un remue-méninges pendant cinq minutes : notez toutes les idées).
\item Choose 3 or 4 strong ideas that you can illustrate with an example from your community.
\item Order the ideas: advantages first, then disadvantages (or the reverse).
\item Note one linker and one example next to each idea.
\end{enumerate}
\end{methode}

\courssub{The structure of a composition}

A Bac composition has three parts, like a funnel: it starts wide (general idea), narrows down to your arguments, and opens again at the end with your conclusion.

\begin{popfigure}
\begin{center}
\begin{tikzpicture}[font=\small, scale=0.9]
% Funnel shape - four layers
\fill[popblueL] (-4.5,4.5) -- (4.5,4.5) -- (3,3) -- (-3,3) -- cycle;
\fill[poporangeL] (-3,3) -- (3,3) -- (2,1.5) -- (-2,1.5) -- cycle;
\fill[popgreenL] (-2,1.5) -- (2,1.5) -- (1.5,-1.5) -- (-1.5,-1.5) -- cycle;
\fill[poppinkL] (-1.5,-1.5) -- (1.5,-1.5) -- (2.5,-3) -- (-2.5,-3) -- cycle;
% Labels inside
\node at (0,3.8) {\textbf{INTRODUCTION}};
\node at (0,3.1) {\textbf{Hook + Thesis statement}};
\node at (0,2.2) {\textbf{BODY: paragraph 1}};
\node at (0,1.5) {\textbf{BODY: paragraph 2}};
\node at (0,0.8) {\textbf{BODY: paragraph 3}};
\node at (0,-0.5) {\textbf{CONCLUSION}};
\node at (0,-1.2) {\textbf{Opinion + Recommendation}};
% Side annotations
\node[right, text=popdark, align=left] at (5,3.8) {Hook: a question or a\\ general fact to attract\\ the reader (accroche)};
\node[right, text=popdark, align=left] at (3.5,3.1) {Your main idea,\\ your answer to the topic (thèse)};
\node[right, text=popdark, align=left] at (2.5,1.5) {One idea + one example\\ + one linker per paragraph};
\node[right, text=popdark, align=left] at (3,-1.2) {Your opinion +\\ a recommendation (avis + conseil)};
\draw[->, popmuted, thick] (4.8,3.8) -- (4,3.8);
\draw[->, popmuted, thick] (3.3,3.1) -- (2.5,3.1);
\draw[->, popmuted, thick] (2.3,1.5) -- (2,1.5);
\draw[->, popmuted, thick] (2.8,-1.2) -- (2.2,-1.2);
\end{tikzpicture}
\end{center}
\legende{The funnel structure of a composition: from a general hook to the thesis, through the body paragraphs (one main idea each), to the conclusion with your personal opinion.}
\end{popfigure}

\begin{definition}[The three parts of a composition]
The \cle{introduction} presents the subject with a \emph{hook} (une accroche: a question, a surprising fact) and ends with the \emph{thesis statement} (la thèse: your main idea). The \cle{body} (le développement) contains two or three paragraphs, each developing one argument with an example. The \cle{conclusion} gives your personal opinion and often a recommendation.
\end{definition}

\begin{propriete}[Golden rule of paragraphing]
Each body paragraph contains \textbf{one and only one main idea}, introduced by a topic sentence and connected to the rest of the text by a \cle{linker} (connecteur logique: \emph{firstly, moreover, however, on the other hand, in addition}). Two different ideas in the same paragraph confuse the reader and cost marks in organisation.
\end{propriete}

\begin{methode}[How to build a body paragraph: T-E-E-L]
\begin{enumerate}
\item \textbf{T -- Topic sentence:} state your idea clearly. \emph{Firstly, ICTs offer cheap entertainment to young people.}
\item \textbf{E -- Explanation:} explain what you mean. \emph{A small data bundle gives access to hours of music and films.}
\item \textbf{E -- Example:} give a concrete example from your community. \emph{For instance, in Yaoundé, many students watch comedies on their phones after school instead of going to the cinema.}
\item \textbf{L -- Linking sentence:} prepare the next paragraph. \emph{However, this entertainment also has a dark side.}
\end{enumerate}
\end{methode}

\courssub{Using lesson vocabulary and grammar}

Your composition must reuse the vocabulary of this lesson and the grammar you revised. Here is a toolbox.

\begin{center}
\begin{tabularx}{\linewidth}{|l|X|}
\hline
\rowcolor{popblueL} \textbf{Function} & \textbf{Useful expressions (Expressions utiles)} \\
\hline
Adding ideas (Ajouter des idées) & firstly, secondly, moreover, in addition, furthermore \\
\hline
Contrasting (Contraster / Nuancer) & however, on the other hand, whereas, although, but \\
\hline
Giving examples (Donner des exemples) & for example, for instance, such as \\
\hline
Comparing (Comparer) & more relaxing than, cheaper than, the most popular \\
\hline
Advising / Recommending - Modals (Conseiller - Modaux) & should, ought to, must, mustn't, had better \\
\hline
Concluding (Conclure) & in conclusion, to sum up, all in all \\
\hline
Lesson vocabulary (Vocabulaire de la leçon) & screen time (temps d'écran), data bundle (forfait data), social network (réseau social), video game (jeu vidéo), streaming (streaming), leisure (loisir), entertainment (divertissement), addiction (dépendance), to relax (se détendre), recreational (récréatif) \\
\hline
\end{tabularx}
\end{center}

\begin{exemplebox}[A model conclusion]
\emph{« In conclusion, although ICTs offer entertainment, students should limit their screen time to succeed at school. »}\\
Notice the structure: a concluding linker (\emph{in conclusion}), a concession (\emph{although}), and a recommendation with a modal (\emph{should}).
\end{exemplebox}

\begin{attention}[Common grammar traps (Pièges grammaticaux fréquents)]
\begin{itemize}
\item \emph{advices} $\rightarrow$ wrong: \emph{advice} is uncountable (indénombrable). Write \emph{some advice} or \emph{a piece of advice}.
\item \emph{more cheaper} $\rightarrow$ wrong: write \emph{cheaper} or \emph{less expensive} (pas de double comparatif).
\item \emph{must to limit} $\rightarrow$ wrong: after a modal, use the base verb (base verbale sans \emph{to}): \emph{must limit}, \emph{should use}.
\item Do not mix tenses: if you describe a general situation, stay in the present simple (présent simple).
\item \emph{informations} $\rightarrow$ wrong: \emph{information} is uncountable. Write \emph{some information}.
\end{itemize}
\end{attention}

\courssub{Guided writing: from a collective paragraph to a full draft}

Let us first draft one body paragraph \emph{together}, following the T-E-E-L method, on the idea « too much screen time harms studies ».

\begin{exoresolu}[Writing a body paragraph together]
Build a body paragraph on the idea: « recreational ICTs can harm school results ».
\tcblower
\textbf{Topic sentence:} On the other hand, spending too much time on recreational ICTs can seriously harm students' results.\\
\textbf{Explanation:} When a teenager plays video games or chats on WhatsApp late at night, he sleeps less and cannot concentrate in class the next morning.\\
\textbf{Example:} For instance, in many quarters of Bafoussam, some students fail their exams because they spend their evenings in game centres instead of doing their homework.\\
\textbf{Linking sentence:} Therefore, young people must learn to control their screen time.\\
\textbf{Full paragraph:} \emph{On the other hand, spending too much time on recreational ICTs can seriously harm students' results. When a teenager plays video games or chats on WhatsApp late at night, he sleeps less and cannot concentrate in class the next morning. For instance, in many quarters of Bafoussam, some students fail their exams because they spend their evenings in game centres instead of doing their homework. Therefore, young people must learn to control their screen time.}
\end{exoresolu}

Now it is your turn. Write your \cle{full draft} (brouillon complet) individually, following your outline. Do not count words while writing; count them at the end and adjust. \textbf{Time management tip:} In the exam, spend 10 minutes planning, 25 minutes writing, 5 minutes checking.

\begin{savaistu}[How is the composition marked at the Bac technique?]
At the Baccalauréat technique, your composition is marked on four criteria: the \textbf{content} (did you answer the question with relevant ideas?), the \textbf{organisation} (clear paragraphs, linkers, logical order), the \textbf{language} (correct grammar and rich vocabulary) and the \textbf{mechanics} (spelling, punctuation, capital letters). A composition with brilliant ideas but poor spelling cannot obtain full marks: all four criteria count!
\end{savaistu}

\courssub{Peer correction with a checklist}

Before giving your final copy to the teacher, exchange your draft with a classmate and correct each other's work with this \cle{checklist} (liste de vérification).

\begin{center}
\begin{tabularx}{\linewidth}{|l|X|}
\hline
\rowcolor{popgreenL} \textbf{Criterion (Critère)} & \textbf{Questions to ask (Questions à se poser)} \\
\hline
Content (Contenu) & Does the text answer the topic? Are there advantages \emph{and} disadvantages? Are examples from the community? \\
\hline
Organisation (Organisation) & Is there an introduction with a hook and a thesis? One idea per paragraph? Linkers? A conclusion with an opinion? \\
\hline
Vocabulary (Vocabulaire) & Are at least five lesson words used correctly (screen time, data bundle, recreational\ldots)? \\
\hline
Grammar (Grammaire) & Present simple for general facts? Correct comparatives? Modals followed by the base verb? \\
\hline
Mechanics (Forme) & Spelling, capital letters at the beginning of sentences, full stops, commas? Between 250 and 300 words? \\
\hline
\end{tabularx}
\end{center}

\begin{exercice}[Your turn to write]
Write a 250--300 word composition on the following topic: \emph{« ICTs are more useful for recreation than for learning. Do you agree? Discuss, giving examples from your community. »} Follow the steps: analyse the topic, brainstorm, build an outline, draft, check with the checklist, then write your final copy.
\end{exercice}

\begin{corrige}[Exercice 1]
Elements of a good answer. \textbf{Introduction:} hook (a question about young people and phones) + thesis (e.g. « In my opinion, ICTs serve both purposes, but recreation dominates in my community »). \textbf{Body paragraph 1:} recreational uses (music, films, games) with an example. \textbf{Body paragraph 2:} learning uses (online lessons, research, typing practice) with an example. \textbf{Body paragraph 3 (optional):} why recreation dominates (cheap bundles, entertainment culture). \textbf{Conclusion:} balanced opinion + recommendation with \emph{should} (e.g. « Students should use their phones for learning as much as for fun »). Expected language: linkers (\emph{firstly, however, on the other hand, in conclusion}), comparatives (\emph{more useful than}), modals (\emph{should, mustn't}), lesson vocabulary, 250--300 words, correct spelling and punctuation.
\end{corrige}

\begin{aretenir}[To remember]
A successful composition follows the funnel: hook $\rightarrow$ thesis $\rightarrow$ body paragraphs (one idea each, built with T-E-E-L) $\rightarrow$ conclusion with opinion and recommendation. Always analyse the topic first, never copy it as your introduction, use linkers and lesson vocabulary, and finish by checking content, organisation, language and mechanics.
\end{aretenir}

\coursec{Integration activities, evaluation and remediation}

In the previous sections, you learnt how to write a composition on using ICTs as recreational resources at home or in the community. You now have vocabulary, grammar tools and writing skills. This final section brings everything together: you will use \emph{all} your skills (speaking, listening, reading and writing) in one integrated project, take a Bac-type evaluation, and then improve your weak points through remediation workshops. This is the heart of the \cle{competency-based approach} (\emph{approche par les compétences}, APC): learning is only complete when you can use it in a real-life situation.

\courssub{Integrated situation: ``The digital life of my neighbourhood''}

\begin{definition}[Integrated situation]
An \cle{integrated situation} (\emph{situation d'intégration}) is a complex, realistic task that forces you to combine several skills and pieces of knowledge at the same time, exactly as in real life. It is not a simple exercise: it has a context, a role for you, a task, and a final product.
\end{definition}

\textbf{The situation.} Your neighbourhood (\emph{quartier}) wants to create a ``Digital Life Charter'' to help young people use phones and the internet wisely. Your class has been invited to contribute. In groups of four, you must:

\begin{enumerate}
\item \textbf{Speaking:} interview two classmates about their phone and internet subscriptions (Which package? How much per month? Why this one?) and express your own preferences.
\item \textbf{Listening:} listen to a short recorded text (read by the teacher) about a student who uses ICTs to improve his school results, and answer five comprehension questions.
\item \textbf{Reading:} read a short article about how young Cameroonians use ICTs for recreation (music, games, social media, football streaming) and pick out the main ideas.
\item \textbf{Writing:} write a paragraph (8--10 lines) summarising your findings and giving one piece of advice for balanced ICT use.
\end{enumerate}

\begin{methode}[How to succeed in an integrated task]
\begin{enumerate}
\item \textbf{Read the whole situation first.} Identify the context, your role, the task and the expected product.
\item \textbf{List the skills you need.} Here: vocabulary of subscriptions, linkers, comparatives, paragraph structure.
\item \textbf{Share the work in the group,} but make sure every member practises every skill.
\item \textbf{Draft, check, improve.} Never present your first version.
\item \textbf{Present clearly:} speak loudly, look at the audience, respect the time limit.
\end{enumerate}
\end{methode}

\courssub{Group task: poster or short presentation on balanced ICT use}

Each group designs an A3 poster (or a 3-minute oral presentation) advising students on \cle{balanced ICT use} (\emph{usage équilibré des TIC}). Your poster must contain:

\begin{itemize}
\item a catchy title (e.g. \emph{``Smart phone, smart student!''});
\item at least \textbf{three pieces of advice} using imperatives and modals (\emph{Limit your screen time. You should not chat during lessons. You must protect your password.});
\item at least \textbf{two linkers} (\emph{however, moreover, first of all, in conclusion});
\item at least \textbf{one comparative} (\emph{Studying online is more useful than playing games all day.});
\item simple drawings or symbols (no complicated artwork needed).
\end{itemize}

\begin{exemplebox}[Sample advice sentences for the poster]
\begin{itemize}
\item \emph{First of all, use the internet to learn, not only to play.}
\item \emph{Social media is fun; however, too much of it harms your results.}
\item \emph{A cheap package with less data is better than an expensive one you cannot renew.}
\item \emph{Moreover, always switch off your phone one hour before sleeping.}
\end{itemize}
\end{exemplebox}

\courssub{Summative evaluation: Bac-type paper}

The \cle{summative evaluation} (\emph{évaluation sommative}) checks what you have mastered at the end of the lesson. It follows the format of the Baccalauréat technique English paper.

\begin{exemplebox}[Indicative marking scheme (barème indicatif)]
\begin{center}
\begin{tabularx}{\linewidth}{|l|X|c|}
\hline
\rowcolor{popblueL} \textbf{Part} & \textbf{Content} & \textbf{Marks} \\
\hline
Reading comprehension & Text on ICTs in daily life + 5 questions & /10 \\
\hline
Grammar & 10 items: tenses, comparatives, linkers, question tags & /5 \\
\hline
Guided composition & 80--100 words on recreational use of ICTs & /5 \\
\hline
\textbf{Total} & & \textbf{/20} \\
\hline
\end{tabularx}
\end{center}
\end{exemplebox}

\begin{methode}[Managing your time during the exam]
\begin{enumerate}
\item \textbf{Read the whole paper first} (5 minutes): identify easy and difficult questions.
\item \textbf{Share your time} according to the marks: for a 2-hour paper marked out of 20, allow about 50 minutes for comprehension, 25 minutes for grammar and 35 minutes for the composition.
\item \textbf{Start with what you know best} to gain confidence and marks quickly.
\item \textbf{Keep the last 10 minutes to proofread} (\emph{relire}) your copy: verb endings, agreements, spelling, punctuation.
\item \textbf{Never leave a question blank:} an attempted answer can earn partial marks.
\end{enumerate}
\end{methode}

\begin{attention}[Frequent error: neglecting proofreading]
Many students lose 2 to 3 marks simply because they do not reread their work. Check especially:
\begin{itemize}
\item \textbf{subject--verb agreement:} \emph{He play} $\rightarrow$ \emph{He plays};
\item \textbf{tenses:} do not mix past and present without reason;
\item \textbf{spelling of unit vocabulary:} \emph{subscription, package, data, recreational};
\item \textbf{capital letters and full stops} at the beginning and end of sentences.
\end{itemize}
A composition with good ideas but careless mistakes is marked down. Proofreading is not a luxury: it is part of the exam.
\end{attention}

\begin{exoresolu}[Bac-type grammar item: comparatives and linkers]
\textbf{Question.} Complete the sentences with the correct word from the box: \emph{cheaper --- however --- the most --- moreover}.

1. The weekly package is \dots than the monthly one, but it finishes quickly.

2. Mobile data is expensive; \dots, students still find money for it.

3. \dots, a good connection helps you do research for your projects.

4. This is \dots popular package in our neighbourhood.
\tcblower
\textbf{Solution.}
1. \emph{cheaper} --- comparative of superiority: \emph{cheap} + \emph{-er} + \emph{than}.

2. \emph{however} --- the semicolon and the contrast (expensive / still buy) require a linker of opposition.

3. \emph{Moreover} --- the sentence adds a new argument, so we need a linker of addition.

4. \emph{the most} --- superlative of a long adjective: \emph{the most popular}.
\end{exoresolu}

\begin{exercice}[Practice: mini Bac-type test (20 minutes)]
\textbf{A. Grammar.} Choose the correct option.

1. I prefer the night package \dots it is cheaper. (because / but / so)

2. The \dots you chat at night, the more tired you are in class. (more / most / much)

3. She \dots her subscription yesterday. (renews / renewed / renew)

\textbf{B. Guided composition.} In 60--80 words, describe how you use the internet for recreation and give one advice to a friend. Use at least two linkers.
\end{exercice}

\begin{corrige}[Practice: mini Bac-type test]
\textbf{A.} 1. \emph{because} (cause). 2. \emph{more} (parallel comparative: \emph{the more\dots the more\dots}). 3. \emph{renewed} (past tense, marker \emph{yesterday}).

\textbf{B. Sample answer.} \emph{In my free time, I use the internet to listen to music and watch football highlights. Moreover, I chat with my friends on WhatsApp. However, I know that too much screen time is dangerous for my studies. My advice to you is simple: finish your homework first, and then enjoy your phone. In conclusion, the internet is a good servant but a bad master.}
\end{corrige}

\courssub{Remediation workshops: working on weak points}

After marking, the teacher identifies the class's weak points and organises short \cle{remediation workshops} (\emph{ateliers de remédiation}). Here are three typical workshops.

\textbf{Workshop 1: Linkers.} Sort linkers into three families: \emph{addition} (and, moreover, also), \emph{opposition} (but, however, although), \emph{conclusion} (so, therefore, in conclusion). Then rewrite five sentences, replacing \emph{and} or \emph{but} with a richer linker.

\textbf{Workshop 2: Comparatives.} Drill the three forms: \emph{-er than} (short adjectives), \emph{more \dots than} (long adjectives), \emph{as \dots as} (equality). Transform sentences: \emph{The monthly package costs 2\,000 F. The weekly package costs 500 F.} $\rightarrow$ \emph{The weekly package is cheaper than the monthly one.}

\textbf{Workshop 3: Paragraph writing.} Rebuild a mixed-up paragraph: identify the topic sentence, the supporting sentences and the concluding sentence, then rewrite it with linkers.

\begin{savaistu}[Remediation and the APC in Cameroon]
Targeted remediation after each evaluation is a key principle of the \cle{competency-based approach} (APC) promoted by MINESEC in Cameroon. The idea is simple: an evaluation is not the end of learning. It is a \emph{diagnosis} that shows what each student must still work on, so that nobody is left behind before the Probatoire or Baccalauréat technique.
\end{savaistu}

\courssub{Self-assessment grid: ``Can I\dots?''}

Tick (\checkmark) the column that corresponds to your real level. Be honest: this grid is for \emph{you}.

\begin{popfigure}
\begin{center}
\begin{tabularx}{\linewidth}{|>{\raggedright\arraybackslash}X|c|c|c|}
\hline
\rowcolor{popblueL} \textbf{Can-do statement} & \textbf{Yes} & \textbf{A little} & \textbf{Not yet} \\
\hline
I can express my preferences on phone/internet packages. & \fbox{} & \fbox{} & \fbox{} \\
\hline
I can compare two service packages. & \fbox{} & \fbox{} & \fbox{} \\
\hline
I can understand a text on ICTs and answer questions. & \fbox{} & \fbox{} & \fbox{} \\
\hline
I can use linkers (however, moreover) correctly. & \fbox{} & \fbox{} & \fbox{} \\
\hline
I can write a structured paragraph of 8--10 lines. & \fbox{} & \fbox{} & \fbox{} \\
\hline
I can proofread and correct my own mistakes. & \fbox{} & \fbox{} & \fbox{} \\
\hline
\end{tabularx}
\end{center}
\legende{Self-assessment grid (grille d'auto-évaluation): tick one box per line. Any ``Not yet'' answer shows you which remediation workshop to attend.}
\end{popfigure}

\courssub{Homework: interview and report}

\begin{exercice}[Homework: family interview]
Interview a family member (father, mother, elder brother or sister) about their phone or internet subscription. Ask: \emph{Which network do you use? Which package do you subscribe to? How much do you spend per month? Why do you prefer it? Would you like to change?}

Then \textbf{write a report of 8--10 sentences} presenting the answers, using reported speech (\emph{She said that\dots}), at least one comparative and two linkers.
\end{exercice}

\begin{corrige}[Homework: family interview (sample report)]
\emph{I interviewed my mother about her phone subscription. She uses MTN and she subscribes to a monthly package of 2\,000 F. She said that this package was cheaper than buying daily bundles. Moreover, it gives her enough data for WhatsApp and her small business. She prefers this network because the connection is better in our quarter. However, she complained that the data finishes too fast. She told me that she would like a bigger package, but it is too expensive. In my opinion, her package is more economical than mine. In conclusion, my mother uses her phone wisely, mainly for work and family.}
\end{corrige}

\begin{aretenir}[To remember]
\begin{itemize}
\item An integrated situation combines speaking, listening, reading and writing around one real-life task.
\item In the exam: read the whole paper, share your time according to the marks, and always proofread.
\item Evaluation is followed by remediation: identify your weak points (linkers, comparatives, paragraph writing) and work on them.
\item Use the self-assessment grid honestly to know where you stand before the Baccalauréat technique.
\end{itemize}
\end{aretenir}

\newpage

\coursec{Activité d'intégration}

\courssub{Objectifs de l'activité}

By the end of this integration activity, students will be able to:

\begin{itemize}
\item \textbf{Compare} three real telephone/internet service packages using \cle{comparatives} and \cle{superlatives} (cheaper than, the best value for money, more data than).
\item \textbf{Express preferences} with appropriate structures (I would rather, I prefer, in my opinion, the best option is).
\item \textbf{Read and interpret} numerical data in a table (prices in FCFA, data volumes in MB/GB, validity periods).
\item \textbf{Listen} to a short recording about using ICTs to enhance learning and extract key information.
\item \textbf{Write} a structured 150-word composition justifying a choice and giving advice on the responsible recreational use of ICTs.
\item \textbf{Speak} in front of the class to present and defend a group decision.
\end{itemize}

This activity mobilises all the competences of the lesson: vocabulary of \cle{service packages} (\emph{vocabulaire des forfaits}), grammar revision, listening, reading and writing, in one authentic situation.

\courssub{Situation}

\textbf{Choosing the best bundle and advising my community}

The Students' Cooperative of Government Technical High School Nkol-Eton, Yaoundé, has received a subsidy. The cooperative manages the school's computer club, where 45 students come every afternoon to do research, prepare their Probatoire and Baccalauréat projects, learn word processing, and relax after classes.

The cooperative has a \textbf{budget of exactly 10,000 FCFA per month} for an internet connection. The president of the cooperative, Amina, has collected three real subscription offers (\emph{offres d'abonnement}) from the providers available in the neighbourhood:

\begin{center}
\begin{tabularx}{\linewidth}{|l|X|X|X|}
\hline
\rowcolor{popblueL} \textbf{Criteria} & \textbf{Offer A: MTN Go Data} & \textbf{Offer B: Orange Internet Max} & \textbf{Offer C: Camtel Home Start} \\
\hline
Monthly price & 9,500 FCFA & 10,000 FCFA & 7,000 FCFA \\
\hline
Data volume & 15 GB & 20 GB & 30 GB \\
\hline
Validity & 30 days & 30 days & 30 days \\
\hline
Network coverage at school & Excellent (4G) & Good (4G) & Limited (fibre not yet installed; 3G only) \\
\hline
Extra services & Free WhatsApp & Free access to e-learning platform & Unlimited night browsing (10 p.m.--6 a.m.) \\
\hline
Connection speed & Fast & Fast & Slow at peak hours \\
\hline
\end{tabularx}
\end{center}

In addition, Amina interviewed a student of the club, Boris, who recorded a one-minute audio message explaining how he uses ICTs to learn: he watches mathematics tutorials on his phone, downloads past Baccalauréat papers, and uses a dictionary app to improve his English.

The cooperative must choose \textbf{one single offer} and justify the choice in writing. The members must also give \textbf{two pieces of advice} to young people of the Nkol-Eton neighbourhood on how to use ICTs for recreation (games, music, films, social media) in a responsible way: protecting their eyes, managing their time, avoiding fake news and respecting others online.

\textbf{Your task:} in groups of four, compare the three offers, listen to Boris's message, write a 150-word composition justifying your choice and giving your two pieces of advice, then present your decision orally. The class will vote for the best proposal.

\courssub{Consignes}

\begin{enumerate}
\item \textbf{Reading the data:} Read the table carefully. For each offer, calculate the price of one gigabyte (price $\div$ data volume). Which offer gives the cheapest gigabyte?
\item \textbf{Comparing:} Write four sentences comparing the offers, using a comparative or a superlative in each one (for example: \emph{Offer C is cheaper than Offer A, but its coverage is worse.})
\item \textbf{Expressing preferences:} In your group, discuss the offers using at least three preference structures: \emph{I would rather\ldots, I prefer\ldots because\ldots, In my opinion\ldots, The best option is\ldots} Note that price is not the only criterion: think of coverage and speed too.
\item \textbf{Listening:} Listen to Boris's recording twice. Answer: (a) What three ICT tools does Boris use to learn? (b) How does he improve his English? (c) Which offer would help him most, and why?
\item \textbf{Writing:} Write a composition of about 150 words with three paragraphs: (1) presentation of the situation and the chosen offer; (2) two justified reasons for the choice (price per GB, coverage, validity, extra services); (3) two pieces of advice for the responsible recreational use of ICTs in your neighbourhood.
\item \textbf{Speaking:} Choose a spokesperson. Present your decision to the class in two minutes maximum. The class votes for the clearest and best-argued proposal.
\end{enumerate}

\courssub{Ressources à mobiliser}

\begin{aretenir}[Key structures and vocabulary]
\begin{itemize}
\item \textbf{Comparatives:} cheaper than, faster than, more expensive than, better/worse than; \textbf{superlatives:} the cheapest, the most reliable, the best value for money.
\item \textbf{Preferences:} I would rather + verb; I prefer + noun/-ing; In my opinion; As far as I am concerned; The best option is\ldots
\item \textbf{Subscription vocabulary:} bundle (\emph{forfait}), data volume (\emph{volume de données}), validity (\emph{validité}), coverage (\emph{couverture}), to subscribe (\emph{souscrire}), top-up (\emph{recharge}), provider (\emph{fournisseur}).
\item \textbf{ICT vocabulary:} to browse (\emph{naviguer}), to download (\emph{télécharger}), tutorial (\emph{tutoriel}), e-learning platform, screen time (\emph{temps d'écran}), fake news (\emph{fausses informations}).
\item \textbf{Link words for writing:} first, moreover, however, finally, to conclude.
\item \textbf{Calculation:} unit price $= \dfrac{\text{total price}}{\text{quantity}}$.
\end{itemize}
\end{aretenir}

\courssub{Démarche et corrigé commenté}

\begin{exoresolu}[Model answer: choosing and justifying]
\textbf{Step 1 -- Reading the data (price per GB).}

We compute the unit price for each offer:
\begin{align*}
\text{Offer A:} \quad & \frac{9500}{15} \approx 633 \text{ FCFA/GB}\\
\text{Offer B:} \quad & \frac{10000}{20} = 500 \text{ FCFA/GB}\\
\text{Offer C:} \quad & \frac{7000}{30} \approx 233 \text{ FCFA/GB}
\end{align*}
Offer C gives the cheapest gigabyte, but its coverage at the school is limited and the speed is slow at peak hours -- precisely when the club meets.

\tcblower
\textbf{Step 2 -- Comparing (examples of sentences).}

\begin{itemize}
\item Offer C is cheaper than Offer A and Offer B, but its coverage is the worst of the three.
\item Offer B is more expensive than Offer C, but it gives more data than Offer A.
\item Offer A has the best network coverage at school.
\item Offer B is the most expensive offer, but it offers the best balance between data and reliability.
\end{itemize}

\textbf{Step 3 -- Expressing preferences (group discussion).}

\emph{``I would rather choose Offer B because the coverage is good and the e-learning platform is free.''} -- \emph{``I prefer Offer A: it is the most reliable, and free WhatsApp helps us to share documents.''} -- \emph{``In my opinion, Offer C is risky because the club works in the afternoon, at peak hours.''}

\textbf{Step 4 -- Listening (Boris's message).}

(a) Boris uses three tools: video tutorials in mathematics, downloaded past examination papers, and a dictionary app. (b) He improves his English with the dictionary app. (c) Offer B would help him most, because the free e-learning platform gives him direct access to lessons and past papers without consuming his data volume.

\textbf{Step 5 -- Model composition (about 150 words).}

\emph{Our computer club has a budget of 10,000 FCFA per month for an internet connection. After comparing the three offers, our group has chosen Offer B, Orange Internet Max.}

\emph{First, this offer gives 20 GB for 10,000 FCFA, that is 500 FCFA per gigabyte, which is cheaper than Offer A. Moreover, the 4G coverage at our school is good and the speed is fast, unlike Offer C, which is very slow at peak hours when the club is open. Finally, the free access to the e-learning platform is a real advantage for students like Boris, who download past papers and watch tutorials.}

\emph{To conclude, we advise the young people of our neighbourhood to limit their screen time to two hours of games and films per day, and to check information before sharing it, in order to avoid fake news. ICTs are excellent tools if we use them responsibly.}

(Word count: approximately 150 words.)

\textbf{Step 6 -- Oral presentation and vote.}

The spokesperson presents the decision in two minutes, using the comparative sentences and the preference structures. The class votes; the teacher checks that every group used at least three target structures and one correct calculation.
\end{exoresolu}

\courssub{Bilan de l'activité}

This activity has shown that the competences of the lesson work together in real life. The students have discovered that choosing a service package is not only a question of price: they had to combine a mathematical calculation (price per GB), a critical reading of technical criteria (coverage, speed, validity), and a reasoned expression of preference in English. The listening task reminded them that ICTs are powerful tools for learning -- tutorials, past papers, dictionary apps -- and the writing task led them to formulate responsible rules for recreational use: limiting screen time, checking information, respecting others online. Finally, the oral presentation and the class vote developed fluency, argumentation and citizenship: defending a collective decision in front of others is exactly the kind of real-life situation (\emph{situation authentique}) that the Baccalauréat technique values.

\newpage

\coursec{Méthodes et exercices résolus}

\coursec{Media and Communication}

\courssub{Subscribing to Service Packages: Vocabulary and Speaking}

\begin{methode}[Expressing Preferences for Service Packages]
\textbf{Step 1: Identify your needs.} List what you use most: calls, SMS, mobile data, international minutes, streaming, social media access.
\textbf{Step 2: Compare packages.} Look at price, validity period, data volume, call rates, bonus offers (free WhatsApp, Facebook, YouTube).
\textbf{Step 3: Use preference structures.}
\begin{itemize}
    \item \cle{I'd rather} + base verb (+ than + base verb) — \emph{Je préférerais}
    \item \cle{I prefer} + -ing / noun (+ to + -ing / noun) — \emph{Je préfère}
    \item \cle{I'd like to subscribe to} + package name — \emph{Je voudrais m'abonner à}
    \item \cle{What does the package include?} — \emph{Que comprend le forfait ?}
    \item \cle{Is there a...?} / \cle{Does it come with...?} — \emph{Est-ce qu'il y a... ? / Est-ce que cela inclut... ?}
\end{itemize}
\textbf{Step 4: Ask for clarification.} \cle{Could you explain the difference between...?} — \emph{Pourriez-vous expliquer la différence entre... ?}
\textbf{Step 5: Confirm and subscribe.} \cle{I'll take the [name] package, please.} / \cle{Please activate the [name] plan for me.}
\end{methode}

\begin{exoresolu}[Choosing a Mobile Data Package]
\textbf{Task:} You are at a telecom agency (MTN / Orange / Camtel). You need a monthly package for heavy internet use (online classes, video calls, downloading course materials). You have a budget of 5\,000 FCFA. The agent presents three packages. Choose the best one and justify your choice orally.

\textbf{Package A:} 2\,500 FCFA — 2 GB data + 30 min calls + unlimited SMS (valid 30 days).
\textbf{Package B:} 5\,000 FCFA — 10 GB data + 100 min calls + unlimited WhatsApp/Facebook (valid 30 days).
\textbf{Package C:} 3\,500 FCFA — 5 GB data + 50 min calls + unlimited social media (valid 15 days).

\textbf{Instructions:} Write a short dialogue (6–8 exchanges) between you and the agent. Use preference structures and ask at least two clarification questions.

\tcblower
\textbf{Solution détaillée}

\textbf{Reasoning:}
\begin{itemize}
    \item Package A: 2 GB is insufficient for video calls and downloading course materials.
    \item Package C: 5 GB for 15 days = 10 GB/month equivalent, but validity is only 15 days; you would need to renew twice (7\,000 FCFA/month), over budget.
    \item Package B: 10 GB for 30 days at 5\,000 FCFA fits the budget exactly, includes unlimited WhatsApp/Facebook (useful for class groups), and 100 min calls for emergencies.
\end{itemize}
\textbf{Choice:} Package B.

\textbf{Dialogue:}
\begin{itemize}
    \item \textbf{Agent:} Good morning! How can I help you today?
    \item \textbf{You:} Good morning. \cle{I'd like to subscribe to} a monthly internet package for my studies.
    \item \textbf{Agent:} We have three options. Package A at 2\,500 FCFA with 2 GB, Package B at 5\,000 FCFA with 10 GB, and Package C at 3\,500 FCFA with 5 GB for 15 days.
    \item \textbf{You:} \cle{What does Package B include exactly?}
    \item \textbf{Agent:} 10 GB data, 100 minutes of calls, and unlimited WhatsApp and Facebook for 30 days.
    \item \textbf{You:} \cle{Does it come with} unlimited YouTube for educational videos?
    \item \textbf{Agent:} No, YouTube uses the main data allowance. But WhatsApp and Facebook are free.
    \item \textbf{You:} \cle{I prefer Package B} because I need lots of data for online classes. \cle{I'd rather have} 10 GB for a full month than 5 GB for two weeks.
    \item \textbf{Agent:} Perfect. \cle{I'll activate Package B} for you right away.
\end{itemize}
\textbf{Conclusion:} Package B is the only one that meets the data needs, budget, and validity requirements.
\end{exoresolu}

\courssub{Grammar General Revision}

\begin{methode}[Key Grammar Points for Revision (Bac Technique)]
\textbf{1. Tenses for Narration and Description}
\begin{itemize}
    \item \cle{Present Simple:} habits, facts, timetables — \emph{I check my emails every morning.}
    \item \cle{Present Continuous:} actions now, temporary situations — \emph{She is downloading the file right now.}
    \item \cle{Past Simple:} finished past actions — \emph{We bought a new router last week.}
    \item \cle{Present Perfect:} past action with present result / experience — \emph{I have lost my password. / Have you ever used Zoom?}
\end{itemize}
\textbf{2. Modals for Functions}
\begin{itemize}
    \item \cle{Can / Could:} ability, permission, requests — \emph{Can you help me?}
    \item \cle{Should / Ought to:} advice — \emph{You should update your antivirus.}
    \item \cle{Must / Have to:} obligation — \emph{You must protect your data.}
    \item \cle{May / Might:} possibility — \emph{The network might be down.}
\end{itemize}
\textbf{3. Passive Voice (Present/Past)}
\cle{Subject + be + past participle} — \emph{The software is updated automatically. / The message was sent yesterday.}
\textbf{4. Reported Speech (Backshifting)}
\cle{Direct: "I need a new SIM card." → Reported: He said he needed a new SIM card.}
\textbf{5. Relative Clauses}
\cle{who (people), which (things), where (places), whose (possession)} — \emph{The technician who fixed my phone is skilled.}
\textbf{6. Conditionals (Types 0, 1, 2)}
\begin{itemize}
    \item Type 0: \cle{If + present, present} (facts) — \emph{If you press the button, the screen lights up.}
    \item Type 1: \cle{If + present, will + base} (real future) — \emph{If the Wi-Fi fails, I will use mobile data.}
    \item Type 2: \cle{If + past, would + base} (hypothetical) — \emph{If I had 5G, I would download faster.}
\end{itemize}
\textbf{7. Connectors for Writing}
\cle{However, Moreover, Furthermore, Therefore, Consequently, In addition, On the other hand.}
\end{methode}

\begin{exoresolu}[Grammar Revision: Mixed Exercises]
\textbf{Task:} Complete the following exercises. They reflect typical Bac Technique grammar items.

\textbf{A. Put the verbs in brackets into the correct tense.}
1. Look! That man \_\_\_\_\_\_ (try) to connect to the Wi-Fi.
2. My brother \_\_\_\_\_\_ (buy) a new smartphone last month.
3. We \_\_\_\_\_\_ (not / finish) the online course yet.
4. While I \_\_\_\_\_\_ (download) the file, the power \_\_\_\_\_\_ (go) out.
5. By the time you arrive, I \_\_\_\_\_\_ (set up) the router.

\textbf{B. Rewrite using the passive voice.}
6. They install fibre optic cables in our neighbourhood.
7. Someone has stolen my laptop.
8. The company will launch a new 5G plan next year.

\textbf{C. Complete with the correct modal (can, could, should, must, might).}
9. You \_\_\_\_\_\_ not share your password with anyone. (strong obligation)
10. \_\_\_\_\_\_ you please send me the link? (polite request)
11. The connection \_\_\_\_\_\_ be slow because of the rain. (possibility)
12. Students \_\_\_\_\_\_ use the school Wi-Fi for research only. (rule)

\textbf{D. Rewrite in reported speech.}
13. "I am learning coding," she said.
14. "We will repair the network tomorrow," the technician said.

\textbf{E. Join the sentences using a relative pronoun.}
15. The app is very useful. I downloaded it yesterday.
16. The engineer works for Camtel. His daughter is in my class.

\tcblower
\textbf{Solution détaillée}

\textbf{A. Tenses}
1. \textbf{is trying} — Present Continuous (action happening now, indicated by "Look!").
2. \textbf{bought} — Past Simple (finished action with "last month").
3. \textbf{have not finished} / \textbf{haven't finished} — Present Perfect (action not completed yet, "yet" marker).
4. \textbf{was downloading} / \textbf{went} — Past Continuous (long action interrupted) + Past Simple (short interrupting action).
5. \textbf{will have set up} — Future Perfect (action completed before a future time "by the time you arrive").

\textbf{B. Passive Voice}
6. \textbf{Fibre optic cables are installed in our neighbourhood.} (Present Simple passive: are + past participle)
7. \textbf{My laptop has been stolen.} (Present Perfect passive: has been + past participle)
8. \textbf{A new 5G plan will be launched next year.} (Future Simple passive: will be + past participle)

\textbf{C. Modals}
9. \textbf{must} — Strong obligation/prohibition (internal rule, safety).
10. \textbf{Could} — Polite request (more polite than "Can").
11. \textbf{might} — Possibility (less certain than "may").
12. \textbf{must} — Rule/regulation (external obligation).

\textbf{D. Reported Speech (Backshifting)}
13. \textbf{She said (that) she was learning coding.} (Present Continuous → Past Continuous)
14. \textbf{The technician said (that) they would repair the network the next day / the following day.} (Will → Would; Tomorrow → The next day)

\textbf{E. Relative Clauses}
15. \textbf{The app which/that I downloaded yesterday is very useful.} (Object relative clause: "I downloaded the app" → "which/that I downloaded")
16. \textbf{The engineer whose daughter is in my class works for Camtel.} (Possessive: "His daughter" → "whose daughter")

\textbf{Conclusion:} Mastery of these grammar points is essential for the Bac Technique English exam, especially for writing and transformation exercises.
\end{exoresolu}

\courssub{Listening: ICTs to Enhance Work and Learning}

\begin{methode}[Effective Listening Strategy for ICT Topics]
\textbf{Step 1: Predict content from keywords.} Before listening, read the questions. Identify key terms: \cle{remote work}, \cle{e-learning}, \cle{video conferencing}, \cle{cloud storage}, \cle{collaboration tools}, \cle{digital skills}, \cle{productivity}.
\textbf{Step 2: Listen for gist (first listening).} Grasp the main idea: Who? What? Context? (e.g., a teacher explaining Google Classroom, an employee describing working from home).
\textbf{Step 3: Listen for detail (second listening).} Note specific information: numbers, names of apps, advantages/disadvantages, recommendations.
\textbf{Step 4: Recognise signpost language.}
\begin{itemize}
    \item \cle{First / Firstly / To start with} — \emph{Premièrement}
    \item \cle{Another advantage is...} / \cle{Moreover...} — \emph{Un autre avantage...}
    \item \cle{However / On the other hand...} — \emph{Cependant...}
    \item \cle{For example / For instance...} — \emph{Par exemple...}
    \item \cle{In conclusion / To sum up...} — \emph{En conclusion...}
\end{itemize}
\textbf{Step 5: Infer meaning of unknown words.} Use context: \cle{bandwidth} (internet speed), \cle{lag} (delay), \cle{glitch} (small technical problem), \cle{seamless} (smooth, without interruption).
\textbf{Step 6: Answer in full sentences.} Use your own words where possible.
\end{methode}

\begin{exoresolu}[Listening Comprehension: Remote Work Testimony]
\textbf{Task:} Listen to the following text (read by teacher or audio) and answer the questions.

\textbf{Text:}
"Good afternoon. My name is Marie Ngo, and I work as a graphic designer for a marketing agency in Douala. Since the pandemic, I have been working remotely three days a week. At first, it was challenging. My internet connection was unstable, and I had frequent power cuts. I missed the face-to-face interaction with my colleagues. However, my company provided us with laptops and a monthly internet allowance of 10\,000 FCFA. We use Zoom for daily meetings, Slack for instant messaging, and Google Drive for file sharing. These tools make collaboration seamless. I can access my projects from anywhere. Last month, I even worked from my village in the West Region for two weeks. The only downside is the occasional lag during video calls when the network is congested. Overall, remote work has improved my work-life balance. I save time and money on transport. I would recommend it to anyone who has a reliable internet connection and good self-discipline."

\textbf{Questions:}
1. What is Marie's job and where does she work?
2. How many days a week does she work remotely?
3. What two problems did she face at the beginning?
4. Name three digital tools her company uses.
5. What did the company provide to support remote work?
6. Where did she work from last month?
7. What is the main disadvantage she mentions?
8. What two benefits does she highlight in conclusion?
9. What two conditions does she give for successful remote work?

\tcblower
\textbf{Solution détaillée}

\textbf{Answers with justifications from text:}
1. \textbf{She is a graphic designer working for a marketing agency in Douala.} (Directly stated: "My name is Marie Ngo, and I work as a graphic designer for a marketing agency in Douala.")
2. \textbf{Three days a week.} (Directly stated: "I have been working remotely three days a week.")
3. \textbf{Unstable internet connection and frequent power cuts.} (Directly stated: "My internet connection was unstable, and I had frequent power cuts.")
4. \textbf{Zoom (video meetings), Slack (instant messaging), Google Drive (file sharing).} (Directly stated: "We use Zoom for daily meetings, Slack for instant messaging, and Google Drive for file sharing.")
5. \textbf{Laptops and a monthly internet allowance of 10\,000 FCFA.} (Directly stated: "my company provided us with laptops and a monthly internet allowance of 10\,000 FCFA.")
6. \textbf{Her village in the West Region.} (Directly stated: "Last month, I even worked from my village in the West Region for two weeks.")
7. \textbf{Occasional lag during video calls when the network is congested.} (Directly stated: "The only downside is the occasional lag during video calls when the network is congested.")
8. \textbf{Improved work-life balance; saves time and money on transport.} (Directly stated: "Overall, remote work has improved my work-life balance. I save time and money on transport.")
9. \textbf{Reliable internet connection and good self-discipline.} (Directly stated: "I would recommend it to anyone who has a reliable internet connection and good self-discipline.")

\textbf{Conclusion:} This listening task tests the ability to extract specific factual information and understand advantages/disadvantages of ICT in the workplace. Key vocabulary: \cle{remote work} (travail à distance), \cle{allowance} (allocation/indemnité), \cle{seamless} (fluide/sans accroc), \cle{lag} (latence/ralentissement), \cle{congested} (saturé), \cle{self-discipline} (autodiscipline).
\end{exoresolu}

\courssub{Reading: ICTs as Recreational Resources}

\begin{methode}[Reading Strategy for ICT Leisure Texts]
\textbf{Step 1: Skim for main idea.} Read title, first and last paragraph. Identify text type: blog post, article, forum discussion, review. Identify topic: gaming, streaming, social media, content creation, online communities.
\textbf{Step 2: Scan for specific details.} Look for: names of platforms (Netflix, YouTube, TikTok, Twitch, Steam, Discord), devices (smartphone, console, PC, smart TV), activities (streaming movies, multiplayer gaming, vlogging, podcasting), time spent, costs.
\textbf{Step 3: Analyse structure.} Identify paragraphs: Introduction (hook + thesis), Body paragraphs (one idea each: social connection, skill development, relaxation, risks), Conclusion (summary + opinion/recommendation).
\textbf{Step 4: Vocabulary in context.} Guess meaning from co-text:
\begin{itemize}
    \item \cle{binge-watch} = watch many episodes in one sitting — \emph{regarder en rafale}
    \item \cle{streamer} = person who broadcasts live — \emph{streamer}
    \item \cle{multiplayer} = game with multiple players — \emph{multijoueur}
    \item \cle{offline} = not connected — \emph{hors ligne}
    \item \cle{screen time} = time spent on devices — \emph{temps d'écran}
    \item \cle{digital wellbeing} = healthy tech habits — \emph{bien-être numérique}
\end{itemize}
\textbf{Step 5: Answer comprehension questions.} Distinguish: \cle{Right there} (explicit), \cle{Think and search} (inference), \cle{Author and me} (opinion).
\textbf{Step 6: Critical thinking.} Evaluate: Is the author biased? Are examples relevant to Cameroon context?
\end{methode}

\begin{exoresolu}[Reading Comprehension: ICT for Leisure in Cameroon]
\textbf{Task:} Read the text and answer the questions.

\textbf{Text:}
\textbf{How Young Cameroonians Use ICT for Fun}
"In Yaoundé, Douala, and even smaller towns like Bafoussam or Bertoua, cybercafés are no longer the only places to go online. With affordable smartphones and data bundles, young people access entertainment anytime. Seventeen-year-old Kevin, a student in Terminale, spends his evenings on TikTok and YouTube. 'I follow tech reviewers and football highlights,' he says. 'I also learn video editing from free tutorials.' His friend Amina prefers streaming series on Netflix via her father's account. 'We binge-watch Korean dramas on weekends,' she laughs.
Online gaming is booming. PUBG Mobile, Free Fire, and Call of Duty Mobile are popular because they run on mid-range phones. Cybercafés have transformed into gaming hubs where youths pay 200 FCFA per hour to play on PCs with better graphics. 'It's social,' explains Franck, 19. 'I play with my cousins in France and Canada. We talk on Discord while gaming.'
Content creation is another trend. Many students run Instagram pages or YouTube channels about fashion, comedy, or education. Some earn money through ads and sponsorships. However, parents and teachers worry about excessive screen time. 'My son stays up until 2 a.m.,' complains a mother in Douala. 'His grades dropped.' Experts advise setting limits: no phones during meals, one hour before bed, and using 'Digital Wellbeing' tools to track usage. Balance is key. ICT offers amazing recreational and learning opportunities, but real life — sports, family, sleep — must come first."

\textbf{Questions:}
1. How has internet access changed for young Cameroonians compared to the past?
2. Give two examples of what Kevin does online.
3. Why are mobile games like PUBG Mobile popular in Cameroon?
4. How have cybercafés adapted?
5. What platform does Franck use to talk while gaming?
6. Name two types of content young creators produce.
7. What concern do parents and teachers share?
8. What three practical tips do experts give for digital wellbeing?
9. Find words in the text meaning: (a) watch many episodes at once, (b) centre/hub, (c) sponsorship/parrainage, (d) excessive.
10. According to the text, what should "come first" before ICT recreation?

\tcblower
\textbf{Solution détaillée}

1. \textbf{Internet access is now mobile and affordable (smartphones + data bundles), no longer limited to cybercafés.} (Paragraph 1: "With affordable smartphones and data bundles, young people access entertainment anytime.")
2. \textbf{He follows tech reviewers and football highlights on TikTok/YouTube; he learns video editing from free tutorials.} (Paragraph 2: "'I follow tech reviewers and football highlights,' he says. 'I also learn video editing from free tutorials.'")
3. \textbf{They run on mid-range (affordable) phones, not requiring expensive gaming PCs.} (Paragraph 3: "...popular because they run on mid-range phones.")
4. \textbf{They have transformed into gaming hubs with PCs for hourly rental (200 FCFA/hour).} (Paragraph 3: "Cybercafés have transformed into gaming hubs where youths pay 200 FCFA per hour to play on PCs with better graphics.")
5. \textbf{Discord.} (Paragraph 3: "We talk on Discord while gaming.")
6. \textbf{Fashion, comedy, education (any two).} (Paragraph 4: "Many students run Instagram pages or YouTube channels about fashion, comedy, or education.")
7. \textbf{Excessive screen time affecting sleep and school grades.} (Paragraph 4: "'My son stays up until 2 a.m.,' complains a mother... 'His grades dropped.'")
8. \textbf{No phones during meals; no phones one hour before bed; use 'Digital Wellbeing' tools to track usage.} (Paragraph 4: "Experts advise setting limits: no phones during meals, one hour before bed, and using 'Digital Wellbeing' tools to track usage.")
9. \textbf{Vocabulary in context:}
   (a) \textbf{binge-watch} (Paragraph 2: "'We binge-watch Korean dramas...'")
   (b) \textbf{hub} / \textbf{gaming hub} (Paragraph 3: "Cybercafés have transformed into gaming hubs...")
   (c) \textbf{sponsorship} (Paragraph 4: "...earn money through ads and sponsorships.")
   (d) \textbf{excessive} (Paragraph 4: "...worry about excessive screen time.")
10. \textbf{Real life — sports, family, sleep — must come first.} (Last sentence: "but real life — sports, family, sleep — must come first.")

\textbf{Conclusion:} This text illustrates recreational ICT use in the Cameroonian context. It tests scanning for details, vocabulary inference, and understanding the author's balanced view (opportunities vs. risks). Key terms: \cle{data bundles} (forfaits data), \cle{mid-range phones} (smartphones milieu de gamme), \cle{gaming hub} (lieu de jeux), \cle{content creation} (création de contenu), \cle{digital wellbeing} (bien-être numérique).
\end{exoresolu}

\courssub{Writing: Composition on Recreational Use of ICTs}

\begin{methode}[Writing a Balanced Composition (For/Against or Opinion Essay)]
\textbf{Step 1: Analyse the prompt.} Identify: Topic (ICT for recreation), Task type (advantages/disadvantages, opinion, proposal), Format (essay, article, letter), Length (150–200 words for Bac Technique).
\textbf{Step 2: Brainstorm and organise.} Use a T-chart:
\begin{center}
\begin{tabular}{|l|l|}
\hline
\textbf{Advantages (Pros)} & \textbf{Disadvantages (Cons)} \\ \hline
Relaxation / stress relief & Addiction / excessive screen time \\
Social connection (friends abroad) & Sleep deprivation \\
Skill learning (tutorials, languages) & Cyberbullying / inappropriate content \\
Creativity (editing, coding, art) & Sedentary lifestyle / health issues \\
Income potential (content creation) & Cost of data / devices \\ \hline
\end{tabular}
\end{center}
\textbf{Step 3: Structure the essay.}
\begin{itemize}
    \item \textbf{Introduction (40–50 words):} Hook (fact/quote/question) + Topic sentence + Thesis statement (your position).
    \item \textbf{Body Paragraph 1 (60–80 words):} Advantages with examples (topic sentence + 2–3 supporting sentences + connectors: \cle{Moreover}, \cle{For instance}).
    \item \textbf{Body Paragraph 2 (60–80 words):} Disadvantages with examples (topic sentence + 2–3 supporting sentences + connectors: \cle{However}, \cle{On the other hand}, \cle{For example}).
    \item \textbf{Conclusion (30–40 words):} Restate thesis + summary + recommendation / call to action (\cle{In conclusion}, \cle{To sum up}, \cle{I recommend}).
\end{itemize}
\textbf{Step 4: Use cohesive devices.} \cle{Furthermore}, \cle{In addition}, \cle{Consequently}, \cle{As a result}, \cle{Nevertheless}.
\textbf{Step 5: Proofread (COPS).} \cle{C}apitalisation, \cle{O}rganisation, \cle{P}unctuation, \cle{S}pelling. Check word count.
\end{methode}

\begin{exoresolu}[Writing Task: Advantages and Disadvantages of ICT for Leisure]
\textbf{Task:} "ICTs offer great recreational opportunities for young people, but they also present risks. Write an essay of 150–180 words discussing the advantages and disadvantages of using ICTs for leisure. Give examples from your community."

\textbf{Instructions:} Follow the method above. Write a complete essay.

\tcblower
\textbf{Model Essay (approx. 170 words)}

\textbf{ICT for Leisure: A Double-Edged Sword}

In Cameroon today, smartphones and affordable data have put a world of entertainment in the pockets of young people. From streaming movies to online gaming and content creation, ICTs have transformed how we relax. However, this digital playground has a dark side.

On the one hand, ICTs provide healthy recreation and learning. For instance, many students in Yaoundé use YouTube tutorials to learn video editing, coding, or foreign languages for free. Moreover, multiplayer games like PUBG Mobile allow youths to socialise with friends and relatives abroad, strengthening bonds across distances. Some even earn income: a classmate of mine monetises his comedy sketches on TikTok, helping pay his school fees.

On the other hand, excessive screen time causes serious problems. Many teenagers stay up late binge-watching series, leading to fatigue and poor grades. Cybercafés turned gaming hubs keep youths indoors for hours, reducing physical activity. Worse, exposure to cyberbullying and inappropriate content threatens their mental wellbeing.

In conclusion, ICTs are powerful tools for recreation and self-development, but they require discipline. I recommend that parents and schools promote digital wellbeing: set screen-time limits, encourage offline hobbies, and teach critical thinking online. Balance is the key to enjoying technology without becoming its slave.

\textbf{Analysis of the essay:}
\begin{itemize}
    \item \textbf{Structure:} Clear 4-paragraph structure (Intro, Pros, Cons, Conclusion).
    \item \textbf{Connectors:} \cle{On the one hand}, \cle{For instance}, \cle{Moreover}, \cle{On the other hand}, \cle{Worse}, \cle{In conclusion}.
    \item \textbf{Vocabulary:} \cle{double-edged sword}, \cle{affordable data}, \cle{monetises}, \cle{binge-watching}, \cle{cyberbullying}, \cle{digital wellbeing}, \cle{critical thinking}.
    \item \textbf{Contextualisation:} Examples from Yaoundé, Cameroonian reality (school fees, cybercafés).
    \item \textbf{Word count:} ~170 words (within 150–180 range).
\end{itemize}
\textbf{Conclusion:} A successful Bac essay addresses the prompt directly, uses a clear structure, varied connectors, relevant local examples, and appropriate register.
\end{exoresolu}

\courssub{Integration Activities, Evaluation and Remediation}

\begin{methode}[Integrated Task: Service Package Presentation + Report]
\textbf{Context:} End-of-unit project combining Speaking, Vocabulary, Grammar, Reading, Writing.
\textbf{Task:} In groups of 3–4, design a \cle{"Student Digital Pack"} for a telecom operator (MTN, Orange, Camtel). The pack must include: mobile data, call minutes, SMS, and \cle{value-added services} (educational apps zero-rated, free access to Moodle/Google Classroom, discounted devices).
\textbf{Steps:}
\textbf{1. Research (Reading/Listening):} Visit operator websites / interview a sales agent. Note real packages and prices.
\textbf{2. Design (Speaking/Vocabulary):} Create your ideal pack. Name it. Set price (max 5\,000 FCFA/month). Justify each component using \cle{ICT for learning} arguments.
\textbf{3. Presentation (Speaking/Grammar):} 3-minute pitch to the class. Use: \cle{Passive} ("The pack is designed for..."), \cle{Modals} ("Students can access..."), \cle{Conditionals} ("If a student subscribes, they will get..."), \cle{Comparatives} ("Our pack offers more data than...").
\textbf{4. Report Writing (Writing):} Write a 150-word formal proposal to the operator's marketing director. Structure: Subject line, Introduction, Package details (bullet points), Benefits for students/operator, Call to action, Formal closing.
\textbf{5. Peer Evaluation:} Use a rubric: Content (5), Vocabulary (5), Grammar (5), Fluency/Pronunciation (5), Visuals/Handout (5) = /25.
\textbf{6. Remediation:} Teacher identifies common errors (e.g., \cle{*I am agree} → \cle{I agree}, \cle{*informations} → \cle{information}, \cle{*advices} → \cle{advice}, tense mixing). Mini-lesson + targeted exercises.
\end{methode}

\begin{exoresolu}[Integrated Evaluation: Formal Proposal Letter]
\textbf{Task:} Write the formal proposal letter (Step 4 above) for your group's "Student Digital Pack". Use the following details:
\textbf{Pack Name:} \cle{Campus Connect 5000}
\textbf{Price:} 5\,000 FCFA / month
\textbf{Contents:} 15 GB data (4G/5G), 200 min calls (all networks), Unlimited SMS, Zero-rated access to: Moodle, Google Classroom, Khan Academy, Wikipedia, Coursera, Cameroon Ministry of Education portal. Bonus: 50\% discount on entry-level smartphone with 12-month subscription.
\textbf{Target:} Marketing Director, MTN Cameroon.

\tcblower
\textbf{Model Formal Proposal}

\begin{center}
\begin{tabular}{ll}
\textbf{From:} & Students' Representative Council, Government Technical High School, Yaoundé \\
\textbf{To:} & Marketing Director, MTN Cameroon \\
\textbf{Date:} & 15 October 2025 \\
\textbf{Subject:} & Proposal for "Campus Connect 5000" — Student Digital Package \\
\end{tabular}
\end{center}

\vspace{0.5cm}

Dear Sir/Madam,

We write to propose a dedicated mobile package for secondary and higher education students in Cameroon: the \textbf{Campus Connect 5000}.

At 5\,000 FCFA per month, this pack offers 15 GB of 4G/5G data, 200 minutes of calls to all networks, and unlimited SMS. Its unique value lies in \textbf{zero-rated access to educational platforms}: Moodle, Google Classroom, Khan Academy, Wikipedia, Coursera, and the Ministry of Education's digital portal. This ensures uninterrupted learning even when students exhaust their main data bundle.

Furthermore, a 50\% discount on an entry-level smartphone for any 12-month subscription would lower the device barrier for many youths.

This package benefits MTN by securing long-term young customers, enhancing brand loyalty, and supporting national digital education goals. We request a meeting to discuss this partnership.

We look forward to your favourable response.

Yours faithfully,

\textit{Ngo Marie} \\
President, Students' Representative Council

\textbf{Evaluation Checklist (Teacher's Rubric):}
\begin{itemize}
    \item \textbf{Format:} Correct formal letter layout (addresses, date, subject, salutation, closing) — /5
    \item \textbf{Content:} All package details included, clear benefits for both parties — /5
    \item \textbf{Vocabulary:} Technical terms used correctly (zero-rated, entry-level, subscription, bundle, partnership) — /5
    \item \textbf{Grammar:} Passive (\emph{is offered}, \emph{would be lowered}), Modals (\emph{ensures}, \emph{would secure}), Conditionals (\emph{If students subscribe, they will benefit}) — /5
    \item \textbf{Cohesion:} Connectors (\emph{Furthermore}, \emph{This package benefits... by...}) — /5
\end{itemize}
\textbf{Total: /25}
\end{exoresolu}

\begin{attention}[Les 5 erreurs qui coûtent des points au Bac]
\begin{enumerate}
    \item \textbf{Confusion \cle{advice} (invariable) vs \cle{advise} (verbe).} \emph{Faux:} "He gave me a good advice." \rightarrow \textbf{Correct:} "He gave me \cle{good advice} / a good \cle{piece of advice}." \emph{Faux:} "I advice you to..." \rightarrow \textbf{Correct:} "I \cle{advise} you to..."
    \item \textbf{Oubli du \cle{s} à la 3\textsuperscript{e} personne du singulier au Present Simple.} \emph{Faux:} "She \cle{use} her phone for studying." \rightarrow \textbf{Correct:} "She \cle{uses}..." (Règle: He/She/It + verbe + s).
    \item \textbf{Mauvaise formation du Passive (oublier \cle{be} ou mauvais participe passé).} \emph{Faux:} "The package \cle{launched} last week." \rightarrow \textbf{Correct:} "The package \cle{was launched} last week." \emph{Faux:} "The data \cle{is use}." \rightarrow \textbf{Correct:} "The data \cle{is used}."
    \item \textbf{Traduction littérale de \cle{"Je suis d'accord"}.} \emph{Faux:} "\cle{I am agree}." \rightarrow \textbf{Correct:} "\cle{I agree}." (Agree est un verbe, pas un adjectif).
    \item \textbf{Connecteurs mal placés ou ponctuation incorrecte.} \emph{Faux:} "However, it is expensive. But it is good." \rightarrow \textbf{Correct:} "\cle{However}, it is expensive; \cle{nevertheless}, it is good." Ou: "It is expensive, \cle{but} it is good." (Ne pas commencer une phrase par "But" dans l'écrit formel; utiliser \cle{However} / \cle{Nevertheless} + virgule).
\end{enumerate}
\end{attention}

\newpage

\coursec{Exercices}

\courssub{Je vérifie mes connaissances}

\begin{aretenir}[Vocabulaire clé : Service Packages \textit{(Forfaits de service)}]
\textbf{Termes essentiels} pour s'abonner aux services téléphone/internet :
\begin{itemize}
\item \cle{Data plan} / \cle{Data package} : forfait de données internet (Mo/Go) \textit{(Forfait data)}
\item \cle{Subscription fee} / \cle{Monthly fee} : frais d'abonnement mensuels \textit{(Frais d'abonnement)}
\item \cle{Coverage area} / \cle{Network coverage} : zone de couverture du réseau \textit{(Zone de couverture)}
\item \cle{Roaming charges} : frais d'itinérance (hors réseau/détranger) \textit{(Frais d'itinérance)}
\item \cle{Bundle} / \cle{Package deal} : offre groupée (internet + appels + SMS) \textit{(Forfait groupé)}
\item \cle{Unlimited calls} / \cle{Unlimited minutes} : appels illimités \textit{(Appels illimités)}
\item \cle{Prepaid} / \cle{Postpaid} : prépayé / postpayé \textit{(Prépayé / Postpayé)}
\item \cle{Mobile Money} / \cle{MoMo} : paiement mobile (MTN MoMo, Orange Money) \textit{(Argent mobile)}
\item \cle{ID card} / \cle{Proof of address} : carte d'identité / justificatif de domicile \textit{(CNI / Justificatif de domicile)}
\end{itemize}
\end{aretenir}

\begin{exercice}[Vocabulary Check: Service Packages]
Match each term (1--6) with its correct definition (A--F). Write the letter in the box provided.

\begin{center}
\begin{tabularx}{\linewidth}{|c|X|c|X|}\hline
\textbf{Term} & \textbf{Definition} & \textbf{Term} & \textbf{Definition} \\\hline
1. Data plan & \hfill \framebox[3cm]{} & 4. Subscription fee & \hfill \framebox[3cm]{} \\\hline
2. Roaming charges & \hfill \framebox[3cm]{} & 5. Bundle & \hfill \framebox[3cm]{} \\\hline
3. Coverage area & \hfill \framebox[3cm]{} & 6. Unlimited calls & \hfill \framebox[3cm]{} \\\hline
\end{tabularx}
\end{center}

\begin{itemize}
\item[A.] The geographic region where the network signal is available. \textit{(Zone de couverture)}
\item[B.] A package combining several services (internet, calls, SMS) at a reduced price. \textit{(Forfait groupé)}
\item[C.] The monthly amount paid to keep the service active. \textit{(Frais d'abonnement)}
\item[D.] Extra costs for using your phone abroad or outside your home network. \textit{(Frais d'itinérance)}
\item[E.] The amount of internet data (MB/GB) included in your offer. \textit{(Forfait data)}
\item[F.] No limit on the duration or number of voice calls. \textit{(Appels illimités)}
\end{itemize}
\end{exercice}

\begin{aretenir}[Grammaire : Modaux de politesse et conseil + Comparatifs]
\textbf{Modaux courants au guichet télécom :}
\begin{itemize}
\item \cle{Can / Could} : capacité, permission, demande polie (\textit{Can I help you? Could I pay by MoMo?})
\item \cle{Would like} : souhait/préférence (\textit{I would like to subscribe...})
\item \cle{Should / Ought to} : conseil (\textit{You should take the Premium pack.})
\item \cle{Must / Have to} : obligation forte (\textit{You must show your ID.})
\item \cle{Will} : futur, prédiction, engagement (\textit{The speed will guarantee...})
\end{itemize}
\textbf{Comparatifs (révision) :}
\begin{itemize}
\item Court adjectif (1 syllabe) : \textit{fast $\to$ faster (than), the fastest}
\item Long adjectif ($\ge$2 syllabes) : \textit{expensive $\to$ more expensive (than), the most expensive}
\item Irréguliers : \textit{good $\to$ better $\to$ best ; bad $\to$ worse $\to$ worst ; much/many $\to$ more $\to$ most}
\end{itemize}
\end{aretenir}

\begin{exercice}[True / False Justified: ICTs at Work]
Read the statements below. Write \textbf{True} or \textbf{False} and justify your answer with a short phrase from the context of using ICTs to enhance work/learning.

\begin{enumerate}
\item Cloud storage allows accessing files from any device with an internet connection. \hfill \framebox[1.5cm]{} \\
Justification: \dotfill
\item Video conferencing tools require all participants to be in the same physical room. \hfill \framebox[1.5cm]{} \\
Justification: \dotfill
\item A strong password should contain personal information like your date of birth. \hfill \framebox[1.5cm]{} \\
Justification: \dotfill
\item Collaborative platforms (e.g., Google Docs) allow multiple users to edit a document simultaneously. \hfill \framebox[1.5cm]{} \\
Justification: \dotfill
\end{enumerate}
\end{exercice}

\begin{exercice}[Definitions: Recreational ICTs]
Define the following terms in your own words (one sentence each). Use the French gloss if needed.

\begin{enumerate}
\item \textbf{Streaming} \textit{(lecture en continu)} : \dotfill
\item \textbf{Social media} \textit{(réseaux sociaux)} : \dotfill
\item \textbf{Online gaming} \textit{(jeux en ligne)} : \dotfill
\item \textbf{Digital wellbeing} \textit{(bien-être numérique)} : \dotfill
\end{enumerate}
\end{exercice}

\begin{exercice}[Grammar Quick Check: Modals \& Comparatives]
Complete the dialogue between a Customer (C) and a Sales Agent (A) with the correct form of the words in brackets.

\begin{enumerate}
\item A: Good morning! How \_\_\_\_\_\_\_\_ (can / I / help) you today?
\item C: I \_\_\_\_\_\_\_\_ (would like / to subscribe) to a new internet package.
\item A: We have two offers. The Premium pack is \_\_\_\_\_\_\_\_ (fast) than the Basic one, but it costs more.
\item C: Which one \_\_\_\_\_\_\_\_ (you / recommend) for a student who watches online courses?
\item A: I \_\_\_\_\_\_\_\_ (suggest) the Premium pack. It has \_\_\_\_\_\_\_\_ (much) data.
\item C: Okay. \_\_\_\_\_\_\_\_ (Could / I / pay) by Mobile Money?
\item A: Yes, you \_\_\_\_\_\_\_\_ (can). You \_\_\_\_\_\_\_\_ (must / show) your ID card first.
\end{enumerate}
\end{exercice}

\courssub{Je m'entraîne}

\begin{methode}[Réussir un role-play au guichet télécom]
\textbf{Étapes :}
\begin{enumerate}
\item \textbf{Saluer} et identifier son besoin (internet, appels, budget, zone).
\item \textbf{Poser des questions} précises : prix, data, couverture, promotions, mode de paiement.
\item \textbf{Comparer} les offres (comparatifs : cheaper, faster, more data, better coverage).
\item \textbf{Exprimer sa préférence} : \textit{I'd rather choose... / I'll take this one.}
\item \textbf{Finaliser} : donner pièces d'identité, payer, récupérer le reçu.
\end{enumerate}
\textbf{Expressions utiles :} \textit{What does the package include? / Is there a promotion? / Does it cover my village? / Here is your receipt.}
\end{methode}

\begin{exercice}[Speaking Practice: Choosing a Package]
Work in pairs. Student A is a customer at a telecom agency (MTN / Orange / Camtel). Student B is the sales agent.
Use the prompts below to role-play a conversation (approx. 1 minute). Record key information in the table.

\begin{center}
\begin{tabularx}{\linewidth}{|l|X|}\hline
\textbf{Customer Needs} & \textbf{Agent's Offer} \\\hline
Wants internet for research + WhatsApp & \dotfill \\\hline
Budget: 5,000 FCFA / month & \dotfill \\\hline
Needs good coverage in rural area & \dotfill \\\hline
Prefers \textit{Mobile Money} payment & \dotfill \\\hline
\end{tabularx}
\end{center}

\textbf{Useful expressions:} \textit{I'd like to subscribe to... / What does the package include? / Is there a promotion? / I'd rather choose... / Here is your receipt.}
\end{exercice}

\begin{savaistu}[Contexte camerounais : Opérateurs télécom]
Au Cameroun, trois opérateurs principaux dominent le marché : \textbf{MTN Cameroon}, \textbf{Orange Cameroun} et \textbf{Camtel}. Les forfaits data populaires incluent \textit{MTN Pulse}, \textit{Orange Max} ou \textit{Camtel Fibre}. Le paiement par \textbf{Mobile Money} (MoMo, Orange Money) est la norme. La couverture 4G s'étend dans les zones urbaines et périurbaines ; la fibre optique (Camtel, MTN Business) se déploie dans les grandes villes (Yaoundé, Douala, Bafoussam, Garoua).
\end{savaistu}

\begin{exercice}[Listening Grid: ICTs for Learning]
Listen to the audio track (dialogue between a university student and an IT trainer) and complete the summary grid below.

\begin{center}
\begin{tabularx}{\linewidth}{|l|X|}\hline
\textbf{Information needed} & \textbf{Notes} \\\hline
Student's field of study & \dotfill \\\hline
Main ICT tool used for research & \dotfill \\\hline
Name of the recommended platform & \dotfill \\\hline
Advantage 1 (mentioned by trainer) & \dotfill \\\hline
Advantage 2 (mentioned by trainer) & \dotfill \\\hline
Potential risk / challenge warned about & \dotfill \\\hline
\end{tabularx}
\end{center}
\textit{Instruction: Écoutez deux fois. La première fois pour l'idée générale, la seconde pour les détails.}
\end{exercice}

\begin{exercice}[Reading Comprehension: ICTs for Leisure]
Read the short text and answer the questions.

\textit{In many Cameroonian communities, cybercafés have been replaced by smartphones as the main gateway to digital leisure. Young people spend hours on TikTok, Facebook, and WhatsApp. They watch comedy sketches, follow influencers, and play games like PUBG or Free Fire. While these activities provide entertainment and social connection, excessive screen time leads to sleep deprivation and reduced physical activity. Parents and educators encourage "digital detox" weekends to restore balance.}

\begin{enumerate}
\item According to the text, what has replaced cybercafés? \dotfill
\item List three recreational activities mentioned. \dotfill
\item What are two negative effects of excessive use? \dotfill
\item What solution do parents propose? \dotfill
\item Find a word in the text meaning "lack of something needed" (para 1). \dotfill
\end{enumerate}
\end{exercice}

\begin{exercice}[Writing: Guided Paragraph -- Digital Leisure]
Write a paragraph of \textbf{80--100 words} on the topic: \textit{``The role of smartphones in the leisure time of students in your neighbourhood.''}
Use the outline and linkers provided.

\textbf{Outline:}
\begin{enumerate}
\item Topic sentence: Smartphones are the main tool for leisure today.
\item Activities: social media, streaming, gaming (use \textit{Firstly, Moreover}).
\item Negative impacts: addiction, health issues (use \textit{However, Consequently}).
\item Recommendation: moderation, parental control (use \textit{Finally, In conclusion}).
\end{enumerate}

\textbf{Your paragraph:}
\dotfill
\dotfill
\dotfill
\dotfill
\end{exercice}

\courssub{Je me prépare au Bac}

\begin{exercice}[Bac Type: Reading Comprehension -- Social Media \& Youth]
Read the following text carefully and answer the questions that follow.

\textbf{Social Media and Cameroonian Youth: A Double-Edged Sword}

In Cameroon today, the smartphone has become an extension of the young person's hand. According to a 2023 study by the Ministry of Posts and Telecommunications, over 85\% of urban youth aged 15--24 own a smartphone with internet access. The primary use is not academic research or professional networking, but social media. Platforms like Facebook, WhatsApp, TikTok, and Instagram dominate screen time, averaging 4 to 6 hours daily.

This digital immersion offers undeniable benefits. Young entrepreneurs use WhatsApp Business and Facebook Marketplace to launch micro-enterprises -- selling clothes, cosmetics, or phone accessories -- without a physical shop. Students create WhatsApp groups to share notes, past exam papers, and tutorial videos, effectively building peer-to-peer learning communities. Furthermore, TikTok has birthed a new generation of Cameroonian comedians and artists who monetize their content, gaining national and even international recognition.

However, the shadows are lengthening. Cyberbullying is rampant; fake news spreads faster than verified information, sometimes inciting tribal hatred or panic. "Sexting" and exposure to inappropriate content threaten the moral fabric of minors. The most insidious danger is addiction: the dopamine loop of likes and notifications disrupts sleep patterns, reduces concentration in class, and fuels anxiety. A teacher in Douala lamented: ``My students know every TikTok trend but cannot summarize a textbook chapter.''

The solution is not prohibition but education. Digital literacy must enter the curriculum. Parents should negotiate "screen-free" zones and hours at home. Telecom operators could offer affordable "educational bundles" prioritizing access to learning platforms. Ultimately, the smartphone is a tool; its impact depends entirely on the intention of the user.

\textbf{Questions:}

\begin{enumerate}
\item \textbf{True or False?} Justify each answer by quoting a specific phrase from the text.
\begin{enumerate}
\item The majority of Cameroonian urban youth use smartphones mainly for academic research.
\item Social media allows some young people to earn money.
\item Cyberbullying is rare on Cameroonian social networks.
\item The author suggests banning smartphones for students.
\end{enumerate}

\item \textbf{Vocabulary in Context.} Find in the text words or expressions that mean:
\begin{enumerate}
\item A situation with both advantages and disadvantages (title/para 1): \dotfill
\item Deep involvement in an activity (para 1): \dotfill
\item Buying and selling goods online (para 2): \dotfill
\item Widespread / uncontrolled (para 3): \dotfill
\item The state of being dependent on a habit (para 3): \dotfill
\end{enumerate}

\item \textbf{Open Questions.} Answer in your own words (complete sentences).
\begin{enumerate}
\item Identify two positive uses of social media mentioned in paragraph 2.
\item What are three negative consequences of addiction mentioned in paragraph 3?
\item According to the last paragraph, who shares the responsibility for solving these problems?
\end{enumerate}

\item \textbf{Grammar / Sentence Transformation.} Rewrite the following sentences as indicated.
\begin{enumerate}
\item ``Students spend too much time on TikTok,'' the teacher said. \\
\textit{(Reported speech: The teacher complained that...)}
\item Young people \textbf{should} verify information before sharing it. \\
\textit{(Rewrite using: \textbf{It is essential that...}})
\item Although the internet is useful, it can be dangerous. \\
\textit{(Rewrite using: \textbf{Despite...}})
\item ``Do you have an educational data plan?'' the student asked the agent. \\
\textit{(Reported question: The student asked the agent...)}
\end{enumerate}

\item \textbf{Summary.} In \textbf{not more than 50 words}, summarize the author's view on the impact of social media on youth and the proposed way forward.
\end{enumerate}
\end{exercice}

\begin{exercice}[Bac Type: Grammar in Context -- Telecom Agency Dialogue]
Complete the dialogue below by choosing the correct option (A, B, C, or D) for each gap (1--10). Write the letter in the answer grid.

\textbf{Context:} A customer enters a \textit{Camtel} agency to subscribe to a fibre optic offer.

\textbf{Agent:} Good morning, Sir. Welcome to Camtel. How (1) \_\_\_\_\_\_\_\_ I help you?
\textbf{Customer:} Good morning. I (2) \_\_\_\_\_\_\_\_ like to get information about your fibre packages.
\textbf{Agent:} Certainly. We have three main offers. The "Essentiel" offer gives 10 Mbps for 15,000 FCFA. The "Confort" is (3) \_\_\_\_\_\_\_\_ at 50 Mbps for 25,000 FCFA. The "Premium" offers 100 Mbps.
\textbf{Customer:} Which one (4) \_\_\_\_\_\_\_\_ you recommend for a family of four who stream videos and work from home?
\textbf{Agent:} I (5) \_\_\_\_\_\_\_\_ strongly advise the "Premium" pack. It is (6) \_\_\_\_\_\_\_\_ expensive, but the speed (7) \_\_\_\_\_\_\_\_ guarantee smooth 4K streaming and video calls for everyone simultaneously.
\textbf{Customer:} That sounds good. (8) \_\_\_\_\_\_\_\_ I pay by Orange Money?
\textbf{Agent:} Yes, Sir. You (9) \_\_\_\_\_\_\_\_ pay by Mobile Money, card, or cash. However, you (10) \_\_\_\_\_\_\_\_ provide your National ID card and a proof of address for the contract.

\begin{center}
\begin{tabular}{|c|c|c|c|c|}\hline
\textbf{Gap} & \textbf{A} & \textbf{B} & \textbf{C} & \textbf{D} \\\hline
1 & can & must & will & shall \\\hline
2 & would & will & can & must \\\hline
3 & fastest & faster & the fastest & more fast \\\hline
4 & do & shall & would & will \\\hline
5 & would & should & can & must \\\hline
6 & more & most & the most & much \\\hline
7 & will & can & must & shall \\\hline
8 & Can & Must & Should & Would \\\hline
9 & can & must & should & would \\\hline
10 & have to & mustn't & don't have to & can \\\hline
\end{tabular}
\end{center}

\textbf{Answer Grid:}
\begin{center}
\begin{tabular}{|c|c|c|c|c|c|c|c|c|c|c|}\hline
Gap & 1 & 2 & 3 & 4 & 5 & 6 & 7 & 8 & 9 & 10 \\\hline
Answer & & & & & & & & & & \\\hline
\end{tabular}
\end{center}
\end{exercice}

\begin{exercice}[Bac Type: Guided Composition -- Data Commentary \& Essay]
The bar chart below shows the average daily screen time (in hours) of students in a Cameroonian technical high school, divided by activity.

\begin{popfigure}
\begin{tikzpicture}
\begin{axis}[
    ybar,
    bar width=18pt,
    width=\linewidth,
    height=8cm,
    ymin=0, ymax=5,
    ylabel={Average Hours per Day},
    symbolic x coords={Social Media, Video Streaming, Online Gaming, Educational Use, Other},
    xtick=data,
    x tick label style={rotate=30, anchor=east, font=\small},
    ymajorgrids=true,
    grid style=dashed,
    legend style={at={(0.5,-0.25)}, anchor=north, legend columns=-1},
    nodes near coords,
    nodes near coords align={vertical},
    popcurve/.style={mark=none, thick, popblue},
]
\addplot[popblue!80] coordinates {(Social Media, 2.8) (Video Streaming, 1.5) (Online Gaming, 1.2) (Educational Use, 0.7) (Other, 0.3)};
\legend{Screen Time}
\end{axis}
\end{tikzpicture}
\legende{Figure 1 : Average daily screen time by activity (Terminale Tech students, 2024).}
\end{popfigure}

\textbf{Part A: Data Commentary (approx. 80 words)}
Write 5 sentences interpreting the chart. You \textbf{must} use:
\begin{itemize}
\item At least two comparatives (e.g., higher than, the most, less... than).
\item At least two linkers (e.g., however, whereas, consequently, in contrast).
\item The present simple tense.
\end{itemize}

\textbf{Part B: Guided Composition (approx. 250 words)}
\textbf{Topic:} \textit{``The Impact of ICTs on the Recreational Habits of Students in Your Community.''}

\textbf{Write a composition following this structure:}

\begin{enumerate}
\item \textbf{Introduction (approx. 40 words):} Define ICTs in leisure context; state the general trend (increase); announce the plan (advantages / drawbacks / recommendations).
\item \textbf{Body Paragraph 1 -- Benefits (approx. 70 words):} Social connection, access to culture/talent discovery, skill acquisition (digital literacy), stress relief. Use linkers: \textit{Firstly, Furthermore, For instance}.
\item \textbf{Body Paragraph 2 -- Drawbacks (approx. 70 words):} Addiction, health issues (eyes, sleep, posture), cyber risks (bullying, scams), academic decline. Use linkers: \textit{However, Moreover, As a result}.
\item \textbf{Body Paragraph 3 -- Solutions / Recommendations (approx. 40 words):} Parental control, school digital literacy programs, self-regulation (screen time apps), affordable educational bundles. Use linkers: \textit{To address this, Therefore, Finally}.
\item \textbf{Conclusion (approx. 30 words):} Summarize main idea; final thought on balance/responsible use.
\end{enumerate}

\textbf{Marking Criteria (Bac Standard):} Content/Relevance (6), Organization/Coherence (4), Grammar/Accuracy (5), Vocabulary/Register (5). Total: 20 marks.
\end{exercice}



\newpage

\coursec{Corrigés des exercices}

\courssub{Corrigés de la section « Je vérifie mes connaissances »}

\begin{corrige}[Exercice 1 — Vocabulary Check: Service Packages]
\textbf{Answers:}
\begin{center}
\begin{tabular}{|c|c|l|}\hline
\textbf{Term} & \textbf{Letter} & \textbf{Explanation} \\\hline
1. Data plan & \textbf{E} & Amount of internet data (MB/GB) included in the offer. \\\hline
2. Roaming charges & \textbf{D} & Extra costs when using the phone abroad or outside the home network. \\\hline
3. Coverage area & \textbf{A} & Geographic region where the network signal is available. \\\hline
4. Subscription fee & \textbf{C} & Monthly amount paid to keep the service active. \\\hline
5. Bundle & \textbf{B} & Package combining several services at a reduced price. \\\hline
6. Unlimited calls & \textbf{F} & No limit on the duration or number of voice calls. \\\hline
\end{tabular}
\end{center}
\textbf{Points de vigilance :} Do not confuse \emph{data plan} (E) with \emph{subscription fee} (C): the data plan is the \emph{content} of the offer (how many GB), the fee is the \emph{price} paid. \emph{Roaming} (D) only applies outside the home network, not abroad only.
\end{corrige}

\begin{corrige}[Exercice 2 — True / False Justified: ICTs at Work]
\begin{enumerate}
\item \textbf{True.} Justification: cloud storage (Google Drive, OneDrive) keeps files on remote servers, so they can be accessed ``from any device with an internet connection''.
\item \textbf{False.} Justification: video conferencing (Zoom, Google Meet) is designed precisely for \emph{remote} meetings — participants join ``from different locations''.
\item \textbf{False.} Justification: a strong password must \emph{never} contain personal information (date of birth, name) because it is easy to guess; it should mix ``letters, numbers and symbols''.
\item \textbf{True.} Justification: collaborative platforms allow ``real-time collaboration'' — several users edit the same document simultaneously.
\end{enumerate}
\textbf{Points de vigilance :} The justification must be a short phrase, not a full paragraph; one relevant argument is enough. For item 3, the expected key idea is ``personal information makes a password weak''.
\end{corrige}

\begin{corrige}[Exercice 3 — Definitions: Recreational ICTs]
\textbf{Suggested answers (accept any grammatically correct equivalent):}
\begin{enumerate}
\item \textbf{Streaming} is watching or listening to media content (videos, music) in real time over the internet without downloading the file.
\item \textbf{Social media} are online platforms (Facebook, WhatsApp, TikTok) where users create profiles, share content and interact with others.
\item \textbf{Online gaming} is playing video games over the internet, often with or against other players around the world.
\item \textbf{Digital wellbeing} means keeping a healthy balance between screen time and offline life in order to protect physical and mental health.
\end{enumerate}
\textbf{Points de vigilance :} A definition needs a \emph{general class} + a \emph{specific feature} (``Streaming is \emph{a way of watching}... \emph{without downloading}''). Avoid circular definitions (``Streaming is when you stream'').
\end{corrige}

\begin{corrige}[Exercice 4 — Grammar Quick Check: Modals \& Comparatives]
\textbf{Answers:}
\begin{enumerate}
\item How \textbf{can I help} you today? (offre de service, inversion du modal)
\item I \textbf{would like to subscribe} to a new internet package. (\emph{would like} + \emph{to} + infinitif)
\item The Premium pack is \textbf{faster} than the Basic one. (comparatif d'adjectif court : fast $\to$ faster + \emph{than})
\item Which one \textbf{would you recommend}? (demande de conseil polie)
\item I \textbf{would suggest} the Premium pack. It has \textbf{more} data. (\emph{data} est ici indénombrable $\to$ \emph{more}, pas \emph{much} à l'affirmatif)
\item \textbf{Could I pay} by Mobile Money? (demande de permission polie, inversion)
\item Yes, you \textbf{can}. You \textbf{must show} your ID card first. (\emph{must} + infinitif sans \emph{to}, obligation)
\end{enumerate}
\textbf{Points de vigilance :} After a modal (\emph{can, could, must, would}), always use the \textbf{bare infinitive}: \emph{must show}, never \emph{must to show}. \emph{Would like} is followed by \emph{to} + infinitive. Comparative of a short adjective: \emph{faster}, never \emph{more faster}.
\end{corrige}

\courssub{Corrigés de la section « Je m'entraîne »}

\begin{corrige}[Exercice 5 — Speaking Practice: Choosing a Package]
\textbf{Sample dialogue (modèle de production attendue):}

\textbf{Agent:} Good morning, welcome to MTN. How can I help you?\\
\textbf{Customer:} Good morning. I would like to subscribe to an internet package. I need data for research and WhatsApp.\\
\textbf{Agent:} What is your monthly budget?\\
\textbf{Customer:} About 5\,000 FCFA per month.\\
\textbf{Agent:} Then I recommend the Pulse bundle: 5 GB, unlimited WhatsApp, for 4\,500 FCFA. It is cheaper than the Max offer and it has more data than the Basic plan.\\
\textbf{Customer:} Does it cover my village? I live in a rural area.\\
\textbf{Agent:} Yes, our 4G coverage reaches your area.\\
\textbf{Customer:} Perfect. I'd rather choose this one. Can I pay by Mobile Money?\\
\textbf{Agent:} Of course. Please show your ID card... Here is your receipt.

\textbf{Sample completed grid:}
\begin{center}
\begin{tabularx}{\linewidth}{|l|X|}\hline
\textbf{Customer Needs} & \textbf{Agent's Offer} \\\hline
Internet for research + WhatsApp & Pulse bundle: 5 GB + unlimited WhatsApp \\\hline
Budget: 5\,000 FCFA / month & 4\,500 FCFA / month (within budget) \\\hline
Good coverage in rural area & 4G coverage confirmed in the village \\\hline
Mobile Money payment & Accepted (MoMo), ID card required \\\hline
\end{tabularx}
\end{center}
\textbf{Points de vigilance (évaluation orale) :} greeting and closing (1 pt), at least two precise questions (2 pts), one comparative (1 pt), clear expression of preference with \emph{I'd rather / I'll take} (1 pt), fluency and pronunciation (1 pt).
\end{corrige}

\begin{corrige}[Exercice 6 — Listening Grid: ICTs for Learning]
\textbf{Expected answers (adapt to the actual audio track):}
\begin{center}
\begin{tabularx}{\linewidth}{|l|X|}\hline
\textbf{Information needed} & \textbf{Notes} \\\hline
Student's field of study & Computer Science (or Engineering / Medicine, according to the audio) \\\hline
Main ICT tool used for research & Laptop + Google Scholar / academic databases \\\hline
Recommended platform & Moodle / Google Classroom / Coursera / Khan Academy \\\hline
Advantage 1 & Access to up-to-date resources; flexibility (self-paced learning) \\\hline
Advantage 2 & Cost-effective; collaboration; certification \\\hline
Potential risk / challenge & Distraction, unreliable sources, plagiarism, poor connectivity, cybersecurity \\\hline
\end{tabularx}
\end{center}
\textbf{Points de vigilance :} First listening = general idea (who speaks, about what); second listening = details. Write \emph{notes}, not full sentences. Accept any plausible platform name actually heard in the recording.
\end{corrige}

\begin{corrige}[Exercice 7 — Reading Comprehension: ICTs for Leisure]
\begin{enumerate}
\item \textbf{Smartphones} have replaced cybercafés as the main gateway to digital leisure.
\item Three recreational activities: watching comedy sketches, following influencers, playing games (PUBG, Free Fire) — also accepted: using TikTok / Facebook / WhatsApp.
\item Two negative effects: \textbf{sleep deprivation} and \textbf{reduced physical activity}.
\item Parents and educators propose \textbf{``digital detox'' weekends} to restore balance.
\item The word is \textbf{deprivation} (in ``sleep deprivation'' = lack of sleep).
\end{enumerate}
\textbf{Points de vigilance :} Answers must come \emph{from the text}, not from personal opinion. For item 5, ``deprivation'' is the noun; ``deprived'' (adjective) is not in the text.
\end{corrige}

\begin{corrige}[Exercice 8 — Writing: Guided Paragraph — Digital Leisure]
\textbf{Model paragraph (approx. 95 words):}

Smartphones are the main tool for leisure today among students in my neighbourhood. \textbf{Firstly}, they spend hours on social media like WhatsApp and TikTok, chatting and watching videos. \textbf{Moreover}, streaming movies and playing online games such as PUBG are very popular. \textbf{However}, this excessive use leads to addiction and health problems like eye strain and poor sleep. \textbf{Consequently}, their school performance often drops. \textbf{Finally}, parents should set screen-time limits and encourage outdoor activities. \textbf{In conclusion}, smartphones are useful for entertainment, but they require responsible use and moderation.

\textbf{Barème indicatif (5 pts) :} respect of the outline and word count (1 pt), the required linkers used correctly (2 pts), grammar accuracy (1 pt), vocabulary related to recreational ICTs (1 pt).
\textbf{Points de vigilance :} One single paragraph (no bullet points); present simple for habits; do not exceed 100 words.
\end{corrige}

\courssub{Corrigés de la section « Je me prépare au Bac »}

\begin{corrige}[Exercice 9 — Bac Type: Reading Comprehension — Social Media \& Youth]
\textbf{1. True or False with quotations (4 $\times$ 1 pt):}
\begin{enumerate}
\item \textbf{False.} Quote: ``The primary use is not academic research or professional networking, but social media.''
\item \textbf{True.} Quote: ``Young entrepreneurs use WhatsApp Business and Facebook Marketplace to launch micro-enterprises'' / ``artists who monetize their content''.
\item \textbf{False.} Quote: ``Cyberbullying is rampant.''
\item \textbf{False.} Quote: ``The solution is not prohibition but education.''
\end{enumerate}
\textbf{Point de vigilance :} a True/False answer \emph{without a quotation} earns half the mark at the Bac — the justification is compulsory.

\textbf{2. Vocabulary in context (5 $\times$ 0,5 pt):}
\begin{enumerate}
\item \textbf{A double-edged sword} (title)
\item \textbf{Immersion} (``digital immersion'', para 1)
\item \textbf{Launch micro-enterprises} / \textbf{monetize their content} (para 2)
\item \textbf{Rampant} (para 3)
\item \textbf{Addiction} (para 3)
\end{enumerate}

\textbf{3. Open questions (3 $\times$ 1 pt) — suggested answers:}
\begin{enumerate}
\item Two positive uses: launching micro-enterprises without a physical shop (WhatsApp Business, Facebook Marketplace); building peer-to-peer learning communities (sharing notes, past papers, tutorial videos); also accepted: monetizing creative content on TikTok.
\item Three consequences of addiction: it disrupts sleep patterns, reduces concentration in class, and fuels anxiety.
\item Responsibility is shared by schools (digital literacy in the curriculum), parents (screen-free zones and hours), telecom operators (affordable educational bundles) — and ultimately the users themselves.
\end{enumerate}

\textbf{4. Sentence transformation (4 $\times$ 1 pt):}
\begin{enumerate}
\item The teacher complained that students \textbf{spent} too much time on TikTok. (backshift: spend $\to$ spent)
\item It is essential that young people \textbf{verify} information before sharing it. (subjunctive/base verb after ``it is essential that'')
\item \textbf{Despite} the usefulness of the internet, it can be dangerous. / Despite being useful, the internet can be dangerous. (\emph{despite} + noun or -ing, never \emph{despite of})
\item The student asked the agent \textbf{if/whether he or she had} an educational data plan. (reported yes/no question: \emph{if} + statement word order, backshift have $\to$ had)
\end{enumerate}
\textbf{Points de vigilance :} in reported questions, no inversion and no question mark (``asked if he had'', never ``asked if had he''); \emph{despite} is never followed by a full clause with \emph{although}-structure.

\textbf{5. Summary (2 pts) — model (48 words):}

Social media offers Cameroonian youth economic opportunities and learning communities, but it also causes addiction, cyberbullying and academic decline. The author advocates education rather than prohibition: digital literacy at school, parental screen limits and affordable educational bundles. The smartphone's impact depends on the user's intention.

\textbf{Barème indicatif total : 12 pts} (T/F: 4 ; Vocab: 2,5 ; Open questions: 3 ; Transformations: 4 — ajusté à 12 ; ou ramener à /20 selon la grille de l'\'etablissement).
\end{corrige}

\begin{corrige}[Exercice 10 — Bac Type: Grammar in Context — Telecom Agency Dialogue]
\textbf{Answer grid:}
\begin{center}
\begin{tabular}{|c|c|c|c|c|c|c|c|c|c|c|}\hline
Gap & 1 & 2 & 3 & 4 & 5 & 6 & 7 & 8 & 9 & 10 \\\hline
Answer & \textbf{A} & \textbf{A} & \textbf{B} & \textbf{C} & \textbf{A} & \textbf{A} & \textbf{A} & \textbf{A} & \textbf{A} & \textbf{A} \\\hline
\end{tabular}
\end{center}

\textbf{Detailed justifications:}
\begin{enumerate}
\item \textbf{A — can}: ``How can I help you?'' is the standard service formula.
\item \textbf{A — would}: ``I would like to...'' expresses a polite wish.
\item \textbf{B — faster}: comparison between two offers (Essentiel vs Confort) $\to$ comparative \emph{faster}; \emph{the fastest} would require comparing all three.
\item \textbf{C — would}: ``Which one would you recommend?'' — polite request for advice.
\item \textbf{A — would}: ``I would strongly advise...'' — polite recommendation.
\item \textbf{A — more}: \emph{expensive} is a long adjective $\to$ comparative \emph{more expensive}.
\item \textbf{A — will}: prediction about a future result (``the speed will guarantee...'').
\item \textbf{A — Can}: asking for permission/possibility to pay by Orange Money.
\item \textbf{A — can}: the agent states a possibility (several payment methods).
\item \textbf{A — have to}: external/contractual obligation (provide ID and proof of address).
\end{enumerate}
\textbf{Points de vigilance :} gap 3 is the classic trap: with only \emph{two} items compared, use the comparative (\emph{faster}), not the superlative. Gap 10: \emph{mustn't} (interdiction) and \emph{don't have to} (absence d'obligation) both contradict the context — a contract \emph{requires} documents.

\textbf{Barème indicatif :} 10 $\times$ 1 pt = 10 pts ; no justification required in the exam grid, but the reasoning above should be explained in class.
\end{corrige}

\begin{corrige}[Exercice 11 — Bac Type: Guided Composition — Data Commentary \& Essay]
\textbf{Part A — Data Commentary (model, approx. 80 words):}

Social Media is \textbf{the most} time-consuming activity, at \num{2,8} hours per day. \textbf{However}, Video Streaming (\num{1,5} h) and Online Gaming (\num{1,2} h) are significantly \textbf{lower than} Social Media. \textbf{In contrast}, Educational Use (\num{0,7} h) is \textbf{less than half} of the time spent on social media. \textbf{Consequently}, recreational activities dominate screen time, totalling \num{5,5} hours against only \num{0,7} hours for learning. \textbf{Whereas} Other activities are negligible (\num{0,3} h), this imbalance raises serious concerns about students' digital wellbeing.

\textbf{Vérification des données :} $\num{2,8} + \num{1,5} + \num{1,2} = \num{5,5}$ h de loisirs ; $\num{0,7}$ h éducatif ; $\num{2,8} \div \num{0,7} = 4$ (les réseaux sociaux représentent 4 fois le temps éducatif) ; total général $= \num{5,8}$ h/jour. Tous les chiffres cités correspondent au graphique.

\textbf{Barème Part A (4 pts) :} two correct comparatives (1 pt), two linkers correctly used (1 pt), accurate data reported (1 pt), present simple and coherence (1 pt).

\textbf{Part B — Guided Composition (model, approx. 250 words):}

\textbf{Introduction.} Information and Communication Technologies (ICTs), such as smartphones, social media and streaming platforms, now shape how students spend their free time. In my community, screen time has increased dramatically in recent years. This essay examines the benefits and drawbacks of this trend before proposing ways to promote balanced use.

\textbf{Benefits.} \textbf{Firstly}, ICTs strengthen social connection: students keep in touch with friends and family through WhatsApp or Facebook, even across long distances. \textbf{Furthermore}, they give access to culture and talent discovery; \textbf{for instance}, many young Cameroonians showcase comedy, music or dance on TikTok and gain national recognition. Recreational use also builds digital literacy — navigating apps, editing videos, managing an online identity — skills that are valuable for future employment.

\textbf{Drawbacks.} \textbf{However}, excessive use brings serious risks. \textbf{Moreover}, addiction to likes and notifications disrupts sleep, causes eye strain and poor posture. \textbf{As a result}, academic performance declines as homework time shrinks. Cyberbullying, scams and exposure to inappropriate content also threaten students' safety and mental health.

\textbf{Solutions.} \textbf{To address this}, parents should activate parental controls and negotiate screen-free meals and bedtimes. \textbf{Therefore}, schools must integrate digital wellbeing into the curriculum. \textbf{Finally}, telecom operators could offer affordable educational bundles prioritizing learning platforms.

\textbf{Conclusion.} In conclusion, ICTs enrich students' leisure but require conscious regulation. Balance, not prohibition, is the key: used wisely, the smartphone empowers; used blindly, it enslaves.

\textbf{Barème Bac (20 pts) :} Content/Relevance — all five parts of the plan addressed (6) ; Organization/Coherence — clear paragraphs, linkers, logical progression (4) ; Grammar/Accuracy — tenses, modals, comparatives, agreement (5) ; Vocabulary/Register — unit lexis (streaming, bundle, digital wellbeing...), formal tone (5).

\textbf{Points de vigilance :} respect the word count ($\pm$ 10 \%) ; every body paragraph must open with a linker from the list ; do not copy whole sentences from the chart commentary into the essay ; conclusion must not introduce a new argument.
\end{corrige}

\newpage

\coursec{Fiche bilan}

\begin{popfigure}
\begin{tikzpicture}[
  mindmap,
  grow cyclic,
  every node/.style={concept, execute at begin node=\hskip0pt},
  concept color=popblue,
  level 1/.append style={level distance=4.5cm, sibling angle=360/6},
  level 2/.append style={level distance=3cm, sibling angle=360/4},
  text=white,
  font=\small\bfseries
]
\node[concept, scale=1.2] {Media \& Communication}
  child[concept color=poporange] { node[concept, align=center]{Subscribing\\ Packages} }
  child[concept color=popgreen] { node[concept, align=center]{Grammar\\ Revision} }
  child[concept color=poppurple] { node[concept, align=center]{Listening:\\ ICT Work/Learn} }
  child[concept color=poppink] { node[concept, align=center]{Reading:\\ ICT Recreation} }
  child[concept color=popgold] { node[concept, align=center]{Writing:\\ Composition} }
  child[concept color=popdark] { node[concept, align=center]{Integration\\ \& Evaluation} };
\end{tikzpicture}
\legende{Mindmap of Lesson 05: Media and Communication}
\end{popfigure}

\begin{bilan}

\courssub{Key Vocabulary by Theme}

\begin{center}
\begin{tabularx}{\linewidth}{|l|X|}\hline
\rowcolor{popblueL}
\textbf{Theme} & \textbf{Key Words \& Expressions (English / Français)} \\ \hline
\cle{Subscribing to Packages} & \textbf{subscribe to} (s'abonner), \textbf{service package} (forfait), \textbf{data plan} (forfait data), \textbf{prepaid/postpaid} (prépayé/postpayé), \textbf{bundle} (pack combiné), \textbf{coverage} (couverture réseau), \textbf{bandwidth} (bande passante), \textbf{unlimited} (illimité), \textbf{roaming} (itinérance), \textbf{contract/terms} (contrat/conditions), \textbf{cancel/renew} (résilier/renouveler) \\ \hline
\cle{ICT for Work \& Learning} & \textbf{enhance} (améliorer), \textbf{remote work} (télétravail), \textbf{e-learning} (e-learning), \textbf{video conferencing} (visioconférence), \textbf{cloud storage} (stockage cloud), \textbf{collaboration tools} (outils collaboratifs), \textbf{productivity} (productivité), \textbf{upskilling} (montée en compétences), \textbf{digital literacy} (culture numérique), \textbf{connectivity} (connectivité) \\ \hline
\cle{ICT for Recreation} & \textbf{streaming} (streaming), \textbf{binge-watching} (visionnage intensif), \textbf{gaming} (jeu vidéo), \textbf{social media} (réseaux sociaux), \textbf{content creation} (création de contenu), \textbf{scrolling} (défilement), \textbf{screen time} (temps d'écran), \textbf{digital wellbeing} (bien-être numérique), \textbf{addiction} (dépendance), \textbf{leisure} (loisir) \\ \hline
\end{tabularx}
\end{center}

\courssub{Grammar General Revision: Essential Points}

\begin{center}
\begin{tabularx}{\linewidth}{|l|X|l|}\hline
\rowcolor{popblueL}
\textbf{Area} & \textbf{Key Rules / Structures} & \textbf{Example} \\ \hline
\cle{Tenses} & Present Simple (habits), Present Continuous (now/future arranged), Past Simple (finished), Present Perfect (link past-present), Future (will/going to) & \emph{I \textbf{check} emails daily. / She \textbf{is uploading} a file.} \\ \hline
\cle{Modals} & \textbf{can/could} (ability), \textbf{should/must} (advice/obligation), \textbf{may/might} (possibility), \textbf{would} (preference/polite request) & \emph{\textbf{Would} you like a bigger data plan? / You \textbf{must} read the terms.} \\ \hline
\cle{Conditionals} & Zero (facts), First (real future), Second (unreal present), Third (unreal past) & \emph{If I \textbf{subscribe}, I \textbf{get} a discount. (First)} \\ \hline
\cle{Passive Voice} & be + V3 (focus on action/receiver) & \emph{The package \textbf{is offered} by the provider.} \\ \hline
\cle{Reported Speech} & Tense backshift, pronoun/time changes & \emph{He said he \textbf{wanted} a new plan.} \\ \hline
\cle{Relative Clauses} & who/which/that/whose/where (defining vs non-defining) & \emph{The app \textbf{which} I use is free.} \\ \hline
\cle{Comparatives/Superlatives} & -er/more + than, the -est/most (choices) & \emph{This plan is \textbf{cheaper than} the other.} \\ \hline
\cle{Prepositions} & \textbf{on} the internet, \textbf{via} email, \textbf{through} an app, \textbf{for} a purpose, \textbf{with} a tool & \emph{I pay \textbf{via} mobile money.} \\ \hline
\end{tabularx}
\end{center}

\courssub{Express Methods}

\begin{methode}[Speaking: Expressing Preferences for Service Packages]
\begin{enumerate}
\item \textbf{Greet \& Identify Need:} \emph{``Hello, I'd like to subscribe to an internet package.''}
\item \textbf{Ask Key Questions:} \emph{``How much data does it include?''}, \emph{``Is it prepaid or postpaid?''}, \emph{``What is the coverage like in my area?''}
\item \textbf{Compare Options:} Use comparatives. \emph{``Package A is cheaper, but Package B has better speed.''}
\item \textbf{State Preference:} \emph{``I'd prefer the unlimited plan because I stream a lot.''}, \emph{``I'd rather choose the prepaid option.''}
\item \textbf{Confirm Details:} Price, duration, renewal terms. \emph{``So, it's 5000 FCFA per month, renewable automatically?''}
\item \textbf{Close:} \emph{``I'll take this one. How do I activate it?''}
\end{enumerate}
\end{methode}

\begin{methode}[Writing: Composition on Recreational ICT Use]
\begin{enumerate}
\item \textbf{Analyse the Topic:} Identify key terms (e.g., \emph{advantages, disadvantages, youth, community}).
\item \textbf{Outline:} Introduction (Hook + Thesis), Body 1 (Benefits: access, creativity, connection), Body 2 (Drawbacks: addiction, health, isolation), Body 3 (Solutions/Balance: screen time limits, parental control), Conclusion (Summary + Opinion/Recommendation).
\item \textbf{Draft:} Use connectors (\emph{Furthermore, However, Consequently, For instance}). Use vocabulary from the lesson.
\item \textbf{Revise:} Check grammar (tenses, modals, conditionals), spelling, punctuation, word count.
\item \textbf{Final Polish:} Ensure formal/semi-formal register. Clear paragraphs.
\end{enumerate}
\end{methode}

\end{bilan}

\begin{aretenir}[Checklist avant le Bac]
$\square$ I can name at least 10 words related to \cle{subscribing to phone/internet packages}.\par
$\square$ I can \textbf{compare two service packages} using comparatives and justify my choice orally.\par
$\square$ I master the \textbf{formation and use of all key tenses} (Present, Past, Perfect, Future).\par
$\square$ I can use \textbf{modals} correctly for advice (\emph{should}), obligation (\emph{must}), preference (\emph{would rather}).\par
$\square$ I can transform sentences into \textbf{Passive Voice} and \textbf{Reported Speech}.\par
$\square$ I understand listening texts about \textbf{ICTs for work/learning} (keywords: remote work, e-learning, cloud).\par
$\square$ I can extract main ideas from a reading text on \textbf{ICTs for recreation} (streaming, gaming, wellbeing).\par
$\square$ I can write a structured \textbf{composition (150--200 words)} on recreational ICT use with intro, body, conclusion.\par
$\square$ I use \textbf{linking words} (addition, contrast, cause/effect, example) to organise my writing.\par
$\square$ I know the \textbf{difference between prepaid/postpaid, limited/unlimited, bandwidth/coverage}.\par
$\square$ I can participate in a \textbf{role-play: customer vs. sales agent} for a subscription scenario.\par
\end{aretenir}

\findecours

\end{document}
